% Lesson 33 — Vocabulary: Words and expressions related to environmental issues like climate change and global warming. # Semaine / Période Chapitre Titre Date effective Visa de l'enseignant — Anglais Tle A — Cours signé OnBuch+
% Programme officiel MINESEC. Compiler avec : tectonic EN33-vocabulary-words-and-expressions-related-to-environment.tex
\newcommand{\DOCMATIERE}{Anglais}
\newcommand{\DOCNIVEAU}{Tle A}
\newcommand{\DOCMODULE}{Chapter 3 : Environment, Well-being and Health}
\newcommand{\DOCLECON}{EN33}
\newcommand{\DOCTITRE}{Lesson 33 — Vocabulary: Words and expressions related to environmental issues like climate change and global warming. \# Semaine / Période Chapitre Titre Date effective Visa de l'enseignant}
% OnBuch+ — préambule « cours signés » ENRICHI (compilé avec tectonic / XeLaTeX)
% Système visuel « Pop » d'OnBuch+ : fond crème (jamais blanc), bordures encre
% épaisses, ombres « dures » décalées, titres Archivo Black en capitales.
% Version MATHÉMATIQUES : graphes (pgfplots/tikz), boîtes pédagogiques
% (définition, propriété, méthode, attention, le savais-tu, exercice résolu,
% exercice, TP, bilan), page de garde et signature OnBuch+ ancrée en bas.
%
% Macros attendues dans main.tex AVANT \input{preamble} :
%   \DOCMATIERE  \DOCNIVEAU  \DOCMODULE  \DOCLECON (numéro)  \DOCTITRE
\documentclass[11pt]{article}

\usepackage{fontspec}
\usepackage[english,french]{babel} % français = langue principale ; l'anglais via \eng{..} / textebox
\usepackage[a4paper,margin=2cm,top=3.4cm,bottom=2.8cm,headheight=1.4cm,footskip=1.3cm]{geometry}
\usepackage{amsmath,amssymb,amsfonts}
\usepackage{mathtools} % \xmapsto, \coloneqq... souvent utilisés par les modèles
\usepackage{cancel} % simplifications barrées (bilans de forces)
% \abs{z} (module, valeur absolue) et \norm{u} (norme), utilisés par les modèles
\providecommand{\abs}{}\let\abs\relax\DeclarePairedDelimiter{\abs}{\lvert}{\rvert}
\providecommand{\norm}{}\let\norm\relax\DeclarePairedDelimiter{\norm}{\lVert}{\rVert}
\usepackage[table]{xcolor} % [table] : fournit aussi \rowcolors (alternance de lignes)
\usepackage{pagecolor}
\usepackage{tikz}
\usetikzlibrary{arrows.meta,positioning,calc,shapes.geometric,shapes.misc,shapes.arrows,decorations.pathmorphing,decorations.pathreplacing,decorations.markings,patterns,shadows,mindmap,trees,fit,backgrounds,matrix,intersections,angles,quotes}
\usepackage{pgfplots}
\usepackage[version=4]{mhchem} % équations nucléaires \ce{^{235}_{92}U}
\usepackage[european,siunitx]{circuitikz} % schémas électriques
\pgfplotsset{compat=1.18}
\usepgfplotslibrary{fillbetween,groupplots} % aires entre courbes (calcul intégral)
\usepackage{enumitem}
\usepackage{fancyhdr}
% [most] charge toutes les bibliothèques tcolorbox (listings, minted, raster,
% documentation...) et gonfle inutilement la mémoire TeX sur un document long
% (~300 boîtes) : on ne charge que ce qui est réellement utilisé.
\usepackage[skins,breakable]{tcolorbox}
\usepackage{graphicx}
\usepackage{colortbl}
\usepackage{booktabs}
\usepackage{tabularx}
\usepackage{diagbox} % case d'en-tête diagonale des tableaux à double entrée
% Colonnes extensibles centrée / gauche / droite (C, L, R), souvent utilisées par les modèles
\newcolumntype{C}{>{\centering\arraybackslash}X}
\newcolumntype{L}{>{\raggedright\arraybackslash}X}
\newcolumntype{R}{>{\raggedleft\arraybackslash}X}
\usepackage{array}
\usepackage{multicol}
\usepackage{multirow}
\usepackage{siunitx}
% parse-numbers=false : accepte aussi \SI{2,5 \times 10^{-3}}{...}
\sisetup{locale=FR,output-decimal-marker={,},group-minimum-digits=4,parse-numbers=false}
\DeclareSIUnit\ohm{\ensuremath{\symup{\Omega}}} % U+2126 absent de la police
\DeclareSIUnit\division{div} % réglages d'oscilloscope (ms/div, V/div)
\DeclareSIUnit\tr{tr} % tours (vitesse de rotation, ex. \SI{1500}{\tr\per\minute})
\DeclareSIUnit\molar{M} % concentration molaire (ex. \SI{3}{\molar}, \SI{1}{\micro\molar})
\DeclareSIUnit\an{an} % année (le modèle confond parfois avec \year, un compteur TeX) — ex. \SI{5}{\kilo\meter\per\an}
\DeclareSIUnit\FCFA{FCFA} % franc CFA (ex. \SI{500}{\FCFA\per\kilo\gram})
\DeclareSIUnit\degreeBrix{\textdegree Brix} % teneur en sucre d'un sirop/jus (réfractométrie)
\DeclareSIUnit\month{mois} % le modèle utilise parfois \month, un compteur TeX, comme unité de durée
\usepackage{qrcode}
% \degree : commande fréquemment utilisée par les modèles pour un angle ou une
% température en dehors de \SI{}{} (non fournie par les paquets chargés).
\providecommand{\degree}{^{\circ}}
% \qed (fin de démonstration), utilisé par les modèles sans amsthm
\providecommand{\qed}{\hfill$\square$}
% \barre : double barre des valeurs interdites dans un tableau de variations (tkz-tab)
\providecommand{\barre}{\ensuremath{\|}}
% environnement proof (amsthm non chargé) utilisé par les modèles
\ifdefined\proof\else\newenvironment{proof}[1][Démonstration]{\par\noindent\textit{#1.}\ }{\qed\par}\fi
\usepackage{lastpage}
\usepackage{hyperref}
\hypersetup{hidelinks}

% ============================================================
% Palette « Pop » OnBuch+ — mêmes hex que la classe P (plus_ui.dart)
% ============================================================
\definecolor{poporange}{HTML}{F59321}
\definecolor{popdark}{HTML}{E07A0C}
\definecolor{popcream}{HTML}{FFF6E8}
\definecolor{popink}{HTML}{1C1714}
\definecolor{poppurple}{HTML}{7A5AE0}
\definecolor{popgreen}{HTML}{2E9E6A}
\definecolor{poppink}{HTML}{E0457B}
\definecolor{popblue}{HTML}{2F7DE1}
\definecolor{popgold}{HTML}{C98A12}
\definecolor{popmuted}{HTML}{8A7F76}
\definecolor{popline}{HTML}{EADFD2}
\definecolor{popgray}{HTML}{8A7F76}
\colorlet{popgrey}{popgray}
% Alias tolérants (le modèle nomme parfois les couleurs différemment)
\colorlet{popwhite}{white}
\colorlet{popblack}{popink}
\colorlet{popred}{poppink}
\colorlet{popyellow}{popgold}
% Teintes claires (fonds de tableaux/encadrés) — le modèle nomme parfois
% les couleurs avec un suffixe L (ex. popblueL) sans les avoir définies.
\colorlet{poporangeL}{poporange!20}
\colorlet{popdarkL}{popdark!20}
\colorlet{poppurpleL}{poppurple!20}
\colorlet{popgreenL}{popgreen!20}
\colorlet{poppinkL}{poppink!20}
\colorlet{popblueL}{popblue!20}
\colorlet{popgoldL}{popgold!20}
\colorlet{popredL}{popred!20}
\colorlet{popgrayL}{popgray!20}
% Teintes claires pour fonds de tableaux / zones
\colorlet{popblueL}{popblue!10!white}
\colorlet{poporangeL}{poporange!14!white}
\colorlet{popgreenL}{popgreen!12!white}
\colorlet{poppurpleL}{poppurple!12!white}
\colorlet{poppinkL}{poppink!12!white}
\colorlet{popgoldL}{popgold!14!white}

\pagecolor{popcream}

% ============================================================
% Typographie
% ============================================================
\defaultfontfeatures{Ligatures=TeX}
\setmainfont{PlusJakartaSans-Medium.ttf}[Path=../../../fonts/,BoldFont=PlusJakartaSans-Bold.ttf]
% Maths en sans-serif (Fira Math) avec chiffres et lettres droites de Plus
% Jakarta, pour que les formules se fondent dans le texte.
\usepackage[math-style=ISO,bold-style=ISO]{unicode-math}
\setmathfont{FiraMath-Regular.otf}[Path=../../../fonts/]
\setmathfont{PlusJakartaSans-Medium.ttf}[Path=../../../fonts/,range={up/num,up/{Latin,latin},\mathup}]
\usepackage{icomma} % virgule décimale sans espace en mode math
% Caractères mathématiques Unicode absents des polices de texte (Plus Jakarta,
% Archivo Black) : rendus par leur équivalent mathématique, en texte comme en
% formule — sinon ils s'affichent en carré vide (ex. « Arithmétique dans ℤ »).
\usepackage{newunicodechar}
\newunicodechar{ℝ}{\ensuremath{\mathbb{R}}}
\newunicodechar{ℤ}{\ensuremath{\mathbb{Z}}}
\newunicodechar{ℂ}{\ensuremath{\mathbb{C}}}
\newunicodechar{ℕ}{\ensuremath{\mathbb{N}}}
\newunicodechar{ℚ}{\ensuremath{\mathbb{Q}}}
\newunicodechar{𝔻}{\ensuremath{\mathbb{D}}}
\newunicodechar{α}{\ensuremath{\alpha}}
\newunicodechar{β}{\ensuremath{\beta}}
\newunicodechar{γ}{\ensuremath{\gamma}}
\newunicodechar{δ}{\ensuremath{\delta}}
\newunicodechar{ε}{\ensuremath{\varepsilon}}
\newunicodechar{θ}{\ensuremath{\theta}}
\newunicodechar{λ}{\ensuremath{\lambda}}
\newunicodechar{μ}{\ensuremath{\mu}}
\newunicodechar{σ}{\ensuremath{\sigma}}
\newunicodechar{τ}{\ensuremath{\tau}}
\newunicodechar{φ}{\ensuremath{\varphi}}
\newunicodechar{ψ}{\ensuremath{\psi}}
\newunicodechar{ω}{\ensuremath{\omega}}
\newunicodechar{Σ}{\ensuremath{\Sigma}}
% majuscules produites par \MakeUppercase dans les titres (α-aminés -> Α)
\newunicodechar{Α}{\ensuremath{\alpha}}
\newunicodechar{Β}{\ensuremath{\beta}}
\newunicodechar{Γ}{\ensuremath{\Gamma}}
\newunicodechar{Δ}{\ensuremath{\Delta}}
\newunicodechar{Ε}{\ensuremath{\varepsilon}}
\newunicodechar{Θ}{\ensuremath{\Theta}}
\newunicodechar{Λ}{\ensuremath{\Lambda}}
\newunicodechar{Μ}{\ensuremath{\mu}}
\newunicodechar{Τ}{\ensuremath{\tau}}
\newunicodechar{Φ}{\ensuremath{\Phi}}
\newunicodechar{Ψ}{\ensuremath{\Psi}}
\newunicodechar{Ω}{\ensuremath{\Omega}}
\newunicodechar{↦}{\ensuremath{\mapsto}}
\newunicodechar{∈}{\ensuremath{\in}}
\newunicodechar{∉}{\ensuremath{\notin}}
\newunicodechar{∧}{\ensuremath{\wedge}}
\newunicodechar{⇔}{\ensuremath{\Leftrightarrow}}
\newunicodechar{⟺}{\ensuremath{\Longleftrightarrow}}
\newunicodechar{⇒}{\ensuremath{\Rightarrow}}
\newunicodechar{⟹}{\ensuremath{\Longrightarrow}}
\newunicodechar{∘}{\ensuremath{\circ}}
\newunicodechar{∪}{\ensuremath{\cup}}
\newunicodechar{∩}{\ensuremath{\cap}}
\newunicodechar{≡}{\ensuremath{\equiv}}
\newunicodechar{⊂}{\ensuremath{\subset}}
\newunicodechar{⊥}{\ensuremath{\perp}}
\newunicodechar{∃}{\ensuremath{\exists}}
\newunicodechar{∀}{\ensuremath{\forall}}
\newunicodechar{∛}{\ensuremath{\sqrt[3]{\,}}}
\newunicodechar{∜}{\ensuremath{\sqrt[4]{\,}}}
\newunicodechar{⃗}{\ensuremath{{}^{\to}}}
\newunicodechar{ⁿ}{\ifmmode^{n}\else\textsuperscript{\ensuremath{n}}\fi}
\newunicodechar{ˣ}{\ifmmode^{x}\else\textsuperscript{\ensuremath{x}}\fi}
\newunicodechar{ᵇ}{\ifmmode^{b}\else\textsuperscript{\ensuremath{b}}\fi}
\newunicodechar{ʳ}{\ifmmode^{r}\else\textsuperscript{\ensuremath{r}}\fi}
\newunicodechar{ᵘ}{\ifmmode^{u}\else\textsuperscript{\ensuremath{u}}\fi}
\newunicodechar{ʸ}{\ifmmode^{y}\else\textsuperscript{\ensuremath{y}}\fi}
\newunicodechar{ᵃ}{\ifmmode^{a}\else\textsuperscript{\ensuremath{a}}\fi}
\newunicodechar{ᵉ}{\ifmmode^{e}\else\textsuperscript{\ensuremath{e}}\fi}
\newunicodechar{ᵏ}{\ifmmode^{k}\else\textsuperscript{\ensuremath{k}}\fi}
\newunicodechar{ᵖ}{\ifmmode^{p}\else\textsuperscript{\ensuremath{p}}\fi}
\newunicodechar{ᴺ}{\ifmmode^{N}\else\textsuperscript{\ensuremath{N}}\fi}
\newunicodechar{ᶜ}{\ifmmode^{c}\else\textsuperscript{\ensuremath{c}}\fi}
\newunicodechar{ʰ}{\ifmmode^{h}\else\textsuperscript{\ensuremath{h}}\fi}
\newunicodechar{ᵠ}{\ifmmode^{q}\else\textsuperscript{\ensuremath{q}}\fi}
\newunicodechar{ₙ}{\ifmmode_{n}\else\textsubscript{\ensuremath{n}}\fi}
\newunicodechar{ₐ}{\ifmmode_{a}\else\textsubscript{\ensuremath{a}}\fi}
\newunicodechar{ᵢ}{\ifmmode_{i}\else\textsubscript{\ensuremath{i}}\fi}
\newunicodechar{ᵣ}{\ifmmode_{r}\else\textsubscript{\ensuremath{r}}\fi}
\newunicodechar{ₖ}{\ifmmode_{k}\else\textsubscript{\ensuremath{k}}\fi}
\newunicodechar{ᵦ}{\ifmmode_{\beta}\else\textsubscript{\ensuremath{\beta}}\fi}
\newunicodechar{ₓ}{\ifmmode_{x}\else\textsubscript{\ensuremath{x}}\fi}
\newfontfamily\popdisplay{ArchivoBlack-Regular.ttf}[Path=../../../fonts/]
\newfontfamily\poplogo{SpaceGrotesk-Bold.ttf}[Path=../../../fonts/]
\newfontfamily\popsemi{PlusJakartaSans-SemiBold.ttf}[Path=../../../fonts/]
\newfontfamily\popxbold{PlusJakartaSans-ExtraBold.ttf}[Path=../../../fonts/]
\newfontfamily\popxboldls{PlusJakartaSans-ExtraBold.ttf}[Path=../../../fonts/,LetterSpace=3.0]

\setlist[enumerate]{leftmargin=1.7em,itemsep=5pt,topsep=4pt}
\setlist[itemize]{leftmargin=1.7em,itemsep=4pt,topsep=4pt,label={\color{poporange}\textbullet}}
\setlength{\parindent}{0pt}
\setlength{\parskip}{5pt}
\renewcommand{\arraystretch}{1.25}

% pgfplots : style Pop par défaut
\pgfplotsset{
  every axis/.append style={axis line style={popink,line width=1pt},tick style={popink},
    grid style={popline,line width=0.5pt},label style={font=\small},tick label style={font=\footnotesize}},
  popcurve/.style={poporange,line width=1.8pt,no markers,smooth},
}
% Même style disponible dans un \draw TikZ simple (hors pgfplots)
\tikzset{popcurve/.style={poporange,line width=1.8pt}}
\tikzset{popgrid/.style={popline,very thin}}
\colorlet{popcurve}{poporange} % le nom du style est parfois utilisé comme couleur

% ============================================================
% Pill / squiggle
% ============================================================
\newcommand{\poppill}[3]{%
  \tikz[baseline=-0.55ex]\node[draw=popink,line width=1.2pt,rounded corners=2.6pt,%
    fill=#2,inner xsep=6.5pt,inner ysep=3pt]{%
    \popxboldls\fontsize{7}{7}\selectfont\color{#3}\MakeUppercase{#1}};%
}
\newcommand{\popsquiggle}[1]{%
  \tikz[baseline=-0.4ex]{
    \draw[#1,line width=2.6pt,line cap=round]
      plot [smooth] coordinates {(0,0.09) (0.34,0.01) (0.68,0.21) (1.02,0.01) (1.36,0.10)};
  }%
}
% Mot-clé surligné (marqueur orange sous le texte)
\newcommand{\cle}[1]{{\popxbold\color{popdark}#1}}

% ============================================================
% En-tête / pied
% ============================================================
\pagestyle{fancy}
\fancyhf{}
\renewcommand{\headrulewidth}{0pt}
\renewcommand{\footrulewidth}{0pt}
\lhead{\raisebox{-0.15ex}{\poplogo\fontsize{13}{13}\selectfont\color{popink}OnBuch\color{poporange}+}}
\rhead{\poppill{\DOCMATIERE\ \textbullet\ \DOCNIVEAU\ \textbullet\ Leçon \DOCLECON}{white}{popink}}
\lfoot{\popxbold\fontsize{7.6}{7.6}\selectfont\color{popmuted}\MakeUppercase{Cours signé OnBuch+}\ \textbullet\ {\popsemi\DOCTITRE}}
\rfoot{\tikz[baseline=-0.6ex]\node[rounded corners=8.5pt,fill=popink,minimum height=17pt,minimum width=17pt,inner xsep=5pt,inner sep=0]{\popxbold\fontsize{8}{8}\selectfont\color{popcream}\thepage\,/\,\pageref*{LastPage}};}

% ============================================================
% Titres
% ============================================================
\newcommand{\chaptitre}[2]{%
  \begin{tcolorbox}[enhanced,colback=poporange,colframe=popink,boxrule=2.6pt,arc=11pt,
      drop shadow=popink,left=15pt,right=15pt,top=12pt,bottom=11pt,before skip=3pt,after skip=18pt]
    \poppill{#2}{popink}{popcream}\par\vspace{9pt}
    {\raggedright\hyphenpenalty=10000\exhyphenpenalty=10000\popdisplay\fontsize{22}{24}\selectfont\color{popcream}\MakeUppercase{#1}\par}
  \end{tcolorbox}
}
% Section numérotée : pastille orange avec le numéro + titre Archivo
\newcounter{popsec}
\newcommand{\coursec}[1]{%
  \stepcounter{popsec}\setcounter{popsub}{0}%
  \par\vspace{16pt}\noindent
  \begin{minipage}{\linewidth}
  \tikz[baseline=-0.7ex]\node[draw=popink,line width=1.6pt,fill=poporange,rounded corners=4pt,minimum size=22pt,inner sep=0]{\popdisplay\fontsize{12}{12}\selectfont\color{popcream}\thepopsec};%
  \hspace{8pt}{\popdisplay\fontsize{15}{17}\selectfont\color{popink}\MakeUppercase{#1}}\par
  \vspace{4pt}\noindent{\color{poporange}\rule{36pt}{3.6pt}}
  \end{minipage}\par\nopagebreak\vspace{8pt}
}
\newcounter{popsub}
\newcommand{\courssub}[1]{%
  \stepcounter{popsub}%
  \par\vspace{9pt}\noindent{\popxbold\fontsize{11.5}{13}\selectfont\color{popink}{\color{poporange}\thepopsec.\thepopsub}\hspace{6pt}#1}\par\nopagebreak\vspace{4pt}
}

% ============================================================
% Boîtes pédagogiques — même recette : fond blanc, bordure épaisse colorée,
% ombre dure encre, badge plein. Titre optionnel [#1].
% ============================================================
\tcbset{popbase/.style={breakable,enhanced,colback=white,boxrule=2.3pt,arc=9pt,
  shadow={3pt}{-3pt}{0pt}{popink},left=14pt,right=14pt,top=11pt,bottom=11pt,before skip=11pt,after skip=15pt,
  fontupper=\popsemi\fontsize{10.6}{14.5}\selectfont\color{popink},
  fontlower=\popsemi\fontsize{10.6}{14.5}\selectfont\color{popink}}}
% Coche dessinée (le glyphe ✓ n'existe pas dans les polices du document)
\newcommand{\popcheck}[1]{\tikz[baseline=-0.5ex]\draw[#1,line width=1.6pt,line cap=round,line join=round](-0.13,0.0)--(-0.04,-0.09)--(0.13,0.1);}
\providecommand{\coche}{\popcheck{popgreen}}
\providecommand{\resultat}[1]{\par\smallskip\noindent\fcolorbox{popgreen}{popgreenL}{\parbox{\dimexpr\linewidth-2\fboxsep-2\fboxrule\relax}{\textbf{Résultat.} #1}}\par\smallskip}
\newcommand{\popboxhead}[3]{\poppill{#1}{#2}{white}\ifx\relax#3\relax\else\hspace{6pt}{\popxbold\fontsize{10.6}{12}\selectfont\color{popink}#3}\fi\par\vspace{7pt}}

\newtcolorbox{aretenir}[1][]{popbase,colframe=popgreen,before upper={\popboxhead{À retenir}{popgreen}{#1}}}
% Texte en anglais (document de lecture, dialogue, extrait) : cadre
% d'étude. Un titre optionnel (source, auteur) s'affiche dans l'en-tête.
\newtcolorbox{textebox}[1][]{popbase,colframe=popblue,before upper={\popboxhead{Text}{popblue}{#1}\begin{otherlanguage}{english}},after upper={\end{otherlanguage}}}
% Marques ✓ / ✗ (forme correcte / fautive) souvent utilisées par les modèles
\providecommand{\cross}{\textcolor{poppink}{\ensuremath{\times}}}
\providecommand{\xmark}{\cross}\providecommand{\crossmark}{\cross}\providecommand{\cmark}{\ensuremath{\checkmark}}
% Ligne de réponse / justification à compléter (inventée par les modèles)
\providecommand{\justification}[1]{\par\noindent\textit{Justification~:} \dotfill\par}
% \textipa (package tipa non chargé) : la police ne couvre pas l'API -> simple passage
\providecommand{\textipa}[1]{#1}
% Soulignement (ulem non chargé) : \uline, \ul
\providecommand{\uline}[1]{\underline{#1}}\providecommand{\ul}[1]{\underline{#1}}
% \glossaire{\item ...} inventée par les modèles : liste de glossaire
\providecommand{\glossaire}[1]{\begin{itemize}#1\end{itemize}}
% \alert (beamer) inventée par les modèles : texte mis en évidence
\providecommand{\alert}[1]{\textbf{\textcolor{poppink}{#1}}}
% variantes ASCII d'ordinaux écrites en macro
\providecommand{\iere}{\textsuperscript{re}}\providecommand{\ieme}{\textsuperscript{e}}\providecommand{\ere}{\textsuperscript{re}}\providecommand{\eme}{\textsuperscript{e}}\providecommand{\er}{\textsuperscript{er}}\providecommand{\ie}{\textsuperscript{e}}
% Ordinaux français écrits en macro par les modèles : 1\ière, 3\ième
\newcommand{\ière}{\textsuperscript{re}}\newcommand{\ième}{\textsuperscript{e}}
% \exemple (inventée par les modèles) : étiquette « Exemple : »
\providecommand{\exemple}{\textbf{Exemple~:}\ }
% \square utilisable aussi hors mode math (cases à cocher)
\AtBeginDocument{\let\squareMath\square\renewcommand{\square}{\ensuremath{\squareMath}}}
% Mot ou phrase en anglais à l'intérieur d'un paragraphe français
\newcommand{\eng}[1]{\ifmmode\text{\textit{#1}}\else\begingroup\selectlanguage{english}\itshape #1\endgroup\fi}
\let\esp\eng % alias toléré
\newtcolorbox{exemplebox}[1][]{popbase,colframe=poporange,before upper={\popboxhead{Exemple}{poporange}{#1}}}
\newtcolorbox{definition}[1][]{popbase,colframe=popblue,before upper={\popboxhead{Définition}{popblue}{#1}}}
\newtcolorbox{propriete}[1][]{popbase,colframe=poppurple,before upper={\popboxhead{Propriété}{poppurple}{#1}}}
\newtcolorbox{methode}[1][]{popbase,colframe=popgold,before upper={\popboxhead{Méthode}{popgold}{#1}}}
\newtcolorbox{attention}[1][]{popbase,colframe=poppink,before upper={\popboxhead{Attention}{poppink}{#1}}}
\newtcolorbox{experience}[1][]{popbase,colframe=popblue,colback=popblueL,before upper={\popboxhead{Expérience}{popblue}{#1}}}
% « Le savais-tu ? » : carte violette pleine (contexte, culture, Cameroun)
\newtcolorbox{savaistu}[1][]{popbase,colback=poppurple,colframe=popink,
  fontupper=\popsemi\fontsize{10.4}{14.2}\selectfont\color{white},
  before upper={\poppill{Le savais-tu\,?}{white}{poppurple}\ifx\relax#1\relax\else\hspace{6pt}{\popxbold\color{white}#1}\fi\par\vspace{7pt}}}
% Exercice résolu : énoncé en haut, \tcblower puis la solution sur fond crème
\newcounter{popexo}\newcounter{popexr}
\newtcolorbox{exoresolu}[1][]{popbase,colframe=poporange,colbacklower=popcream,
  segmentation style={popink,line width=1.2pt,dashed},
  before upper={\stepcounter{popexr}\popboxhead{Exercice résolu \thepopexr}{poporange}{#1}},
  before lower={\poppill{Solution}{popink}{popcream}\par\vspace{6pt}}}
% Exercice (non résolu en place) — la correction est dans la partie Corrigés
\newtcolorbox{exercice}[1][]{popbase,colframe=popink,
  before upper={\stepcounter{popexo}\popboxhead{Exercice \thepopexo}{popink}{#1}}}
% Corrigé
\newtcolorbox{corrige}[1][]{popbase,colframe=popgreen,colback=popgreenL,
  before upper={\popboxhead{Corrigé}{popgreen}{#1}}}
% Objectifs / prérequis / bilan
\newtcolorbox{objectifs}{popbase,colframe=popink,colback=poporangeL,
  before upper={\popboxhead{Objectifs de la leçon}{popink}{}}}
\newtcolorbox{prerequis}{popbase,colframe=popink,colback=popblueL,
  before upper={\popboxhead{Prérequis}{popblue}{}}}
\newtcolorbox{bilan}{popbase,colframe=popink,colback=white,boxrule=2.8pt,
  before upper={\popboxhead{Fiche bilan}{poporange}{L'essentiel à connaître pour l'examen}}}
% Figure encadrée avec légende
\newtcolorbox{popfigure}[1][]{enhanced,breakable=false,colback=white,colframe=popline,boxrule=1.4pt,arc=8pt,
  left=10pt,right=10pt,top=10pt,bottom=8pt,before skip=10pt,after skip=12pt,halign=center,
  lower separated=false,halign lower=center,
  fontlower=\popsemi\fontsize{9}{11.5}\selectfont\color{popmuted},
  #1}
\newcommand{\legende}[1]{\par\vspace{6pt}{\popsemi\fontsize{9}{11.5}\selectfont\color{popmuted}#1}}

% ============================================================
% Signature OnBuch+
% ============================================================
\newcommand{\poplogoinline}[1]{{\poplogo\fontsize{#1}{#1}\selectfont\color{popink}OnBuch\color{poporange}+}}
\newcommand{\signatureonbuch}{%
  \begin{tcolorbox}[enhanced,colback=popink,colframe=popink,boxrule=2.2pt,arc=10pt,
      drop shadow=poporange,left=14pt,right=14pt,top=10pt,bottom=10pt,before skip=0pt,after skip=0pt]
    \tikz[baseline=-0.7ex]\node[circle,fill=popgreen,draw=popcream,line width=1.3pt,minimum size=22pt,inner sep=0]{\popcheck{white}};%
    \hspace{9pt}\begin{minipage}[c]{0.62\linewidth}
      {\popxbold\fontsize{12}{13}\selectfont\color{popcream}Signé\ \textbullet\ {\poplogo OnBuch\color{poporange}+}}\par\vspace{2pt}
      {\raggedright\popsemi\fontsize{8}{10.5}\selectfont\color{popline}\MakeUppercase{Équipe pédagogique OnBuch+}\\\MakeUppercase{Conforme au programme officiel MINESEC}\par}
    \end{minipage}\hfill
    {\poplogo\fontsize{22}{22}\selectfont\color{popcream}OnBuch\color{poporange}+}
  \end{tcolorbox}%
}

% ============================================================
% Page de garde
% #1 titre · #2 matière · #3 niveau · #4 module · #5 n° leçon · #6 durée ·
% #7 description / objectifs (texte ou itemize)
% ============================================================
\newcommand{\pagegarde}[7]{%
  \thispagestyle{empty}
  \newgeometry{a4paper,margin=1.7cm,top=1.6cm,bottom=1.4cm}%
  \begin{tikzpicture}[remember picture,overlay]
    \fill[poporange!18] ([xshift=-3.2cm,yshift=-3.6cm]current page.north east) circle (4.6cm);
    \fill[poppurple!14] ([xshift=2.4cm,yshift=7.6cm]current page.south west) circle (3.8cm);
  \end{tikzpicture}%
  {\poplogo\fontsize{30}{30}\selectfont\color{popink}OnBuch\color{poporange}+}\hfill
  \raisebox{0.6ex}{\poppill{Cours signé \textbullet\ Premium}{popink}{popcream}}\par
  \vspace{1pt}\popsquiggle{poppurple}\par
  \vspace{8pt}
  \begin{tcolorbox}[enhanced,colback=poporange,colframe=popink,boxrule=2.8pt,arc=13pt,
      drop shadow=popink,left=18pt,right=18pt,top=12pt,bottom=13pt,after skip=10pt]
    \poppill{#2 \textbullet\ #3}{popink}{popcream}\hspace{5pt}\poppill{#4}{white}{popink}\par\vspace{8pt}
    {\popdisplay\fontsize{14}{14}\selectfont\color{popink}LEÇON #5}\par\vspace{4pt}
    {\raggedright\hyphenpenalty=10000\exhyphenpenalty=10000\popdisplay\fontsize{25}{27}\selectfont\color{popcream}\MakeUppercase{#1}\par}
    \vspace{8pt}
    {\popxbold\fontsize{9.5}{9.5}\selectfont\color{popink}\MakeUppercase{Durée conseillée : #6}}
  \end{tcolorbox}
  \begin{tcolorbox}[enhanced,colback=white,colframe=popink,boxrule=2.2pt,arc=10pt,
      drop shadow=popink,left=16pt,right=16pt,top=10pt,bottom=10pt,after skip=8pt]
    \poppill{Dans cette leçon}{poppurple}{white}\par\vspace{5pt}
    \popsemi\fontsize{9.3}{12.6}\selectfont\color{popink}#7
  \end{tcolorbox}
  \noindent\begin{minipage}[t]{0.487\linewidth}
    \begin{tcolorbox}[enhanced,colback=popgreen,colframe=popink,boxrule=2.2pt,arc=10pt,
        drop shadow=popink,left=11pt,right=11pt,top=10pt,bottom=10pt,height=3.1cm,valign=center]
      \tcbox[colback=white,colframe=popink,boxrule=1.3pt,arc=5pt,boxsep=0pt,nobeforeafter,
        left=3pt,right=3pt,top=3pt,bottom=3pt,valign=center]{\qrcode[height=1.9cm]{https://play.google.com/store/apps/details?id=com.luvvix.onbuch&hl=fr}}%
      \hspace{8pt}\begin{minipage}[c]{0.5\linewidth}
        {\popdisplay\fontsize{11}{12.5}\selectfont\color{white}\MakeUppercase{Télécharge OnBuch}}\par\vspace{4pt}
        {\raggedright\popsemi\fontsize{8.4}{11}\selectfont\color{white}Gratuit \textbullet\ Google Play\\Tuteur IA, annales, concours, résultats\par}
      \end{minipage}
    \end{tcolorbox}
  \end{minipage}\hfill
  \begin{minipage}[t]{0.487\linewidth}
    \begin{tcolorbox}[enhanced,colback=poppurple,colframe=popink,boxrule=2.2pt,arc=10pt,
        drop shadow=popink,left=11pt,right=11pt,top=10pt,bottom=10pt,height=3.1cm,valign=center]
      \tikz[baseline=-0.7ex]\node[circle,fill=white,draw=popink,line width=1.3pt,minimum size=1.6cm,inner sep=0]{\popdisplay\fontsize{24}{24}\selectfont\color{poppurple}+};%
      \hspace{8pt}\begin{minipage}[c]{0.6\linewidth}
        {\popdisplay\fontsize{11}{12.5}\selectfont\color{white}\MakeUppercase{Passe à OnBuch+}}\par\vspace{4pt}
        {\raggedright\popsemi\fontsize{8.4}{11}\selectfont\color{white}Tous les cours signés, Tuteur IA illimité, fiches PDF — onglet OnBuch+ de l'app\par}
      \end{minipage}
    \end{tcolorbox}
  \end{minipage}
  \vspace{8pt}
  \popcontenu
  \vfill
  \signatureonbuch
  \restoregeometry
  \newpage
}

% Rangée « Ce cours contient » (page de garde)
\newcommand{\poptile}[3]{% #1 couleur #2 gros texte #3 légende
  \begin{tcolorbox}[enhanced,colback=#1,colframe=popink,boxrule=2pt,arc=9pt,shadow={2.5pt}{-2.5pt}{0pt}{popink},
      left=6pt,right=6pt,top=9pt,bottom=9pt,halign=center,valign=center,height=1.75cm,width=\linewidth,nobeforeafter]
    {\popdisplay\fontsize{12.5}{13}\selectfont\color{white}\MakeUppercase{#2}}\par\vspace{3pt}
    {\centering\popsemi\fontsize{7.8}{9.5}\selectfont\color{white}#3\par}
  \end{tcolorbox}}
\newcommand{\popcontenu}{%
  \par\noindent{\popxboldls\fontsize{7.5}{7.5}\selectfont\color{popmuted}CE COURS SIGNÉ CONTIENT}\par\vspace{7pt}
  \noindent\begin{minipage}[t]{0.235\linewidth}\poptile{popblue}{Cours}{complet, illustré, pas à pas}\end{minipage}\hfill
  \begin{minipage}[t]{0.235\linewidth}\poptile{poporange}{Méthodes}{et exercices résolus}\end{minipage}\hfill
  \begin{minipage}[t]{0.235\linewidth}\poptile{poppink}{Exercices}{type examen + corrigés}\end{minipage}\hfill
  \begin{minipage}[t]{0.235\linewidth}\poptile{popgreen}{Fiche bilan}{l'essentiel à retenir}\end{minipage}\par}

% Sommaire
\newtcolorbox{sommaire}{enhanced,colback=white,colframe=popink,boxrule=2.2pt,arc=10pt,
  drop shadow=popink,left=18pt,right=18pt,top=13pt,bottom=13pt,before skip=6pt,after skip=20pt,
  before upper={\poppill{Sommaire}{poppurple}{white}\par\vspace{9pt}\popsemi\fontsize{10.6}{14.5}\selectfont\color{popink}}}

% Signature de fin de cours — poussée tout en bas de la dernière page
\newcommand{\findecours}{%
  \par\vfill\nopagebreak
  \begin{center}\popsquiggle{poporange}\end{center}\vspace{4pt}
  \signatureonbuch
}
\AtBeginDocument{\def\llbracket{\mathopen{[\mkern-3.5mu[}}\def\rrbracket{\mathclose{]\mkern-3.5mu]}}} % intervalles d'entiers (glyphe ⟦ absent de la police)
% Glyphes absents de Fira Math : redéfinis à partir de glyphes disponibles
\AtBeginDocument{%
  \def\setminus{\mathbin{\backslash}}\let\smallsetminus\setminus
  \def\vdots{\mathord{\vbox{\baselineskip3.2pt\lineskiplimit0pt\kern4pt\hbox{.}\hbox{.}\hbox{.}}}}%
  \def\ddots{\mathinner{\mkern1mu\raise6.5pt\hbox{.}\mkern2mu\raise3.6pt\hbox{.}\mkern2mu\raise0.7pt\hbox{.}\mkern1mu}}%
  \def\triangle{\mathord{\scalebox{0.9}{$\symup{\Delta}$}}}%
  \def\nabla{\mathord{\raisebox{\depth}{\scalebox{1}[-1]{$\symup{\Delta}$}}}}%
  \def\checkmark{\popcheck{popdark}}%
  \def\star{\mathbin{\ast}}\def\bigstar{\mathord{\ast}}%
  \def\diamond{\mathbin{\scalebox{0.6}{$\Diamond$}}}%
  \def\bigcup{\mathop{\mathchoice{\raisebox{-0.3em}{\scalebox{1.7}{$\cup$}}}{\scalebox{1.25}{$\cup$}}{\cup}{\cup}}\displaylimits}%
  \def\bigcap{\mathop{\mathchoice{\raisebox{-0.3em}{\scalebox{1.7}{$\cap$}}}{\scalebox{1.25}{$\cap$}}{\cap}{\cap}}\displaylimits}%
  \def\lesseqqgtr{\mathrel{\lessgtr}}\def\bigoplus{\mathop{\oplus}\displaylimits}%
}
\providecommand{\ssi}{si et seulement si} % abréviation fréquente
\AtBeginDocument{\providecommand{\wideparen}{\overparen}} % arc de cercle

\begin{document}

\pagegarde{Lesson 33 — Vocabulary: Words and expressions related to environmental issues like climate change and global warming. \# Semaine / Période Chapitre Titre Date effective Visa de l'enseignant}{Anglais}{Tle A}{Chapter 3 : Environment, Well-being and Health}{EN33}{2 h}{Cette leçon de vocabulaire présente les mots et expressions essentiels pour parler des problèmes environnementaux : changement climatique, réchauffement planétaire, pollution et solutions écologiques. Elle propose champs lexicaux, collocations, faux amis, familles de mots et mise en pratique en contexte, en vue du Bac.
\vspace{4pt}
\begin{multicols}{2}\small
\begin{itemize}
\item À la fin de cette leçon, je sais identifier et définir le vocabulaire clé lié au changement climatique et au réchauffement planétaire.
\item À la fin de cette leçon, je sais classer ce lexique en champs thématiques (causes, effets, solutions).
\item À la fin de cette leçon, je sais employer les collocations correctes (ex. 'reduce emissions', 'tackle climate change').
\item À la fin de cette leçon, je sais utiliser des expressions idiomatiques liées à l'environnement.
\item À la fin de cette leçon, je sais éviter les faux amis et les calques du français (ex. 'pollution' vs 'pollution', 'to protect the environment').
\item À la fin de cette leçon, je sais former des mots de la même famille (pollute, pollution, pollutant, polluted).
\end{itemize}\end{multicols}}

\begin{sommaire}
\begin{multicols}{2}
\begin{enumerate}
\item Champ lexical 1 : Climate change and global warming
\item Champ lexical 2 : Pollution and waste
\item Champ lexical 3 : Solutions and protection
\item Collocations et expressions idiomatiques
\item Faux amis et pièges des francophones
\item Familles de mots, prononciation et emploi en contexte
\item Activité d'intégration
\item Méthodes et exercices résolus
\item Exercices
\item Corrigés des exercices
\item Fiche bilan
\end{enumerate}
\end{multicols}
\end{sommaire}

\begin{objectifs}
\begin{itemize}
\item À la fin de cette leçon, je sais identifier et définir le vocabulaire clé lié au changement climatique et au réchauffement planétaire.
\item À la fin de cette leçon, je sais classer ce lexique en champs thématiques (causes, effets, solutions).
\item À la fin de cette leçon, je sais employer les collocations correctes (ex. 'reduce emissions', 'tackle climate change').
\item À la fin de cette leçon, je sais utiliser des expressions idiomatiques liées à l'environnement.
\item À la fin de cette leçon, je sais éviter les faux amis et les calques du français (ex. 'pollution' vs 'pollution', 'to protect the environment').
\item À la fin de cette leçon, je sais former des mots de la même famille (pollute, pollution, pollutant, polluted).
\item À la fin de cette leçon, je sais prononcer et orthographier correctement les mots difficiles (environment, government, sustainable).
\item À la fin de cette leçon, je sais réemployer ce lexique dans des phrases, dialogues et courts textes à l'oral comme à l'écrit.
\end{itemize}
\end{objectifs}

\begin{prerequis}
\begin{itemize}
\item Vocabulaire général de la nature vu en Première (weather, seasons, animals, plants).
\item Structures de base : present simple, present continuous et futur (will / be going to) pour parler des tendances.
\item Connecteurs logiques simples (because, so, however, therefore).
\item Notions de formation des mots : préfixes et suffixes courants (-tion, -able, -ist, un-, re-).
\end{itemize}
\end{prerequis}

\newpage

\coursec{Champ lexical 1 : Climate change and global warming}

\courssub{Introduction : pourquoi ce vocabulaire ?}

\begin{textebox}[A situation to imagine]
Imagine the long dry season in the Far North of Cameroon: rivers shrink, crops fail and families move to the cities. On the radio, journalists speak of “climate change” and “global warming”. In Buea, students march with placards reading “Save Mount Cameroon!”. To understand the news, join a debate or write a Bac essay on the environment, you need the right English words — and you must avoid the traps of French. This lesson gives you that vocabulary toolkit.
\end{textebox}

\textbf{Problématique.} Comment nommer, en anglais correct, les causes et les effets du changement climatique ? Quels mots utiliser pour décrire la sécheresse à Maroua, les inondations à Douala ou la fonte des glaces au pôle Nord ? Cette première section construit le \cle{champ lexical} du \eng{climate change} (changement climatique) et du \eng{global warming} (réchauffement planétaire) : les mots-clés, les phénomènes, les causes, la prononciation et les pièges à éviter.

\courssub{Observation : un texte pour découvrir le lexique}

\begin{textebox}[Reading text: “A hotter Far North” — adapted from a school radio report]
In the Far North of Cameroon, droughts are becoming longer and more severe. Scientists say this is a direct consequence of climate change. Thirty years ago, the rainy season started in May; today, farmers near Maroua sometimes wait until July for the first heavy rain. When the rain finally comes, it can be so violent that it causes floods which destroy houses and roads.

The main cause of global warming is well known: human activities release greenhouse gases into the atmosphere. The most important of these gases is carbon dioxide, or CO2. It is produced when we burn fossil fuels such as coal, oil and gas. Factories, lorries and old taxis all send exhaust fumes into the air. Deforestation makes the problem worse, because trees naturally absorb CO2. When forests are cut down for firewood or farmland, that carbon stays in the atmosphere.

The effects are visible everywhere. Temperatures are rising, heatwaves are more frequent, and the ice caps at the poles are melting. As a result, sea levels are rising and coastal cities like Douala and Limbe face a greater risk of flooding. In the Sahel region, desertification is turning farmland into desert, and thousands of families must leave their villages.

The good news is that everyone can help. By planting trees, saving energy and reducing waste, each of us can reduce our carbon footprint — the total amount of greenhouse gases that our daily life produces. As one student in Buea said: “The Earth is our only home. There is no planet B.”
\end{textebox}

\textbf{Glossaire (mots difficiles) :}

\begin{center}
\begin{tabularx}{\linewidth}{|X|X|X|}\hline
\rowcolor{popblueL} English & Nature & Français \\ \hline
drought & nom & la sécheresse \\ \hline
severe & adjectif & grave, sévère \\ \hline
flood & nom & l'inondation \\ \hline
greenhouse gas & nom & le gaz à effet de serre \\ \hline
to release & verbe & rejeter, libérer \\ \hline
fossil fuels & nom pluriel & les énergies fossiles \\ \hline
exhaust fumes & nom pluriel & les gaz d'échappement \\ \hline
to absorb & verbe & absorber \\ \hline
heatwave & nom & la vague de chaleur, la canicule \\ \hline
ice caps & nom pluriel & les calottes glaciaires \\ \hline
to melt & verbe & fondre \\ \hline
sea level & nom & le niveau de la mer \\ \hline
desertification & nom & la désertification \\ \hline
carbon footprint & nom & l'empreinte carbone \\ \hline
\end{tabularx}
\end{center}

\textbf{Questions de compréhension (orales ou écrites) :}
\begin{enumerate}
\item What are two consequences of climate change mentioned in the text?
\item Which human activities produce carbon dioxide?
\item Why does deforestation make global warming worse?
\item Which Cameroonian cities are threatened by rising sea levels?
\item What is a “carbon footprint”?
\end{enumerate}

\courssub{Les deux notions fondamentales : climate change et global warming}

\begin{definition}[Global warming]
\eng{Global warming} is \emph{the gradual increase in the Earth's average temperature, mainly caused by greenhouse gases}. Traduction : le \cle{réchauffement planétaire} (ou réchauffement climatique), c'est-à-dire l'augmentation progressive de la température moyenne de la Terre.
\end{definition}

\begin{definition}[Greenhouse effect]
\eng{The greenhouse effect} is \emph{the process by which gases in the atmosphere trap heat from the sun}. Traduction : l'\cle{effet de serre}, c'est-à-dire le processus par lequel les gaz de l'atmosphère retiennent la chaleur du soleil, comme les parois de verre d'une serre.
\end{definition}

\begin{definition}[Climate change]
\eng{Climate change} désigne les \cle{changements durables du climat} d'une région ou de la planète (températures, pluies, vents) sur plusieurs décennies. Le \eng{global warming} est la \emph{cause principale} du \eng{climate change} actuel : le réchauffement provoque sécheresses, inondations et fonte des glaces.
\end{definition}

\begin{propriete}[Vocabulaire des causes : les gaz et les émissions]
\begin{itemize}
\item \eng{greenhouse gases (GHG)} : les gaz à effet de serre. On dit \eng{greenhouse gas emissions} pour « les émissions de gaz à effet de serre ».
\item \eng{carbon dioxide (CO2)} : le dioxyde de carbone. Prononciation : « car-beune daï-ox-aïde ».
\item \eng{emissions} : les émissions, les rejets (toujours au pluriel dans ce sens). Verbe associé : \eng{to emit} (émettre, rejeter).
\item \eng{carbon footprint} : l'empreinte carbone, c'est-à-dire la quantité totale de gaz à effet de serre produite par une personne, une école ou un pays.
\item \eng{fossil fuels} : les énergies fossiles — \eng{coal} (le charbon), \eng{oil} (le pétrole), \eng{gas} (le gaz). Verbe associé : \eng{to burn fossil fuels} (brûler des énergies fossiles).
\item \eng{exhaust fumes} : les gaz d'échappement (des voitures, des motos, des bendskins).
\item \eng{industrial pollution} : la pollution industrielle ; \eng{overconsumption} : la surconsommation.
\end{itemize}
\end{propriete}

\begin{exemplebox}[Exemples traduits]
\eng{Burning fossil fuels releases carbon dioxide into the atmosphere.} « Brûler des énergies fossiles libère du dioxyde de carbone dans l'atmosphère. »

\eng{Old lorries in Douala produce a lot of exhaust fumes.} « Les vieux camions de Douala produisent beaucoup de gaz d'échappement. »

\eng{Every family can reduce its carbon footprint by saving energy.} « Chaque famille peut réduire son empreinte carbone en économisant l'énergie. »

\eng{Factories emit greenhouse gases twenty-four hours a day.} « Les usines émettent des gaz à effet de serre vingt-quatre heures sur vingt-quatre. »
\end{exemplebox}

\courssub{Les phénomènes : causes et effets visibles}

\begin{propriete}[Vocabulaire des phénomènes climatiques]
\begin{itemize}
\item \eng{rising temperatures} : la hausse des températures. Note le participe présent \eng{rising} utilisé comme adjectif : \eng{rising sea levels} (la montée du niveau des mers).
\item \eng{heatwave} : la vague de chaleur, la canicule (un seul mot en anglais !).
\item \eng{drought} : la sécheresse. Prononciation : « draoute ».
\item \eng{flood} : l'inondation ; verbe : \eng{to flood} (inonder). Prononciation : « fleud ».
\item \eng{melting ice caps} : la fonte des calottes glaciaires. Verbe : \eng{to melt} (fondre).
\item \eng{rising sea levels} : l'élévation du niveau des mers.
\item \eng{desertification} : la désertification (les terres agricoles deviennent des déserts).
\item \eng{deforestation} : la déforestation, le déboisement. Verbe : \eng{to deforest} ; contraire : \eng{reforestation} (le reboisement).
\end{itemize}
\end{propriete}

\begin{exemplebox}[Exemples en contexte camerounais]
\eng{The drought in the Far North has destroyed the millet harvest.} « La sécheresse dans l'Extrême-Nord a détruit la récolte de mil. »

\eng{Every rainy season, floods block the streets of Douala.} « Chaque saison des pluies, des inondations bloquent les rues de Douala. »

\eng{The ice caps are melting faster than scientists expected.} « Les calottes glaciaires fondent plus vite que les scientifiques ne le prévoyaient. »

\eng{Deforestation around Mount Cameroon threatens many animal species.} « La déforestation autour du mont Cameroun menace de nombreuses espèces animales. »

\eng{Rising sea levels could flood coastal villages near Limbe.} « L'élévation du niveau des mers pourrait inonder des villages côtiers près de Limbé. »
\end{exemplebox}

\courssub{Tableau récapitulatif : 20 termes essentiels}

\begin{center}
\begin{tabularx}{\linewidth}{|X|X|X|}\hline
\rowcolor{popblueL} English & Français & Exemple / Remarque \\ \hline
climate change & le changement climatique & \eng{Climate change affects us all.} \\ \hline
global warming & le réchauffement planétaire & \eng{Global warming is accelerating.} \\ \hline
greenhouse effect & l'effet de serre & \eng{The greenhouse effect traps heat.} \\ \hline
greenhouse gases (GHG) & les gaz à effet de serre & \eng{We must cut greenhouse gases.} \\ \hline
carbon dioxide (CO2) & le dioxyde de carbone & \eng{Trees absorb carbon dioxide.} \\ \hline
emissions & les émissions, les rejets & \eng{Car emissions pollute the air.} \\ \hline
carbon footprint & l'empreinte carbone & \eng{Reduce your carbon footprint!} \\ \hline
fossil fuels & les énergies fossiles & coal, oil, gas \\ \hline
exhaust fumes & les gaz d'échappement & \eng{Exhaust fumes cause diseases.} \\ \hline
rising temperatures & la hausse des températures & adjectif en \emph{-ing} \\ \hline
heatwave & la canicule & un seul mot \\ \hline
drought & la sécheresse & prononcé « draoute » \\ \hline
flood & l'inondation & prononcé « fleud » \\ \hline
melting ice caps & la fonte des calottes glaciaires & verbe \eng{to melt} \\ \hline
rising sea levels & la montée du niveau des mers & \eng{Sea levels are rising.} \\ \hline
desertification & la désertification & nom en \emph{-tion} \\ \hline
deforestation & la déforestation & verbe \eng{to deforest} \\ \hline
industrial pollution & la pollution industrielle & \eng{Industrial pollution kills fish.} \\ \hline
overconsumption & la surconsommation & \eng{Overconsumption wastes resources.} \\ \hline
to burn / to release & brûler / rejeter & \eng{Burning coal releases CO2.} \\ \hline
\end{tabularx}
\end{center}

\courssub{Weather ou climate ? Une distinction capitale}

\begin{attention}[Ne confonds pas weather et climate]
\eng{Weather} désigne \cle{le temps qu'il fait} : la pluie, le soleil, le vent \emph{aujourd'hui ou cette semaine} (court terme). \eng{Climate} désigne \cle{le climat} d'une région, c'est-à-dire les conditions moyennes observées \emph{sur des années ou des décennies} (long terme).

\eng{The weather is hot and dry in Maroua today.} « Le temps est chaud et sec à Maroua aujourd'hui. »

\eng{The climate of the Far North is becoming drier every decade.} « Le climat de l'Extrême-Nord devient plus sec à chaque décennie. »

Phrase mémorable : \eng{Climate is what you expect; weather is what you get.} « Le climat, c'est ce à quoi on s'attend ; la météo, c'est ce qu'on reçoit. »

Autre piège : \eng{weather} est \emph{incomptable} — on ne dit jamais \emph{a weather} ni \emph{weathers}. On dit \eng{The weather is fine today.}
\end{attention}

\begin{center}
\begin{tabularx}{\linewidth}{|X|X|X|}\hline
\rowcolor{popblueL} & Weather & Climate \\ \hline
Durée & court terme (un jour, une semaine) & long terme (années, décennies) \\ \hline
Français & le temps (qu'il fait), la météo & le climat \\ \hline
Exemple & \eng{It is raining in Buea this morning.} & \eng{Buea has a cool, wet climate.} \\ \hline
Grammaire & nom incomptable & nom comptable : \eng{a tropical climate} \\ \hline
\end{tabularx}
\end{center}

\courssub{Prononciation : les mots pièges}

\begin{propriete}[Prononciation et orthographe]
\begin{itemize}
\item \eng{climate} se prononce « klaï-mit » : le \emph{-mate} final ne se prononce jamais « mat » comme en français.
\item \eng{drought} se prononce « draoute » : le groupe \emph{-ough} se prononce « aou » ici. Attention à l'orthographe : \emph{d-r-o-u-g-h-t}.
\item \eng{flood} se prononce « fleud » : le double \emph{oo} sonne « eu » (comme dans \eng{blood}), jamais « ou ».
\item \eng{earth} se prononce « eurth » ; avec majuscule (\eng{the Earth}) quand on parle de la planète.
\item \eng{atmosphere} se prononce « at-mo-sfieur », avec l'accent tonique sur la première syllabe.
\end{itemize}
\end{propriete}

\courssub{Carte mentale : organiser le champ lexical}

\begin{popfigure}
\begin{center}
\begin{tikzpicture}[mindmap, grow cyclic,
  every node/.style={concept, text=white, align=center},
  root concept/.append style={concept color=popblue, font=\large\bfseries},
  level 1/.append style={level distance=3.6cm, sibling angle=120},
  level 2/.append style={level distance=2.6cm, sibling angle=45},
  cause/.style={concept color=poporange},
  effect/.style={concept color=poppurple},
  key/.style={concept color=popgreen}]
\node[root concept]{CLIMATE CHANGE}
  child[cause]{ node {Causes}
    child { node[align=center]{fossil fuels\\ coal, oil, gas} }
    child { node {deforestation} }
    child { node {exhaust fumes} }
  }
  child[effect]{ node {Effects}
    child { node {drought} }
    child { node {floods} }
    child { node[align=center]{rising sea\\ levels} }
  }
  child[key]{ node {Key words}
    child { node[align=center]{greenhouse\\ effect} }
    child { node[align=center]{CO2\\ emissions} }
    child { node[align=center]{carbon\\ footprint} }
  };
\end{tikzpicture}
\end{center}
\legende{Carte mentale du champ lexical de \eng{climate change} : causes, effets et mots-clés.}
\end{popfigure}

\courssub{Exercice résolu et entraînement}

\begin{exoresolu}[Complète avec le mot qui convient]
Énoncé. Complète chaque phrase avec un mot de la liste : \eng{drought}, \eng{emissions}, \eng{carbon footprint}, \eng{floods}, \eng{weather}, \eng{climate}.

1. The \ldots\ in the Far North lasted eight months and the rivers dried up.
2. We must reduce CO2 \ldots\ from cars and factories.
3. Last week, heavy rain caused terrible \ldots\ in Douala.
4. The \ldots\ of the South West Region is wet and cool.
5. Planting trees is a good way to reduce your \ldots.
6. What is the \ldots\ like in Yaoundé today?
\tcblower
Solution détaillée.
1. \eng{drought} — une période de huit mois sans pluie est une sécheresse.
2. \eng{emissions} — on « réduit les émissions » de CO2 ; le mot est au pluriel.
3. \eng{floods} — la pluie violente provoque des inondations.
4. \eng{climate} — il s'agit des conditions durables d'une région (long terme), pas du temps d'aujourd'hui.
5. \eng{carbon footprint} — planter des arbres réduit l'empreinte carbone.
6. \eng{weather} — « le temps qu'il fait aujourd'hui » = \eng{weather} ; la question-type est \eng{What is the weather like?}
\end{exoresolu}

\begin{exercice}[À toi de jouer]
1. Traduis en anglais : « La déforestation et les énergies fossiles sont les principales causes du réchauffement planétaire. »
2. Complète : \eng{The ice caps are m\ldots\ because of rising t\ldots.}
3. Choisis le bon mot : \eng{The (weather / climate) in the Sahel is becoming hotter and drier.}
4. Trouve le mot : un gaz produit par les voitures et les usines, absorbé par les arbres.
5. Rédige deux phrases sur ta ville avec \eng{flood} et \eng{heatwave}.
\end{exercice}

\begin{corrige}[Exercice « À toi de jouer »]
1. \eng{Deforestation and fossil fuels are the main causes of global warming.}
2. \eng{The ice caps are melting because of rising temperatures.}
3. \eng{climate} (long terme, tendance durable).
4. \eng{carbon dioxide (CO2)}.
5. Exemple : \eng{In Douala, floods often block the roads during the rainy season. Last year, a terrible heatwave made life difficult in Garoua.}
\end{corrige}

\begin{savaistu}[Why “greenhouse”?]
The term \eng{greenhouse effect} compares the Earth to a glass greenhouse that keeps heat inside: sunlight passes through the glass, but the heat cannot escape. In the same way, greenhouse gases let the sun's rays in but trap the heat around our planet. Gardeners in England have used glass greenhouses for centuries to grow tomatoes and flowers in cold weather — scientists simply borrowed the image! And remember the slogan of young climate activists, from Yaoundé to London: \eng{There is no planet B.}
\end{savaistu}

\begin{aretenir}[L'essentiel de la section]
\begin{itemize}
\item \eng{Global warming} (le réchauffement) est la hausse des températures ; \eng{climate change} (le changement climatique) en est la conséquence globale.
\item Causes : \eng{burning fossil fuels}, \eng{exhaust fumes}, \eng{deforestation}, \eng{overconsumption}.
\item Effets : \eng{drought}, \eng{floods}, \eng{heatwaves}, \eng{melting ice caps}, \eng{rising sea levels}, \eng{desertification}.
\item \eng{Weather} = le temps qu'il fait (court terme, incomptable) ; \eng{climate} = le climat (long terme).
\item Prononciation : \eng{climate} « klaï-mit », \eng{drought} « draoute », \eng{flood} « fleud ».
\end{itemize}
\end{aretenir}

\coursec{Champ lexical 2 : Pollution and waste}

Dans la section précédente, nous avons découvert le vocabulaire du changement climatique et du réchauffement de la planète (\eng{climate change}, \eng{global warming}, \eng{greenhouse gases}). Mais le réchauffement climatique n'est pas le seul problème environnemental : chaque jour, à Douala comme à Yaoundé, nous voyons des sachets plastiques dans les rues, de la fumée noire qui sort des pots d'échappement des \eng{bendskins} et des usines qui rejettent des produits chimiques dans les rivières. Pour parler de tout cela en anglais, il faut maîtriser un deuxième champ lexical essentiel : celui de la \cle{pollution} et des \cle{déchets} (\eng{waste}).

\courssub{Observation : un texte pour découvrir le vocabulaire}

Lisons d'abord ce court article de presse original. Souligne mentalement tous les mots qui désignent la pollution ou les déchets.

\begin{textebox}[The Wouri River in Danger — adapted from a local newspaper]
Every day, tonnes of plastic waste are dumped into the Wouri River. This litter kills fish and contaminates drinking water. In Douala, the biggest sources of pollution are easy to identify. First, factories dump chemicals and used oil into the river instead of treating them. Second, old vehicles and taxis produce thick exhaust fumes, which cause serious air pollution in the city centre. Third, many inhabitants throw away their rubbish in the streets because there are not enough dustbins. During the rainy season, this litter blocks the gutters and causes floods.

Water pollution is not the only problem. Soil pollution is increasing around the city: farmers use too many pesticides and fertilisers, and these chemicals contaminate the ground where vegetables grow. Noise pollution is also getting worse because of loud music in bars and the constant traffic.

Fortunately, solutions exist. Some young people organise clean-up campaigns every Saturday morning. They collect plastic bags and bottles and take them to a recycling centre. A local company now recycles single-use plastics and transforms them into paving stones. “We must stop dumping our waste everywhere,” says Marie, a volunteer. “If everyone recycles and stops littering, Douala will become a clean and healthy city.” The city council has also promised to build a modern landfill outside the city and to punish illegal dumping with heavy fines.
\end{textebox}

\textbf{Glossaire}

\begin{center}
\begin{tabularx}{\linewidth}{|l|l|X|}
\hline
\rowcolor{popblueL} \textbf{Mot anglais} & \textbf{Nature} & \textbf{Traduction française} \\
\hline
waste & nom (indénombrable) & déchets \\
\hline
to dump & verbe & déverser, jeter (illégalement) \\
\hline
litter & nom (indénombrable) & détritus (jetés dans la rue ou la nature) \\
\hline
to contaminate & verbe & contaminer, polluer (l'eau, le sol) \\
\hline
exhaust fumes & nom pluriel & gaz d'échappement \\
\hline
rubbish & nom (GB, indénombrable) & ordures, déchets \\
\hline
gutters & nom pluriel & caniveaux \\
\hline
pesticides / fertilisers & noms pluriels & pesticides / engrais \\
\hline
single-use plastics & nom pluriel & plastiques à usage unique \\
\hline
to recycle & verbe & recycler \\
\hline
landfill & nom & décharge (contrôlée) \\
\hline
fine & nom & amende \\
\hline
\end{tabularx}
\end{center}

\textbf{Questions de compréhension}
\begin{enumerate}
\item What are the three biggest sources of pollution in Douala according to the text?
\item What happens to the litter during the rainy season?
\item What do the volunteers do every Saturday morning?
\item What has the city council promised to do?
\end{enumerate}

\courssub{Les types de pollution}

\begin{definition}[Types of pollution]
En anglais, le mot \cle{pollution} est généralement \textbf{indénombrable} quand il désigne le phénomène en général, mais on le précise par un complément : \eng{air pollution} (pollution de l'air), \eng{water pollution} (pollution de l'eau), \eng{soil pollution} (pollution du sol), \eng{noise pollution} (pollution sonore), \eng{plastic pollution} (pollution plastique).
\end{definition}

\begin{exemplebox}[Les cinq types de pollution en contexte]
\begin{itemize}
\item \eng{Air pollution in Douala is caused by old cars and factories.} — La pollution de l'air à Douala est causée par les vieilles voitures et les usines.
\item \eng{Water pollution kills fish in the Wouri River.} — La pollution de l'eau tue les poissons du fleuve Wouri.
\item \eng{Soil pollution comes from pesticides and chemical fertilisers.} — La pollution du sol vient des pesticides et des engrais chimiques.
\item \eng{Noise pollution from bars keeps the neighbours awake at night.} — La pollution sonore des bars empêche les voisins de dormir la nuit.
\item \eng{Plastic pollution is a serious problem on the beaches of Kribi.} — La pollution plastique est un problème sérieux sur les plages de Kribi.
\end{itemize}
\end{exemplebox}

\begin{attention}[\eng{Pollution} est indénombrable]
On ne dit jamais \emph{a pollution} ni \emph{pollutions} au sens général. On dit : \eng{Pollution is a serious problem.} (La pollution est un problème sérieux.) Pour individualiser, utilisez \eng{a type of pollution} ou \eng{a form of pollution} : \eng{Noise pollution is a form of pollution that people often forget.}
\end{attention}

\courssub{Les sources de pollution}

Chaque type de pollution a ses \cle{sources}. Voici les mots essentiels :

\begin{definition}[Sources of pollution]
\begin{itemize}
\item \cle{factories} : les usines (rejettent fumées et produits chimiques) ;
\item \cle{vehicles} : les véhicules (voitures, camions, motos — ils produisent des \eng{exhaust fumes}, gaz d'échappement) ;
\item \cle{pesticides} et \cle{fertilisers} : pesticides et engrais (agriculture intensive) ;
\item \cle{oil spills} : marées noires (déversements accidentels de pétrole en mer) ;
\item \cle{dumping waste} : le dépôt sauvage de déchets.
\end{itemize}
\end{definition}

\begin{exemplebox}[Sources en contexte]
\begin{itemize}
\item \eng{Factories must not dump chemicals into rivers.} — Les usines ne doivent pas déverser de produits chimiques dans les rivières.
\item \eng{Exhaust fumes from vehicles pollute the air we breathe.} — Les gaz d'échappement des véhicules polluent l'air que nous respirons.
\item \eng{An oil spill off the coast of Limbe destroyed the fishing grounds.} — Une marée noire au large de Limbe a détruit les zones de pêche.
\item \eng{Farmers should use fewer pesticides and chemical fertilisers.} — Les agriculteurs devraient utiliser moins de pesticides et d'engrais chimiques.
\end{itemize}
\end{exemplebox}

La figure suivante associe chaque type de pollution à sa source principale :

\begin{popfigure}
\begin{center}
\begin{tikzpicture}[
  box/.style={draw=popblue, fill=popblueL, rounded corners=2pt, align=center, font=\small, minimum width=3.4cm, minimum height=0.8cm},
  src/.style={draw=poporange, fill=poporangeL, rounded corners=2pt, align=center, font=\small, minimum width=3.4cm, minimum height=0.8cm},
  link/.style={-{Stealth[length=2mm]}, thick, popdark}
]
\node[box] (air) at (0,3) {\textbf{Air pollution}};
\node[box] (water) at (0,1.5) {\textbf{Water pollution}};
\node[box] (soil) at (0,0) {\textbf{Soil pollution}};
\node[box] (noise) at (0,-1.5) {\textbf{Noise pollution}};
\node[src] (cars) at (7,3) {vehicles, exhaust fumes};
\node[src] (fact) at (7,1.5) {factories, chemicals, oil spills};
\node[src] (chem) at (7,0) {pesticides, fertilisers};
\node[src] (loud) at (7,-1.5) {loudspeakers, traffic};
\draw[link] (air) -- (cars);
\draw[link] (water) -- (fact);
\draw[link] (soil) -- (chem);
\draw[link] (noise) -- (loud);
\node[font=\small\bfseries, popblue] at (0,4) {Types of pollution};
\node[font=\small\bfseries, poporange] at (7,4) {Sources};
\end{tikzpicture}
\end{center}
\legende{Chaque type de pollution et sa source principale}
\end{popfigure}

\courssub{Le lexique des déchets : rubbish, garbage, trash, litter, waste}

C'est ici que les francophones se trompent le plus souvent, car l'anglais distingue plusieurs mots là où le français n'a que « déchets » ou « ordures ».

\begin{definition}[Waste, rubbish, garbage, trash, litter]
\begin{itemize}
\item \cle{waste} : terme général et soutenu, indénombrable — « déchets » (\eng{household waste}, déchets ménagers ; \eng{toxic waste}, déchets toxiques) ;
\item \cle{rubbish} : ordures, déchets ménagers — mot \textbf{britannique}, indénombrable ;
\item \cle{garbage} et \cle{trash} : ordures — mots \textbf{américains}, indénombrables ;
\item \cle{litter} : détritus jetés \textbf{dans la rue ou la nature} (sachets, canettes, papiers) — indénombrable ;
\item \cle{landfill} : décharge, site d'enfouissement des déchets ;
\item \cle{plastic bags} : sachets plastiques ; \cle{single-use plastics} : plastiques à usage unique.
\end{itemize}
\end{definition}

\begin{attention}[Rubbish (GB) ou garbage/trash (US) ?]
\eng{Rubbish} est le mot \textbf{britannique} ; \eng{garbage} et \eng{trash} sont \textbf{américains}. Un Anglais dit \eng{Put the rubbish in the bin}, un Américain dit \eng{Put the trash in the garbage can}. Attention à \eng{litter} : il désigne uniquement les déchets abandonnés dans l'espace public. Ne traduisez donc pas « les ordures » par \emph{litters} : \eng{litter} est indénombrable et n'a pas de pluriel.
\end{attention}

\begin{center}
\begin{tabularx}{\linewidth}{|X|X|X|}
\hline
\rowcolor{popblueL} \textbf{Français} & \textbf{English} & \textbf{Remarque} \\
\hline
les déchets (terme général) & waste & indénombrable, registre soutenu \\
\hline
les ordures ménagères & rubbish (GB) / garbage, trash (US) & indénombrables dans les deux variétés \\
\hline
les détritus (rue, nature) & litter & indénombrable ; verbe : to litter \\
\hline
la poubelle & dustbin / bin (GB) ; garbage can / trash can (US) & nom dénombrable \\
\hline
la décharge & landfill & aussi « rubbish dump » (GB) \\
\hline
les sachets plastiques & plastic bags & toujours au pluriel en général \\
\hline
les plastiques à usage unique & single-use plastics & paille, gobelet, sachet \\
\hline
jeter (aux ordures) & to throw away & verbe à particule \\
\hline
\end{tabularx}
\end{center}

\courssub{Les verbes associés}

\begin{definition}[Verbes du champ lexical]
\begin{itemize}
\item \cle{to pollute} : polluer (le verbe le plus général) ;
\item \cle{to contaminate} : contaminer (rendre dangereux l'eau, le sol, les aliments — plus précis et plus fort que \eng{pollute}) ;
\item \cle{to dump} : déverser, jeter illégalement (déchets, produits chimiques) ;
\item \cle{to throw away} : jeter (aux ordures) — verbe à particule, séparable : \eng{throw it away} ;
\item \cle{to litter} : jeter des détritus dans l'espace public ;
\item \cle{to recycle} : recycler.
\end{itemize}
\end{definition}

\begin{exemplebox}[Verbes en contexte]
\begin{itemize}
\item \eng{Exhaust fumes pollute the atmosphere.} — Les gaz d'échappement polluent l'atmosphère.
\item \eng{The oil spill contaminated the drinking water of the whole village.} — La marée noire a contaminé l'eau potable de tout le village.
\item \eng{It is forbidden to dump waste in the forest.} — Il est interdit de déverser des déchets dans la forêt.
\item \eng{Don't throw away plastic bottles: recycle them!} — Ne jetez pas les bouteilles plastiques : recyclez-les !
\item \eng{People who litter the streets should pay a fine.} — Les gens qui jettent des détritus dans les rues devraient payer une amende.
\end{itemize}
\end{exemplebox}

\begin{methode}[Mémoriser par association : type + verbe + source]
Pour retenir durablement ce vocabulaire, construisez des « trios » mentaux qui relient un type de pollution, un verbe et une source :
\begin{enumerate}
\item Choisissez un type de pollution : \eng{air pollution}.
\item Associez-lui un verbe : \eng{to pollute the air}.
\item Ajoutez une source concrète : \eng{exhaust fumes from cars}.
\item Formulez une phrase complète : \eng{Exhaust fumes from cars pollute the air.}
\end{enumerate}
Autres trios : \eng{water pollution — to contaminate — factories dumping chemicals} ; \eng{soil pollution — to contaminate — pesticides and fertilisers} ; \eng{noise pollution — to disturb — loudspeakers and traffic}. Répétez chaque trio à voix haute, puis fermez le cahier et réécrivez les phrases de mémoire.
\end{methode}

\courssub{Collocations et expressions idiomatiques}

\begin{definition}[Collocations usuelles sur l'environnement]
Certains mots vont systématiquement ensemble. Les apprendre par blocs évite les calques du français.
\begin{itemize}
\item \cle{to reduce waste} (réduire les déchets) — pas \emph{decrease waste} ;
\item \cle{to sort waste} / \cle{to separate waste} (trier les déchets) ;
\item \cle{to dispose of waste} (se débarrasser des déchets / éliminer les déchets) — \textbf{of} obligatoire ;
\item \cle{carbon footprint} (empreinte carbone) ;
\item \cle{renewable energy} (énergie renouvelable) ;
\item \cle{endangered species} (espèces menacées) ;
\item \cle{to raise awareness} (sensibiliser) — \eng{raise awareness about plastic pollution}.
\end{itemize}
\end{definition}

\begin{exemplebox}[Collocations en phrases]
\begin{itemize}
\item \eng{We must \textbf{reduce our carbon footprint} by using public transport.} — Nous devons réduire notre empreinte carbone en utilisant les transports en commun.
\item \eng{The council encourages people to \textbf{sort their waste} before collection.} — Le conseil municipal encourage les gens à trier leurs déchets avant la collecte.
\item \eng{It is illegal to \textbf{dispose of} toxic waste in the river.} — Il est illégal de se débarrasser de déchets toxiques dans la rivière.
\end{itemize}
\end{exemplebox}

\courssub{Familles de mots, prononciation et orthographe}

\begin{propriete}[Familles de mots (Word families)]
Construire le nom, le verbe, l'adjectif et l'adverbe à partir d'une racine permet d'enrichir son lexique rapidement.

\begin{center}
\begin{tabular}{|l|l|l|l|l|}
\hline
\rowcolor{popblueL} \textbf{Verbe} & \textbf{Nom (chose/action)} & \textbf{Nom (agent)} & \textbf{Adjectif} & \textbf{Adverbe} \\
\hline
pollute & pollution & polluter & polluted / polluting & pollutedly (rare) \\
\hline
contaminate & contamination & contaminant & contaminated & — \\
\hline
recycle & recycling / recycler & recycler & recycled / recyclable & — \\
\hline
conserve & conservation & conservationist & conserved / conservational & — \\
\hline
degrade & degradation & — & degradable / biodegradable & — \\
\hline
sustain & sustainability & — & sustainable & sustainably \\
\hline
\end{tabular}
\end{center}
\end{propriete}

\begin{propriete}[Trois mots à orthographier avec soin]
\begin{itemize}
\item \cle{environment} : attention au \textbf{n} avant le \textbf{m} — en-vi-ron-\textbf{n}-ment. L'erreur classique des francophones est d'écrire \emph{enviroment}. Prononciation : « in-vaïe-ron-ment ».
\item \cle{pollution} : un seul \textbf{l} au début mais double \textbf{l} ? Non : p-o-\textbf{ll}-u-tion, deux \textbf{l} comme en français. Ne pas confondre avec le verbe \eng{to pollute} (deux \textbf{l} aussi).
\item \cle{contamination} : c-o-n-t-a-m-i-n-a-tion — attention à la voyelle \textbf{a} après \eng{cont} (pas \emph{contemination}).
\item \cle{biodegradable} : bio-de-grad-a-ble — ne pas oublier le \textbf{a} après \eng{grad}. Prononciation : « baïo-di-gré-da-beul ».
\end{itemize}
\end{propriete}

\begin{savaistu}[Le Cameroun contre les sachets plastiques]
Cameroon banned the production and importation of non-biodegradable plastic bags in 2014. — Le Cameroun a interdit la production et l'importation des sachets plastiques non biodégradables en 2014. Cette mesure visait à réduire la \eng{plastic pollution}, car les sachets abandonnés bouchent les caniveaux (\eng{gutters}) et provoquent des inondations (\eng{floods}) pendant la saison des pluies. Le mot \eng{biodegradable} se prononce « baïo-di-gré-da-beul » : il vient du préfixe grec \emph{bio-} (vie) et du verbe latin \emph{degradare} (dégrader).
\end{savaistu}

\begin{exoresolu}[Complétez avec le mot qui convient]
Énoncé — Complétez chaque phrase avec le mot correct choisi dans la liste : \eng{litter}, \eng{rubbish}, \eng{landfill}, \eng{contaminated}, \eng{dump}, \eng{recycle}.
\begin{enumerate}
\item In Britain, people put their \ldots\ in the dustbin twice a week.
\item Please do not drop \ldots\ in the school playground.
\item The factory was fined because it continued to \ldots\ chemicals into the river.
\item The oil spill \ldots\ the drinking water of the village.
\item Old fridges and cars are taken to a \ldots\ outside the city.
\item We should \ldots\ glass bottles instead of throwing them away.
\end{enumerate}
\tcblower
Solution détaillée :
\begin{enumerate}
\item \eng{rubbish} — contexte britannique (\eng{dustbin}, \eng{Britain}) : le mot américain serait \eng{trash} ou \eng{garbage}.
\item \eng{litter} — il s'agit de détritus jetés dans un espace public (la cour de l'école).
\item \eng{dump} — déverser illégalement des produits chimiques ; la construction est \eng{to dump something into a river}.
\item \eng{contaminated} — passé simple du verbe \eng{to contaminate} : l'eau potable a été rendue dangereuse.
\item \eng{landfill} — une décharge contrôlée où l'on enfouit les gros déchets.
\item \eng{recycle} — la structure \eng{instead of + V-ing} (\eng{throwing}) impose l'infinitif après \eng{should} : \eng{recycle}.
\end{enumerate}
\end{exoresolu}

\begin{attention}[Pièges des francophones]
\begin{itemize}
\item Ne dites pas \emph{to pollute the water with wastes} : \eng{waste} est indénombrable. Dites \eng{to pollute the water with waste}.
\item \eng{To throw away} est séparable : \eng{Throw it away!} et jamais \emph{Throw away it!}
\item Le français « une décharge sauvage » se dit \eng{an illegal dump} ou \eng{illegal dumping} (l'action), pas \emph{a wild discharge}.
\item Attention au faux ami : \eng{to contaminate} s'emploie pour l'eau, le sol, les aliments, mais pas pour l'air en général — pour l'air, préférez \eng{to pollute}.
\item \eng{Dispose} ne s'emploie jamais seul pour « jeter » : il faut \eng{dispose \textbf{of}}.
\end{itemize}
\end{attention}

\begin{exercice}[À vous de jouer]
\begin{enumerate}
\item Traduisez en anglais : « Les usines ne doivent pas déverser de produits chimiques dans les rivières. »
\item Complétez : \eng{People who \ldots\ in the streets must pay a fine.} (jeter des détritus)
\item Choisissez la variété britannique : \eng{garbage} / \eng{rubbish} / \eng{trash}.
\item Construisez trois phrases complètes avec les trios : \eng{soil pollution — to contaminate — pesticides} ; \eng{air pollution — to pollute — exhaust fumes} ; \eng{water pollution — to dump — factories}.
\item Mettez les mots de la même famille dans le bon ordre (verbe $\to$ nom action $\to$ adjectif) : \eng{sustain, sustainability, sustainable} ; \eng{degrade, degradation, degradable}.
\end{enumerate}
\end{exercice}

\begin{corrige}[Exercice « À vous de jouer »]
\begin{enumerate}
\item \eng{Factories must not dump chemicals into rivers.}
\item \eng{litter} — \eng{People who litter in the streets must pay a fine.}
\item \eng{rubbish} est le mot britannique ; \eng{garbage} et \eng{trash} sont américains.
\item Réponses possibles : \eng{Pesticides contaminate the soil.} / \eng{Exhaust fumes pollute the air.} / \eng{Factories dump waste into rivers and cause water pollution.}
\item \eng{sustain $\to$ sustainability $\to$ sustainable} ; \eng{degrade $\to$ degradation $\to$ degradable} (ou \eng{biodegradable}).
\end{enumerate}
\end{corrige}

\begin{aretenir}[L'essentiel de la section]
\begin{itemize}
\item Cinq types de pollution : \eng{air}, \eng{water}, \eng{soil}, \eng{noise}, \eng{plastic pollution} — \eng{pollution} est indénombrable.
\item Sources : \eng{factories}, \eng{vehicles} (et leurs \eng{exhaust fumes}), \eng{pesticides}, \eng{fertilisers}, \eng{oil spills}, \eng{dumping waste}.
\item Déchets : \eng{waste} (général), \eng{rubbish} (GB) / \eng{garbage, trash} (US), \eng{litter} (détritus dans l'espace public), \eng{landfill} (décharge) — tous indénombrables sauf \eng{landfill}.
\item Verbes : \eng{to pollute}, \eng{to contaminate}, \eng{to dump}, \eng{to throw away} (séparable), \eng{to litter}, \eng{to recycle}, \eng{to dispose of}.
\item Collocations clés : \eng{reduce waste}, \eng{sort waste}, \eng{carbon footprint}, \eng{renewable energy}, \eng{raise awareness}.
\item Familles de mots : \eng{pollute/pollution}, \eng{contaminate/contamination}, \eng{recycle/recycling}, \eng{sustain/sustainability}, \eng{degrade/degradation}.
\item Orthographe : \eng{environment} (avec \textbf{n} avant \textbf{m}), \eng{pollution}, \eng{contamination}, \eng{biodegradable}.
\end{itemize}
\end{aretenir}

\coursec{Champ lexical 3 : Solutions and protection}

Dans la section précédente, nous avons étudié le vocabulaire de la \cle{pollution} et des \cle{déchets} : \eng{air pollution}, \eng{plastic waste}, \eng{toxic fumes}, \eng{landfills}... Nous avons donc dressé le diagnostic du problème. Mais une fois le problème identifié, une question s'impose naturellement : \emph{what can we do about it?} (que pouvons-nous y faire ?). C'est précisément l'objet de cette troisième section : le vocabulaire des \cle{solutions} et de la \cle{protection} de l'environnement. C'est le lexique que l'on retrouve dans les discours officiels, les campagnes de sensibilisation, les accords internationaux et les articles de presse anglophones — autant de documents que vous rencontrerez à l'épreuve d'anglais du baccalauréat.

\courssub{Observation : un texte pour découvrir le lexique en contexte}

Lisons d'abord ce court article de presse original. Soulignez mentalement tous les mots qui expriment une \emph{solution} ou une \emph{action positive} pour l'environnement.

\begin{textebox}[Going green in Bamenda — a local newspaper article]
Students in Bamenda launched an awareness campaign last week to encourage people to plant trees and use solar lamps. The campaign, called “Green Hills”, was organised by the school's environmental club with the support of two local NGOs and the city council.

“Our goal is to raise awareness about climate change,” explains Mirabel, the club's president. “Many people cut down trees for firewood, but nobody replaces them. If we want to protect the environment, we must take action now.”

During the campaign, students distributed biodegradable bags to traders at the main market and taught shopkeepers how to reduce, reuse and recycle plastic waste. They also persuaded the city council to ban single-use plastics in public buildings. “Small actions matter,” says Mr Ndi, a science teacher. “If every family saves energy, plants one tree and avoids burning plastic, our town will become cleaner and healthier.”

The regional government has promised to support the project. Officials announced new measures to preserve the forest around the city and to promote renewable energy, especially solar power, which is cheap and abundant in the region. Environmental activists welcomed the decision but warned that laws alone are not enough. “The government can pass laws,” one activist said, “but real change depends on eco-friendly consumers and local communities.”

The “Green Hills” campaign shows that sustainable development is not only a job for presidents and ministers. It starts in schools, markets and homes — with ordinary people who choose to act.
\end{textebox}

\textbf{Glossaire (mots difficiles) :}

\begin{center}
\begin{tabularx}{\linewidth}{|l|l|X|}
\hline
\rowcolor{popblueL} Mot anglais & Nature & Traduction française \\
\hline
to launch & verbe & lancer (une campagne) \\
\hline
awareness campaign & nom & campagne de sensibilisation \\
\hline
firewood & nom & bois de chauffe \\
\hline
to take action & expression & agir, passer à l'action \\
\hline
single-use plastics & nom & plastiques à usage unique \\
\hline
to ban & verbe & interdire, bannir \\
\hline
to pass laws & expression & voter / adopter des lois \\
\hline
city council & nom & conseil municipal, mairie \\
\hline
\end{tabularx}
\end{center}

\textbf{Questions de compréhension :}
\begin{enumerate}
\item What is the name of the campaign, and who organised it?
\item What two actions does Mirabel say people must take?
\item What did the regional government promise to do?
\item According to the activist, why are laws not enough?
\end{enumerate}

\courssub{Le vocabulaire des actions : les verbes de la protection}

Le texte contient les verbes essentiels du champ lexical des solutions. Observons-les en contexte avant de les systématiser :

\begin{itemize}
\item \eng{to protect the environment} — protéger l'environnement ;
\item \eng{to preserve the forest} — préserver la forêt ;
\item \eng{to save energy} — économiser l'énergie ;
\item \eng{to reduce, reuse and recycle plastic waste} — réduire, réutiliser et recycler les déchets plastiques ;
\item \eng{to plant trees} — planter des arbres ;
\item \eng{to ban single-use plastics} — interdire les plastiques à usage unique.
\end{itemize}

\begin{propriete}[Protect, preserve, save : trois verbes à ne pas confondre]
Ces trois verbes sont proches mais ne s'emploient pas dans les mêmes contextes :
\begin{itemize}
\item \cle{to protect} = protéger contre un danger : \eng{We must protect endangered species.} (Nous devons protéger les espèces menacées.)
\item \cle{to preserve} = préserver, maintenir en l'état, garder intact : \eng{We must preserve forests for future generations.} (Nous devons préserver les forêts pour les générations futures.)
\item \cle{to save} = économiser, ne pas gaspiller (énergie, eau, argent) ou sauver (une vie, la planète) : \eng{Switch off the lights to save energy.} (Éteins les lumières pour économiser l'énergie.)
\end{itemize}
\end{propriete}

\begin{exemplebox}[Exemples traduits]
\eng{We must reduce our carbon footprint by using public transport.} — « Nous devons réduire notre empreinte carbone en utilisant les transports en commun. »

\eng{The school reuses old paper for rough work.} — « L'école réutilise le vieux papier pour le brouillon. »

\eng{In Douala, some companies recycle plastic bottles into building materials.} — « À Douala, certaines entreprises recyclent les bouteilles en plastique pour en faire des matériaux de construction. »

\eng{The government has banned plastic bags in supermarkets.} — « Le gouvernement a interdit les sacs plastiques dans les supermarchés. »
\end{exemplebox}

\begin{definition}[The 3 Rs : Reduce, Reuse, Recycle]
L'expression \cle{the 3 Rs} désigne les trois principes fondamentaux de la gestion des déchets : \cle{Reduce} (réduire sa consommation à la source), \cle{Reuse} (réutiliser les objets au lieu de les jeter), \cle{Recycle} (recycler les matériaux pour fabriquer de nouveaux produits). On dit aussi \eng{the three Rs of waste management}.
\end{definition}

\begin{popfigure}
\begin{tikzpicture}[>=Stealth, node distance=2.6cm]
\node[draw=popgreen, fill=popgreenL, circle, minimum size=1.9cm, align=center, font=\bfseries] (reduce) at (90:2.3) {Reduce};
\node[draw=poporange, fill=poporangeL, circle, minimum size=1.9cm, align=center, font=\bfseries] (reuse) at (210:2.3) {Reuse};
\node[draw=popblue, fill=popblueL, circle, minimum size=1.9cm, align=center, font=\bfseries] (recycle) at (330:2.3) {Recycle};
\draw[->, very thick, popgreen, bend left=25] (reduce) to (reuse);
\draw[->, very thick, poporange, bend left=25] (reuse) to (recycle);
\draw[->, very thick, popblue, bend left=25] (recycle) to (reduce);
\node[align=center, font=\itshape\small, text=popdark] at (0,0) {The 3 Rs\\ of waste\\ management};
\end{tikzpicture}
\legende{Le cycle circulaire des 3 Rs : Reduce, Reuse, Recycle.}
\end{popfigure}

\courssub{Les énergies : renouvelables contre fossiles}

Le texte mentionne \eng{solar power} (l'énergie solaire). Voici le vocabulaire complet des énergies, organisé en deux camps opposés :

\begin{center}
\begin{tabularx}{\linewidth}{|X|X|}
\hline
\rowcolor{popgreenL} Renewable energy (énergies renouvelables) & \cellcolor{poporangeL} Non-renewable energy (énergies non renouvelables) \\
\hline
solar power / solar energy (énergie solaire) & fossil fuels (énergies fossiles) \\
\hline
wind power (énergie éolienne) & coal (charbon) \\
\hline
hydropower (hydroélectricité) & oil / petroleum (pétrole) \\
\hline
geothermal energy (géothermie) & natural gas (gaz naturel) \\
\hline
biomass (biomasse) & nuclear energy (énergie nucléaire) \\
\hline
\end{tabularx}
\end{center}

\begin{exemplebox}[Exemples traduits]
\eng{Cameroon has great potential for hydropower thanks to its many rivers.} — « Le Cameroun a un grand potentiel hydroélectrique grâce à ses nombreux fleuves. »

\eng{Solar lamps are replacing kerosene lamps in many villages.} — « Les lampes solaires remplacent les lampes à pétrole dans de nombreux villages. »

\eng{Burning fossil fuels releases greenhouse gases into the atmosphere.} — « La combustion des énergies fossiles libère des gaz à effet de serre dans l'atmosphère. »
\end{exemplebox}

\begin{attention}[Attention à la prononciation et à la formation des mots]
\begin{itemize}
\item \eng{renewable} se prononce « ri-niou-e-beul » ; le préfixe \cle{re-} exprime le renouvellement, comme dans \eng{reuse} et \eng{recycle}.
\item \eng{non-renewable} prend un trait d'union : le préfixe \cle{non-} marque la négation.
\item Ne dites pas \eng{*solar energy power} : choisissez \eng{solar power} ou \eng{solar energy}, pas les deux.
\item \eng{wind power} : attention, \eng{wind} (le vent) se prononce « ouïnd », alors que le verbe \eng{to wind} (enrouler) se prononce « ouaïnd ».
\end{itemize}
\end{attention}

\courssub{Les acteurs du changement : who does what?}

Protéger l'environnement est l'affaire de nombreux acteurs. Voici le vocabulaire pour les nommer :

\begin{itemize}
\item \cle{environmentalists} (écologistes, défenseurs de l'environnement) : \eng{Environmentalists are calling for stricter laws.} (Les écologistes réclament des lois plus strictes.)
\item \cle{activists} (militants) : \eng{Young activists organised a march in Yaoundé.} (De jeunes militants ont organisé une marche à Yaoundé.)
\item \cle{NGOs} = \emph{non-governmental organisations} (ONG) : \eng{Several NGOs plant trees in the North Region.} (Plusieurs ONG plantent des arbres dans la région du Nord.)
\item \cle{the government} (le gouvernement) : \eng{The government should invest in renewable energy.} (Le gouvernement devrait investir dans les énergies renouvelables.)
\item \cle{local communities} (les communautés locales) : \eng{Local communities manage the forest together.} (Les communautés locales gèrent la forêt ensemble.)
\item \cle{eco-friendly consumers} (les consommateurs écoresponsables) : \eng{Eco-friendly consumers avoid plastic packaging.} (Les consommateurs écoresponsables évitent les emballages plastiques.)
\end{itemize}

\begin{attention}[Faux ami : « militant » et « activist »]
Le français « écologiste » se traduit par \eng{environmentalist} ou \eng{green activist}, jamais par \eng{*ecologist} dans le sens de « militant » : un \eng{ecologist} est d'abord un \emph{scientifique} spécialiste de l'écologie. De même, \eng{the government} est singulier en anglais britannique courant quand on le considère comme un bloc : \eng{The government has promised...} (et non \eng{*the government have} dans un usage scolaire prudent).
\end{attention}

\courssub{Les concepts clés : sustainable development et ses amis}

\begin{definition}[Sustainable development]
\cle{Sustainable development} (développement durable) : \emph{development that meets the needs of the present without compromising the ability of future generations to meet theirs} — un développement qui répond aux besoins du présent sans compromettre la capacité des générations futures à répondre aux leurs. Le nom correspondant est \cle{sustainability} (la durabilité).
\end{definition}

Autour de ce concept central gravitent plusieurs adjectifs et noms indispensables :

\begin{itemize}
\item \cle{eco-friendly} (écologique, respectueux de l'environnement) : \eng{Eco-friendly products do not harm nature.} (Les produits écologiques ne nuisent pas à la nature.)
\item \cle{green} (vert, écologique) : \eng{Green policies favour public transport.} (Les politiques écologiques favorisent les transports en commun.)
\item \cle{biodegradable} (biodégradable) : \eng{Banana leaves are biodegradable packaging.} (Les feuilles de bananier sont un emballage biodégradable.)
\item \cle{conservation} (la conservation, la préservation de la nature) : \eng{Wildlife conservation protects gorillas in the South-West.} (La conservation de la faune protège les gorilles du Sud-Ouest.)
\item \cle{awareness campaigns} (campagnes de sensibilisation) : \eng{Awareness campaigns teach people about recycling.} (Les campagnes de sensibilisation apprennent aux gens le recyclage.)
\end{itemize}

\begin{attention}[« Green » ne veut pas toujours dire « vert »]
Dans ce contexte, l'adjectif \cle{green} signifie « écologique » : \eng{green energy} (énergie propre / verte), \eng{green policies} (politiques écologiques), \eng{green products} (produits écologiques), \eng{to go green} (adopter un mode de vie écologique). Ne traduisez donc pas \eng{green energy} par « énergie verte » au sens littéral dans une rédaction : dites « énergie propre » ou « énergie renouvelable ». Attention aussi au faux ami \eng{*sustainable} : il signifie « durable », pas « soutenable » au sens de « défendable ».
\end{attention}

\courssub{Tableau récapitulatif : 15 termes essentiels}

\begin{center}
\begin{tabularx}{\linewidth}{|l|X|}
\hline
\rowcolor{popblueL} English & Français \\
\hline
to protect the environment & protéger l'environnement \\
\hline
to preserve / conservation & préserver / la conservation \\
\hline
to save energy & économiser l'énergie \\
\hline
to reduce, reuse, recycle (the 3 Rs) & réduire, réutiliser, recycler (les 3 R) \\
\hline
to plant trees & planter des arbres \\
\hline
to ban plastics & interdire les plastiques \\
\hline
renewable energy & énergie renouvelable \\
\hline
solar power / wind power / hydropower & énergie solaire / éolienne / hydroélectricité \\
\hline
fossil fuels & énergies fossiles \\
\hline
sustainable development / sustainability & développement durable / durabilité \\
\hline
eco-friendly / green & écologique \\
\hline
biodegradable & biodégradable \\
\hline
environmentalist / activist / NGO & écologiste / militant / ONG \\
\hline
awareness campaign & campagne de sensibilisation \\
\hline
carbon footprint & empreinte carbone \\
\hline
\end{tabularx}
\end{center}

\courssub{Les collocations des solutions : les verbes qui font agir}

En anglais, certains verbes s'associent presque obligatoirement à certains noms : ce sont les \cle{collocations}. Les apprendre par blocs évite les calques du français.

\begin{center}
\begin{tabularx}{\linewidth}{|l|X|X|}
\hline
\rowcolor{popblueL} Collocation & Traduction & Exemple \\
\hline
to raise awareness & sensibiliser (litt. « élever la conscience ») & \eng{The club raised awareness about deforestation.} \\
\hline
to take action & agir, passer à l'action & \eng{We must take action before it is too late.} \\
\hline
to pass laws & adopter / voter des lois & \eng{Parliament passed a law against plastic bags.} \\
\hline
to reach an agreement & parvenir à un accord & \eng{Countries reached an agreement in Paris in 2015.} \\
\hline
to launch a campaign & lancer une campagne & \eng{They launched a campaign to clean the Wouri river.} \\
\hline
to make a difference & faire la différence, changer les choses & \eng{Every citizen can make a difference.} \\
\hline
\end{tabularx}
\end{center}

\begin{attention}[Pièges de traduction]
\begin{itemize}
\item « Sensibiliser » ne se traduit pas par \eng{*to sensitise} dans l'anglais courant : dites \eng{to raise awareness (of/about)}.
\item « Agir » au sens de « prendre des mesures » se dit \eng{to take action}, pas \eng{*to act actions}.
\item « Parvenir à un accord » se dit \eng{to reach an agreement}, pas \eng{*to arrive at an agreement} (calque du français « arriver à »).
\end{itemize}
\end{attention}

\begin{savaistu}[The Paris Agreement]
\cle{The Paris Agreement} (l'Accord de Paris) est un traité international signé en 2015 par presque tous les pays du monde, dont le Cameroun. It aims to limit global warming to 1.5°C above pre-industrial levels (il vise à limiter le réchauffement climatique à 1,5 degré Celsius au-dessus des niveaux préindustriels). Pour y parvenir, chaque pays s'engage à réduire ses émissions de gaz à effet de serre — \eng{greenhouse gas emissions}. C'est un exemple célèbre de la collocation \eng{to reach an agreement}.
\end{savaistu}

\courssub{Exercice résolu et entraînement}

\begin{exoresolu}[Complétez avec le mot ou l'expression qui convient]
Énoncé — Complete the sentences with: \emph{raise awareness, biodegradable, fossil fuels, take action, solar power, ban}.
\begin{enumerate}
\item Burning \ldots\ldots\ldots\ increases global warming.
\item The mayor decided to \ldots\ldots\ldots\ plastic bags in the market.
\item Students organised a concert to \ldots\ldots\ldots\ about climate change.
\item This packaging is \ldots\ldots\ldots : it disappears naturally in a few months.
\item My village uses \ldots\ldots\ldots\ to light the streets at night.
\item We cannot wait any longer: we must \ldots\ldots\ldots\ now!
\end{enumerate}
\tcblower
Solution détaillée :
\begin{enumerate}
\item \textbf{fossil fuels} — « brûler des énergies fossiles » ; le verbe \eng{to burn} impose un combustible.
\item \textbf{ban} — « interdire » ; après \eng{decided to}, on attend un verbe à l'infinitif.
\item \textbf{raise awareness} — collocation obligatoire pour « sensibiliser » ; \eng{*sensitise} serait un calque du français.
\item \textbf{biodegradable} — la deuxième partie de la phrase (« disparaît naturellement ») définit ce mot.
\item \textbf{solar power} — l'énergie solaire éclaire les rues la nuit grâce aux lampes solaires.
\item \textbf{take action} — collocation signifiant « agir » ; l'urgence (« now! ») la confirme.
\end{enumerate}
\end{exoresolu}

\begin{exercice}[À vous de jouer]
\begin{enumerate}
\item Traduisez en anglais : « Les communautés locales doivent préserver la forêt et planter des arbres. »
\item Traduisez en anglais : « Le gouvernement a adopté une loi pour promouvoir les énergies renouvelables. »
\item Trouvez l'intrus et justifiez : \eng{solar power — wind power — coal — hydropower}.
\item Rédigez trois phrases sur votre école en utilisant \eng{reduce}, \eng{reuse} et \eng{recycle}.
\end{enumerate}
\end{exercice}

\begin{corrige}[Exercice « À vous de jouer »]
\begin{enumerate}
\item \eng{Local communities must preserve the forest and plant trees.}
\item \eng{The government has passed a law to promote renewable energy.}
\item L'intrus est \eng{coal} (le charbon) : c'est une énergie fossile, alors que les trois autres sont des énergies renouvelables.
\item Exemple de production : \eng{Our school should reduce paper waste. We can reuse old exercise books for rough work. Students must recycle plastic bottles instead of burning them.}
\end{enumerate}
\end{corrige}

\begin{aretenir}[L'essentiel de la section]
\begin{itemize}
\item Les verbes d'action : \eng{to protect, to preserve, to save, to reduce, to reuse, to recycle, to plant, to ban}.
\item Les énergies : \eng{renewable energy} (solar, wind, hydro) contre \eng{fossil fuels} (coal, oil, gas).
\item Les acteurs : \eng{environmentalists, activists, NGOs, the government, local communities, eco-friendly consumers}.
\item Les concepts : \eng{sustainable development, eco-friendly, green, biodegradable, conservation, awareness campaigns}.
\item Les collocations à mémoriser par blocs : \eng{to raise awareness, to take action, to pass laws, to reach an agreement, to launch a campaign, to make a difference}.
\item Piège majeur : « sensibiliser » = \eng{to raise awareness}, jamais \eng{*to sensitise}.
\end{itemize}
\end{aretenir}

\coursec{Collocations et expressions idiomatiques}

Dans la section précédente, nous avons étudié le vocabulaire des solutions et de la protection de l'environnement : \eng{recycling}, \eng{renewable energy}, \eng{conservation}. Mais connaître un mot isolément ne suffit pas pour s'exprimer correctement en anglais. Un mot vit en compagnie d'autres mots : on ne dit pas n'importe quel verbe avec n'importe quel nom. Cette section est consacrée aux \cle{collocations} — ces associations naturelles de mots — et aux \cle{expressions idiomatiques} du thème de l'environnement, indispensables pour briller au Bac, à l'écrit comme à l'oral.

\courssub{Qu'est-ce qu'une collocation ?}

\begin{definition}[Collocation]
Une \cle{collocation} est une combinaison habituelle et naturelle de deux mots ou plus, que les anglophones associent spontanément : \eng{heavy rain} (pluie abondante), \eng{make a decision} (prendre une décision), \eng{reduce emissions} (réduire les émissions). En anglais, on dit \eng{strong coffee} et non \eng{powerful coffee} ; de même, on dit \eng{tackle climate change} et non \eng{attack climate change}. La collocation correcte est celle que les natifs emploient : elle ne se devine pas toujours à partir du français.
\end{definition}

\begin{methode}[Comment retenir les collocations]
\begin{enumerate}
\item N'apprends jamais un mot isolément : apprends-le \textbf{par paire} (\eng{verb + noun} : \eng{reduce + emissions} ; \eng{adjective + noun} : \eng{renewable + energy}).
\item Note chaque collocation dans une phrase complète, avec sa traduction.
\item Classe tes paires par thème (problèmes / solutions) dans ton carnet de vocabulaire.
\item Réemploie chaque paire dans une phrase personnelle sur le Cameroun : c'est le meilleur moyen de la mémoriser et de la réutiliser au Bac.
\end{enumerate}
\end{methode}

\courssub{Collocations verbe + nom : les verbes d'action}

Voici les collocations essentielles du thème « environnement ». Observe d'abord le tableau, puis les exemples en contexte.

\begin{center}
\begin{tabularx}{\linewidth}{|X|X|X|}
\hline
\rowcolor{popblueL} Collocation (English) & Traduction & Exemple en contexte \\
\hline
to tackle climate change & s'attaquer au changement climatique & Governments must \textbf{tackle climate change} now. \\
\hline
to combat climate change & combattre le changement climatique & Planting trees helps to \textbf{combat climate change}. \\
\hline
to reduce emissions & réduire les émissions & Factories should \textbf{reduce their emissions}. \\
\hline
to release gases & rejeter / émettre des gaz & Cars \textbf{release harmful gases} into the air. \\
\hline
to cause damage & causer des dégâts & Floods \textbf{cause serious damage} to crops. \\
\hline
to suffer from drought & souffrir de la sécheresse & The Far North \textbf{suffers from drought} every year. \\
\hline
to protect wildlife & protéger la faune sauvage & National parks \textbf{protect wildlife} from hunters. \\
\hline
to save energy & économiser l'énergie & Switch off the lights to \textbf{save energy}. \\
\hline
to protect the environment & protéger l'environnement & Everyone should \textbf{protect the environment}. \\
\hline
\end{tabularx}
\end{center}

\begin{exemplebox}[Exemples traduits]
\eng{Cameroon has suffered from drought in the northern regions for several years.} « Le Cameroun souffre de la sécheresse dans les régions du Nord depuis plusieurs années. »

\eng{If we do not reduce emissions, global warming will cause even more damage.} « Si nous ne réduisons pas les émissions, le réchauffement climatique causera encore plus de dégâts. »

\eng{The government released a plan to protect wildlife in the Dja Reserve.} « Le gouvernement a publié un plan pour protéger la faune sauvage dans la réserve du Dja. »
\end{exemplebox}

\begin{propriete}[« Tackle » et « combat », pas « fight against » au Bac]
En anglais soutenu, on \textbf{tackle} ou on \textbf{combat} climate change — le verbe est suivi directement du nom, sans préposition. Évite \eng{to fight against climate change} dans une rédaction au Bac : si tu veux employer « fight », utilise la \textbf{forme nominale} : \eng{the fight against climate change} (« la lutte contre le changement climatique »). De même : \eng{the fight against pollution}, \eng{the fight against deforestation}.
\end{propriete}

\begin{popfigure}
\begin{center}
\begin{tikzpicture}[
  verbe/.style={rectangle, rounded corners, draw=popblue, fill=popblueL, font=\bfseries, align=center, minimum height=0.8cm},
  nom/.style={rectangle, rounded corners, draw=popgreen, fill=popgreenL, align=center, minimum height=0.8cm},
  fleche/.style={-{Stealth[length=2.5mm]}, thick, poporange}
]
\node[verbe] (v1) at (0,3) {tackle / combat};
\node[nom] (n1) at (6,3) {climate change};
\draw[fleche] (v1) -- (n1);
\node[verbe] (v2) at (0,1.8) {reduce};
\node[nom] (n2) at (6,1.8) {emissions};
\draw[fleche] (v2) -- (n2);
\node[verbe] (v3) at (0,0.6) {protect};
\node[nom] (n3) at (6,0.6) {the environment};
\draw[fleche] (v3) -- (n3);
\node[verbe] (v4) at (0,-0.6) {save};
\node[nom] (n4) at (6,-0.6) {energy};
\draw[fleche] (v4) -- (n4);
\node[font=\itshape, popmuted] at (3,4) {Verb + Noun : les paires à mémoriser};
\end{tikzpicture}
\end{center}
\legende{Quatre collocations « verbe + nom » fondamentales du thème environnement.}
\end{popfigure}

\courssub{Collocations adjectif + nom}

L'adjectif aussi a ses partenaires privilégiés. On dit \eng{global warming} et non \eng{worldwide warming} ; \eng{natural disaster} et non \eng{natural catastrophe} dans le langage courant.

\begin{center}
\begin{tabularx}{\linewidth}{|X|X|X|}
\hline
\rowcolor{popgreenL} Collocation (English) & Traduction & Exemple en contexte \\
\hline
global warming & le réchauffement climatique & \textbf{Global warming} is melting the ice caps. \\
\hline
renewable energy & l'énergie renouvelable & Solar power is a \textbf{renewable energy}. \\
\hline
natural disaster & la catastrophe naturelle & The flood was a terrible \textbf{natural disaster}. \\
\hline
endangered species & les espèces menacées & Elephants are an \textbf{endangered species}. \\
\hline
toxic waste & les déchets toxiques & Factories must not dump \textbf{toxic waste} in rivers. \\
\hline
clean energy & l'énergie propre & Wind farms produce \textbf{clean energy}. \\
\hline
\end{tabularx}
\end{center}

\begin{exemplebox}[Exemples traduits]
\eng{The Cross River gorilla is one of the most endangered species in Africa.} « Le gorille de Cross River est l'une des espèces les plus menacées d'Afrique. »

\eng{Investing in clean energy is the best answer to global warming.} « Investir dans les énergies propres est la meilleure réponse au réchauffement climatique. »
\end{exemplebox}

\courssub{Expressions idiomatiques du thème}

\begin{definition}[Expression idiomatique]
Une \cle{expression idiomatique} est une expression figée dont le sens global ne se déduit pas du sens de chaque mot. Elle ne se traduit jamais mot à mot : il faut la connaître comme un tout.
\end{definition}

\begin{center}
\begin{tabularx}{\linewidth}{|X|X|X|}
\hline
\rowcolor{poppurpleL} Idiom (English) & Équivalent français & Sens réel \\
\hline
the tip of the iceberg & la partie émergée de l'iceberg & la petite partie visible d'un problème bien plus grand \\
\hline
to be in hot water & être dans de beaux draps & être en difficulté, avoir des ennuis \\
\hline
a drop in the ocean & une goutte d'eau dans l'océan & une action très petite face à l'ampleur du problème \\
\hline
to go green & adopter un mode de vie écologique & devenir écolo, changer ses habitudes pour la planète \\
\hline
carbon-neutral & neutre en carbone & qui n'ajoute pas de CO\textsubscript{2} net à l'atmosphère \\
\hline
\end{tabularx}
\end{center}

\begin{exemplebox}[Exemples traduits]
\eng{Recycling one bottle is a drop in the ocean, but every action counts.} « Recycler une bouteille est une goutte d'eau dans l'océan, mais chaque action compte. »

\eng{The factory was in hot water after it released toxic waste into the Wouri river.} « L'usine a eu de gros ennuis après avoir rejeté des déchets toxiques dans le Wouri. »

\eng{Our school decided to go green: no more plastic bags, and trees in the playground.} « Notre école a décidé de passer au vert : plus de sachets plastiques, et des arbres dans la cour. »

\eng{The recent floods in Douala are only the tip of the iceberg.} « Les inondations récentes de Douala ne sont que la partie émergée de l'iceberg. »
\end{exemplebox}

\begin{attention}[« To go green » ne veut pas dire « devenir vert »]
Ne traduis jamais \eng{to go green} par « devenir vert » ! L'expression signifie \textbf{adopter un comportement écologique}. De même, \eng{to be in hot water} ne signifie pas « être dans l'eau chaude » mais « avoir des ennuis ». Les idioms se traduisent par leur équivalent de sens, jamais mot à mot.
\end{attention}

\courssub{Verbes à particule (phrasal verbs) de l'environnement}

Les \cle{phrasal verbs} (verbe + particule) sont omniprésents dans le discours sur l'environnement. La particule change le sens du verbe : \eng{to cut} (couper) devient \eng{to cut down} (abattre / réduire).

\begin{center}
\begin{tabularx}{\linewidth}{|X|X|X|}
\hline
\rowcolor{popgoldL} Phrasal verb & Traduction & Exemple en contexte \\
\hline
to cut down (trees) & abattre (des arbres) ; réduire & Loggers \textbf{cut down} thousands of trees every year. \\
\hline
to use up (resources) & épuiser (des ressources) & We are \textbf{using up} the planet's resources. \\
\hline
to run out of (oil) & être à court de (pétrole) & The world will \textbf{run out of} oil one day. \\
\hline
to throw away & jeter (aux ordures) & Do not \textbf{throw away} plastic bottles: recycle them! \\
\hline
to clean up & nettoyer, assainir & Students \textbf{cleaned up} the beach in Kribi. \\
\hline
\end{tabularx}
\end{center}

\begin{propriete}[Attention à la construction de « run out of »]
\eng{To run out of} se construit avec \textbf{of} + nom : \eng{We are running out of water.} « Nous sommes en train d'épuiser l'eau. » Ne dis jamais \eng{run out water}. À la voix « inversée », on dit aussi : \eng{Water is running out.} « L'eau s'épuise. » (sans \eng{of} quand il n'y a pas de complément).
\end{propriete}

\courssub{Familles de mots et prononciation}

Connaître les dérivés d'un mot de base permet d'enrichir son vocabulaire rapidement. Le tableau ci-dessous regroupe les familles principales du thème. La prononciation est notée entre guillemets à la française.

\begin{center}
\begin{tabularx}{\linewidth}{|X|X|X|X|}
\hline
\rowcolor{popblueL} Verbe & Nom & Adjectif & Prononciation (clé) \\
\hline
to pollute & pollution, pollutant & polluted & « peu-lout », « peu-lou-chn », « peu-lou-tant » \\
\hline
to conserve & conservation, conservationist & conserved & « kon-serve », « kon-ser-vei-chn » \\
\hline
to protect & protection, protector & protective & « pro-tekt », « pro-tek-chn » \\
\hline
to emit & emission, emitter & — & « i-mit », « i-mi-chn » \\
\hline
to sustain & sustainability, sustainer & sustainable & « seus-tein », « seus-te-nai-bi-li-ti » \\
\hline
to recycle & recycling, recycler & recycled & « ri-sai-kl », « ri-sai-kling » \\
\hline
to endanger & endangerment & endangered & « èn-den-djerd » \\
\hline
to warm & warming, warmth & warm & « worm », « worm-ing », « wormth » \\
\hline
\end{tabularx}
\end{center}

\begin{aretenir}[Rappels de prononciation utiles]
\begin{itemize}
\item \eng{drought} se prononce « draut » (le \eng{gh} ne se prononce pas, \eng{ou} = /aʊ/ comme dans \eng{out}).
\item \eng{species} : « spi-chiz » (pluriel identique au singulier).
\item \eng{catastrophe} n'existe pas en anglais courant pour « catastrophe naturelle » ; on dit \eng{disaster} (« di-zas-te ») ou \eng{natural disaster}.
\item \eng{carbon-neutral} : « car-beun niou-tral ».
\end{itemize}
\end{aretenir}

\courssub{Texte d'application : les idioms en contexte}

\begin{textebox}[A warning from Douala]
Last August, the people of Douala watched the water rise in their streets after three days of heavy rain. Houses were damaged, shops were closed, and hundreds of families had to leave their homes. “This is only the tip of the iceberg,” warned Dr. Eyenga, an environmental scientist at the University of Douala. “If we do not reduce our emissions, the floods we saw in Douala will be just the beginning. Global warming is causing serious damage all over the country: the Far North suffers from drought while the coast suffers from floods.”

The government has promised to tackle climate change by investing in renewable energy and by protecting the forests of the South and East regions, where loggers continue to cut down thousands of trees. Some companies have decided to go green: a cement factory near Bonabéri now uses clean energy and has reduced its emissions by half. “It may look like a drop in the ocean,” admits the factory manager, “but if every company does the same, we will save the environment together.”

Young people are also taking action. Last Saturday, students from several high schools cleaned up the banks of the Wouri river and collected two tonnes of plastic waste. “We must stop throwing away plastic everywhere,” said Amina, a Terminale student. “We are running out of time. If we do not protect our environment now, our children will be in hot water.”
\end{textebox}

\textbf{Glossaire.} \emph{to warn} (v.) : avertir ; \emph{logger} (n.) : bûcheron, exploitant forestier ; \emph{by half} : de moitié ; \emph{banks} (n.) : les berges ; \emph{tonne} (n.) : une tonne (orthographe britannique) ; \emph{to take action} : passer à l'action.

\textbf{Comprehension questions.}
\begin{enumerate}
\item What warning does Dr. Eyenga give about the Douala floods?
\item What two measures has the government promised to take?
\item What did the students do last Saturday, and why?
\end{enumerate}

\begin{exoresolu}[Mini-phrases à compléter]
Complète chaque phrase avec la collocation ou l'expression qui convient : \eng{tip of the iceberg} / \eng{reduce emissions} / \eng{cut down} / \eng{go green} / \eng{running out of} / \eng{endangered species}.

1. Factories must \ldots\ if we want to stop global warming.
2. Loggers \ldots\ too many trees in the East region.
3. The mountain gorilla is an \ldots.
4. The flood was terrible, but scientists say it is only the \ldots.
5. We are \ldots\ clean water in many villages.
6. More families should \ldots\ and use solar panels.
\tcblower
\textbf{Solution détaillée.}

1. \textbf{reduce emissions} — collocation verbe + nom : « réduire les émissions ». (On ne dit pas \eng{diminish emissions} dans ce registre.)
2. \textbf{cut down} — phrasal verb : « abattre » des arbres. (\eng{Cut} seul signifie simplement « couper ».)
3. \textbf{endangered species} — collocation adjectif + nom : « espèce menacée ».
4. \textbf{tip of the iceberg} — idiom : la partie visible d'un problème plus vaste.
5. \textbf{running out of} — « être en train d'épuiser » ; attention au \eng{of} obligatoire devant le nom.
6. \textbf{go green} — idiom : « adopter un mode de vie écologique » (et non « devenir vert » !).
\end{exoresolu}

\begin{attention}[Les pièges des francophones — Faux amis et calques]
\begin{itemize}
\item Ne traduis pas « lutter contre » par un verbe + \eng{against} : dis \eng{to tackle climate change}, \eng{to combat pollution}, ou la forme nominale \eng{the fight against climate change}.
\item \eng{To protect} se construit \textbf{sans préposition} : \eng{protect the environment}, jamais \eng{protect against the environment} (sauf \eng{protect wildlife \textbf{from} hunters} : protéger la faune \textbf{des} chasseurs).
\item \eng{A drop in the ocean} : on dit \eng{in}, pas \eng{of} — ne calque pas « une goutte d'eau \emph{dans} l'océan » mot à mot, mais retiens la forme figée anglaise.
\item \eng{To throw away} : la particule \eng{away} est indispensable pour le sens « jeter » ; \eng{to throw} seul signifie « lancer ».
\item \eng{Environment} = environnement (milieu naturel) ; \eng{surroundings} = environs (ce qui entoure). Ne confonds pas.
\item \eng{Catastrophe} existe en anglais mais est plus fort que \eng{disaster} ; pour « catastrophe naturelle », l'usage impose \eng{natural disaster}.
\item \eng{Réduire} $\neq$ \eng{diminish} (qui est plus formel, mathématique) ; au Bac, \eng{reduce emissions / waste / pollution}.
\end{itemize}
\end{attention}

\begin{exercice}[À toi de jouer]
1. Traduis en anglais : « Le gouvernement doit s'attaquer au changement climatique et réduire les émissions des usines. »
2. Trouve l'idiom qui convient : « Nettoyer une seule plage ne suffit pas : c'est \ldots\ dans l'océan. »
3. Construis trois phrases avec : \eng{to run out of}, \eng{to go green}, \eng{toxic waste}.
4. Transforme en forme nominale : \eng{We must fight against deforestation} $\rightarrow$ \eng{The \ldots\ deforestation is essential.}
\end{exercice}

\begin{corrige}[Exercice « À toi de jouer »]
1. \eng{The government must tackle climate change and reduce factory emissions.} (ou \eng{emissions from factories}).
2. \eng{a drop in the ocean} : \eng{Cleaning one beach is not enough: it is a drop in the ocean.}
3. Exemples attendus : \eng{We are running out of clean water.} / \eng{Our school wants to go green.} / \eng{Factories must not dump toxic waste in rivers.}
4. \eng{The fight against deforestation is essential.}
\end{corrige}

\begin{corrige}[Questions de compréhension du texte]
1. Dr. Eyenga warns that the Douala floods are only \eng{the tip of the iceberg} — a small visible part of a much larger problem caused by global warming.
2. The government promised to \eng{invest in renewable energy} and to \eng{protect the forests} of the South and East regions.
3. The students \eng{cleaned up the banks of the Wouri river} and collected two tonnes of plastic waste because they want to stop plastic pollution and protect their environment.
\end{corrige}

\begin{aretenir}[L'essentiel de la section]
\begin{itemize}
\item Une \cle{collocation} s'apprend par paire : \eng{tackle + climate change}, \eng{reduce + emissions}, \eng{renewable + energy}, \eng{endangered + species}.
\item Au Bac, préfère \eng{to tackle / to combat climate change} ou la forme nominale \eng{the fight against climate change}.
\item Les idioms (\eng{the tip of the iceberg}, \eng{a drop in the ocean}, \eng{to go green}, \eng{to be in hot water}) se traduisent par leur équivalent de sens, jamais mot à mot.
\item Les phrasal verbs (\eng{cut down}, \eng{use up}, \eng{run out of}, \eng{throw away}, \eng{clean up}) sont indispensables pour parler de l'environnement naturellement.
\item Maîtrise les familles de mots : \eng{pollute / pollution / polluted}, \eng{protect / protection / protective}, \eng{sustain / sustainability / sustainable}, etc.
\item Attention aux faux amis et calques : \eng{protect} (sans préposition), \eng{natural disaster} (pas \eng{catastrophe}), \eng{reduce} (pas \eng{diminish}), \eng{a drop in the ocean} (pas \eng{of}).
\end{itemize}
\end{aretenir}

\coursec{Faux amis et pièges des francophones}

Dans la section précédente, nous avons appris à combiner les mots de l'environnement en \cle{collocations} naturelles (\eng{to raise awareness}, \eng{to cut emissions}) et à manier quelques expressions idiomatiques. Mais connaître les bons mots ne suffit pas : il faut encore éviter de les assembler « à la française ». C'est précisément l'objet de cette section : repérer les \cle{faux amis} (mots anglais qui ressemblent à des mots français mais ont un sens différent) et les \cle{calques} (traductions mot à mot qui produisent un anglais incorrect). Ce sont les erreurs les plus fréquentes — et les plus coûteuses — chez les candidats francophones au Baccalauréat.

\courssub{Observation : repérer les erreurs dans un texte}

\begin{textebox}[A student's essay — with mistakes]
Last year, our school organised a campaign to protect the nature. I am very sensible to this problem because the warming global affects everyone in Cameroon. Actually, the temperatures are rising and the farmers in my village near Bafoussam are suffering. Many factories make pollution in the rivers, and people do not always have an ecologic behaviour. We must economise our natural resource and recycle our wastes. If we do nothing, the consequences will be terrible for the next generations.
\end{textebox}

\textbf{Glossaire :} \emph{campaign} (n., campagne) ; \emph{to rise} (v., augmenter) ; \emph{behaviour} (n., comportement) ; \emph{waste} (n., déchets).

Ce texte, écrit par un élève francophone, contient \textbf{huit erreurs} de faux amis ou de calques. Saurez-vous toutes les retrouver avant de lire la suite ? Les voici, corrigées :

\begin{center}
\begin{tabularx}{\linewidth}{|X|X|}
\hline
\rowcolor{popblueL} \textbf{Erreur de l'élève} & \textbf{Forme correcte} \\
\hline
to protect the nature & to protect nature / the environment \\
\hline
I am very sensible to this problem & I am very concerned about this problem \\
\hline
the warming global & global warming \\
\hline
Actually, the temperatures are rising & Currently, temperatures are rising \\
\hline
factories make pollution & factories cause pollution / pollute the rivers \\
\hline
an ecologic behaviour & an eco-friendly behaviour \\
\hline
economise our natural resource & save our natural resources \\
\hline
recycle our wastes & recycle our waste \\
\hline
\end{tabularx}
\end{center}

\courssub{Les grands faux amis du thème de l'environnement}

\begin{definition}[Qu'est-ce qu'un faux ami ?]
Un \cle{faux ami} (\eng{false friend} ou \eng{false cognate}) est un mot anglais qui ressemble orthographiquement à un mot français mais dont le sens est différent. Le danger est double : on croit comprendre le mot anglais, et on l'emploie à tort dans sa propre production.
\end{definition}

\begin{propriete}[Sensible ≠ sensitive]
\begin{itemize}
\item \eng{sensible} (anglais) = \textbf{raisonnable, sensé}. \eng{It is sensible to save water during the dry season.} « Il est raisonnable d'économiser l'eau pendant la saison sèche. »
\item Pour dire « sensible » (émotif, réceptif), l'anglais utilise \eng{sensitive} : \eng{She is very sensitive to criticism.} « Elle est très sensible à la critique. »
\item Pour dire « sensible à une cause, préoccupé par », on préfère \eng{concerned about} ou \eng{aware of} : \eng{Young people are concerned about climate change.} « Les jeunes sont sensibles à la question du changement climatique. »
\end{itemize}
\end{propriete}

\begin{attention}[« I am sensible to climate issues » est FAUX]
Pour traduire « je suis sensible aux questions climatiques », ne dites jamais \eng{I am sensible to climate issues} — cela signifierait « je suis raisonnable au sujet des questions climatiques », ce qui n'a pas de sens. Dites :
\begin{itemize}
\item \eng{I am concerned about climate issues.} « Je suis préoccupé par les questions climatiques. »
\item \eng{I care about the environment.} « L'environnement me tient à cœur. »
\item \eng{I am aware of climate issues.} « Je suis conscient des questions climatiques. »
\end{itemize}
\end{attention}

\begin{propriete}[Actually ≠ actuellement]
\begin{itemize}
\item \eng{actually} = \textbf{en fait, en réalité} (sert à corriger ou préciser une idée). \eng{Many people think Douala is the capital, but actually it is Yaoundé.} « Beaucoup pensent que Douala est la capitale, mais en fait c'est Yaoundé. »
\item « Actuellement » se dit \eng{currently}, \eng{at present} ou \eng{at the moment} : \eng{Temperatures are currently rising.} « Les températures augmentent actuellement. »
\item Même piège avec la famille : \eng{actual} = réel, véritable (≠ actuel) ; \eng{current events} = l'actualité.
\end{itemize}
\end{propriete}

\begin{propriete}[Deception ≠ déception]
\begin{itemize}
\item \eng{deception} (anglais) = \textbf{tromperie, duperie}. \eng{The company's green image was pure deception.} « L'image verte de l'entreprise était une pure tromperie. »
\item « Déception » se dit \eng{disappointment} : \eng{The summit's weak results were a great disappointment to activists.} « Les maigres résultats du sommet ont été une grande déception pour les militants. »
\end{itemize}
\end{propriete}

\begin{exemplebox}[Autres faux amis utiles dans ce chapitre]
\begin{itemize}
\item \eng{to achieve} = accomplir, réaliser (≠ « achever » au sens de finir) : \eng{We achieved our recycling target.} « Nous avons atteint notre objectif de recyclage. »
\item \eng{to assist} = aider (≠ assister à) ; « assister à une conférence » = \eng{to attend a conference}.
\item \eng{an issue} = un problème, une question (≠ une issue au sens de sortie) : \eng{Climate change is a global issue.} « Le changement climatique est un problème mondial. »
\item \eng{to demand} = exiger (plus fort que « demander ») ; « demander » = \eng{to ask (for)}.
\item \eng{eventually} = finalement, à la longue (≠ éventuellement = \eng{possibly, if necessary}).
\end{itemize}
\end{exemplebox}

\courssub{Les calques à bannir}

Un \cle{calque} consiste à traduire mot à mot une expression française. Voici les calques les plus fréquents sur le thème de l'environnement.

\begin{propriete}[Nature, warming, ecologic : trois calques classiques]
\begin{enumerate}
\item « Protéger la nature » : ne dites pas \eng{to protect the nature}. En anglais, \eng{nature} dans ce sens ne prend \textbf{pas d'article} : \eng{to protect nature}. On peut aussi dire \eng{to protect the environment}. \eng{We must protect the environment.} « Nous devons protéger l'environnement. »
\item « Le réchauffement global » : l'ordre des mots anglais place l'adjectif \textbf{avant} le nom : \eng{global warming} (jamais \eng{the warming global}). De même : \eng{climate change} « le changement climatique », \eng{renewable energy} « l'énergie renouvelable ».
\item « Écologique » : l'adjectif \eng{ecologic} n'est pratiquement pas utilisé. Dites \eng{ecological} (relatif à l'écologie) ou \eng{eco-friendly} / \eng{environmentally friendly} (respectueux de l'environnement) : \eng{an eco-friendly product} « un produit écologique ».
\end{enumerate}
\end{propriete}

\begin{propriete}[To pollute, pollution, resources, waste]
\begin{itemize}
\item « Faire de la pollution » : ne dites pas \eng{to make pollution}. Dites \eng{to cause pollution} ou utilisez directement le verbe \eng{to pollute} : \eng{Factories pollute the Wouri river.} « Les usines polluent le fleuve Wouri. »
\item « Ressources naturelles » : \eng{resources} s'emploie au \textbf{pluriel} : \eng{natural resources}. \eng{Cameroon is rich in natural resources.} « Le Cameroun est riche en ressources naturelles. »
\item « Les déchets » : \eng{waste} est \textbf{indénombrable} (pas de pluriel) : \eng{household waste} « les déchets ménagers », \eng{toxic waste} « les déchets toxiques ».
\item Orthographe : \eng{environment} s'écrit avec un \textbf{-n-} avant \emph{-ment} : en-vi-ron-\textbf{n}-ment. Prononciation : « in-vaï-reun-ment ». Le nom correspondant est \eng{environmental} (« environnemental »).
\end{itemize}
\end{propriete}

\begin{exemplebox}[Exemples traduits]
\begin{itemize}
\item \eng{We must protect the environment.} « Nous devons protéger l'environnement. » (jamais \eng{the nature})
\item \eng{Plastic waste causes serious pollution in our oceans.} « Les déchets plastiques provoquent une grave pollution dans nos océans. »
\item \eng{We should use our natural resources wisely.} « Nous devrions utiliser nos ressources naturelles judicieusement. »
\item \eng{Global warming is currently the biggest environmental issue.} « Le réchauffement climatique est actuellement le plus grand problème environnemental. »
\item \eng{Solar panels are an eco-friendly source of energy.} « Les panneaux solaires sont une source d'énergie écologique. »
\end{itemize}
\end{exemplebox}

\courssub{Tableau récapitulatif et carte mentale}

\begin{popfigure}
\begin{tikzpicture}[
box/.style={draw=popdark, rounded corners=2pt, align=center, font=\small, inner sep=4pt},
head/.style={box, fill=popblue, text=white, font=\small\bfseries},
bad/.style={box, fill=poppinkL},
good/.style={box, fill=popgreenL},
fr/.style={box, fill=popcream}
]
\node[head] (h1) at (0,0) {Français};
\node[head] (h2) at (5.2,0) {Faux ami / calque};
\node[head] (h3) at (10.4,0) {Forme correcte};
\node[fr] at (0,-1.1) {sensible (à une cause)};
\node[bad] at (5.2,-1.1) {sensible};
\node[good] at (10.4,-1.1) {concerned about / aware of};
\node[fr] at (0,-2.2) {actuellement};
\node[bad] at (5.2,-2.2) {actually};
\node[good] at (10.4,-2.2) {currently / at present};
\node[fr] at (0,-3.3) {une déception};
\node[bad] at (5.2,-3.3) {a deception};
\node[good] at (10.4,-3.3) {a disappointment};
\node[fr] at (0,-4.4) {protéger la nature};
\node[bad] at (5.2,-4.4) {to protect the nature};
\node[good] at (10.4,-4.4) {to protect nature / the environment};
\node[fr] at (0,-5.5) {faire de la pollution};
\node[bad] at (5.2,-5.5) {to make pollution};
\node[good] at (10.4,-5.5) {to cause pollution / to pollute};
\node[fr] at (0,-6.6) {écologique};
\node[bad] at (5.2,-6.6) {ecologic};
\node[good] at (10.4,-6.6) {ecological / eco-friendly};
\end{tikzpicture}
\legende{Tableau des faux amis et calques : français, piège, forme correcte.}
\end{popfigure}

\begin{aretenir}[Les six réflexes anti-calque]
\begin{enumerate}
\item \eng{sensible} = raisonnable ; « sensible » = \eng{sensitive} ou \eng{concerned about}.
\item \eng{actually} = en fait ; « actuellement » = \eng{currently}.
\item \eng{deception} = tromperie ; « déception » = \eng{disappointment}.
\item \eng{to protect nature / the environment} (jamais \eng{the nature}).
\item \eng{to cause pollution / to pollute} (jamais \eng{to make pollution}).
\item \eng{global warming}, \eng{natural resources}, \eng{eco-friendly}, \eng{environment} avec \textbf{-n-}.
\end{enumerate}
\end{aretenir}

\courssub{Méthode et entraînement}

\begin{methode}[Comment éviter les calques en expression écrite]
\begin{enumerate}
\item \textbf{Repérez} dans votre phrase française les expressions que vous vous apprêtez à traduire mot à mot (verbe + nom, adjectif + nom).
\item \textbf{Méfiez-vous} de tout mot anglais qui ressemble à un mot français : c'est un faux ami potentiel.
\item \textbf{Vérifiez} toute expression traduite mot à mot dans un dictionnaire bilingue fiable : cherchez le mot français et lisez les exemples donnés en contexte.
\item \textbf{Préférez} une formulation simple et sûre à une traduction littérale risquée : \eng{I care about the environment} vaut mieux qu'un \eng{sensible} mal employé.
\item \textbf{Relisez} en vous demandant : « Un Anglais dirait-il cela ainsi ? »
\end{enumerate}
\end{methode}

\begin{exoresolu}[Corrigez les erreurs]
\textbf{Énoncé.} Chaque phrase contient un faux ami ou un calque. Corrigez-la.
\begin{enumerate}
\item I am sensible to the problems of deforestation.
\item Actually, many species are disappearing in the Congo Basin.
\item We must protect the nature for our children.
\item Cars make a lot of pollution in Yaoundé.
\item The results of the conference were a deception for everyone.
\item Our country has many natural resource.
\end{enumerate}
\tcblower
\textbf{Solution détaillée.}
\begin{enumerate}
\item \eng{I am concerned about the problems of deforestation.} — \eng{sensible} signifie « raisonnable » ; pour « sensible à une cause », on dit \eng{concerned about}.
\item \eng{Currently, many species are disappearing in the Congo Basin.} — \eng{actually} signifie « en fait » ; « actuellement » se traduit \eng{currently}.
\item \eng{We must protect nature (ou the environment) for our children.} — \eng{nature} ne prend pas d'article dans ce sens.
\item \eng{Cars cause a lot of pollution in Yaoundé.} (ou \eng{Cars pollute Yaoundé.}) — on ne « fait » pas de la pollution en anglais.
\item \eng{The results of the conference were a disappointment for everyone.} — \eng{deception} signifie « tromperie » ; « déception » = \eng{disappointment}.
\item \eng{Our country has many natural resources.} — \eng{resources} s'emploie au pluriel après \eng{many}.
\end{enumerate}
\end{exoresolu}

\begin{exercice}[À vous de jouer]
Traduisez en anglais correct, en évitant les calques :
\begin{enumerate}
\item Les jeunes sont de plus en plus sensibles aux questions environnementales.
\item Actuellement, les températures augmentent dans le nord du Cameroun.
\item Nous devons économiser nos ressources naturelles et recycler nos déchets.
\item Les usines polluent les rivières et provoquent une grave pollution.
\item Ce produit écologique protège l'environnement.
\end{enumerate}
\end{exercice}

\begin{corrige}[Exercice « À vous de jouer »]
\begin{enumerate}
\item \eng{Young people are more and more concerned about environmental issues.}
\item \eng{Currently, temperatures are rising in the north of Cameroon.}
\item \eng{We must save our natural resources and recycle our waste.}
\item \eng{Factories pollute the rivers and cause serious pollution.}
\item \eng{This eco-friendly product protects the environment.}
\end{enumerate}
\end{corrige}

\begin{savaistu}[Les calques coûtent des points au Bac]
Au Baccalauréat, les correcteurs d'expression écrite repèrent immédiatement les calques du français : \eng{the nature}, \eng{actually} pour « actuellement », \eng{to make pollution} sont des « signatures » typiques du candidat francophone. Chaque faux ami mal employé fait perdre des points en \emph{accuracy} (correction de la langue). À l'inverse, employer correctement \eng{currently}, \eng{concerned about}, \eng{eco-friendly} ou \eng{natural resources} montre une maîtrise authentique de l'anglais et valorise votre copie. Apprendre dix paires de faux amis par cœur est l'un des investissements les plus rentables de votre année de Terminale !
\end{savaistu}

\coursec{Familles de mots, prononciation et emploi en contexte}

Dans la section précédente, nous avons appris à débusquer les faux amis et les calques du français qui guettent le candidat au baccalauréat. Nous allons maintenant franchir une étape supplémentaire : au lieu de mémoriser les mots un par un, nous allons apprendre à les regrouper en \cle{familles de mots} (\eng{word families}). Connaître la famille de \eng{pollute}, c'est connaître d'un seul coup quatre ou cinq mots, savoir les prononcer correctement et les employer dans des phrases naturelles. C'est exactement la compétence qui distingue une copie moyenne d'une excellente copie d'expression écrite.

\courssub{Observer avant d'apprendre : un texte pour découvrir les familles}

Lisons d'abord ce court article original. Soulignez mentalement tous les mots construits sur la même base.

\begin{textebox}[Going green at Government High School Buea — a school newspaper article]
\eng{Last year, our school decided to take environmental problems seriously. The students noticed that the schoolyard was polluted with plastic bags and empty bottles. Air pollution from old lorries on the main road was also a major concern. With the help of two environmentalists from a local association, we created a “Green Club”. Our first action was to protect the small forest behind the classrooms: we planted fifty trees and we now recycle paper and plastic every Friday. The club wants to promote sustainable development, because sustainability is the key to our future. Our head teacher says the project is economical too: recycling costs almost nothing, and the school saves money. “Conservation is not expensive,” she explains. “Destroying nature is.” The club's slogan is simple: “Act locally, think globally.”}
\end{textebox}

\textbf{Glossaire :} \eng{schoolyard} (n., cour de récréation) ; \eng{lorries} (n. pl., camions) ; \eng{concern} (n., sujet d'inquiétude) ; \eng{head teacher} (n., proviseur / directeur) ; \eng{to save money} (exp., économiser de l'argent) ; \eng{slogan} (n., slogan).

\textbf{Questions de compréhension :}
\begin{enumerate}
\item \eng{What two pollution problems did the students notice?}
\item \eng{Who helped the school create the Green Club?}
\item \eng{What does the club do every Friday?}
\item \eng{Why does the head teacher say the project is economical?}
\end{enumerate}

Observons : \eng{environmental}, \eng{environmentalists}, \eng{polluted}, \eng{pollution}, \eng{protect}, \eng{sustainable}, \eng{sustainability}, \eng{conservation}, \eng{globally}. Tous ces mots appartiennent à quelques familles seulement. C'est le secret : le vocabulaire de l'environnement est très régulier.

\courssub{Les grandes familles du vocabulaire de l'environnement}

\begin{definition}[Famille de mots]
Une \cle{famille de mots} (\eng{word family}) est un ensemble de mots construits sur la même racine (\eng{root}) à l'aide de préfixes et de suffixes. Chaque membre de la famille a une \cle{nature grammaticale} différente : verbe, nom, adjectif, adverbe. Exemple : \eng{pollute} (verbe), \eng{pollution} (nom), \eng{pollutant} (nom), \eng{polluted} (adjectif).
\end{definition}

\begin{popfigure}
\begin{tikzpicture}[
grow=down,
level distance=1.8cm,
level 1/.style={sibling distance=4.5cm},
every node/.style={draw, rounded corners, align=center, font=\small, minimum width=2.5cm},
edge from parent/.style={draw=popblue, -{Stealth}, thick}
]
\node[fill=popblueL, font=\small\bfseries, align=center]{pollute\\(verb — polluer)}
child { node[fill=popgreenL, align=center]{pollution\\(noun — la pollution)} }
child { node[fill=poporangeL, align=center]{pollutant\\(noun — le polluant)} }
child { node[fill=poppinkL, align=center]{polluted\\(adjective — pollué)} };
\end{tikzpicture}
\legende{Arbre de la famille de \eng{pollute} : une racine, quatre natures grammaticales.}
\end{popfigure}

\begin{center}
\begin{tabularx}{\linewidth}{|l|l|l|X|}
\hline
\rowcolor{popblueL} \textbf{Verbe} & \textbf{Nom} & \textbf{Adjectif} & \textbf{Exemple} \\
\hline
to pollute & pollution / pollutant & polluted & \eng{Factories pollute the river.} \\
\hline
to protect & protection & protective & \eng{We must protect wildlife.} \\
\hline
to conserve & conservation & conservative & \eng{Water conservation is vital.} \\
\hline
to contaminate & contamination & contaminated & \eng{The water is contaminated.} \\
\hline
to recycle & recycling & recyclable & \eng{Glass is recyclable.} \\
\hline
to sustain & sustainability & sustainable & \eng{Sustainable energy is the future.} \\
\hline
— & environment & environmental & \eng{An environmental disaster.} \\
\hline
— & globe & global & \eng{Global warming is real.} \\
\hline
\end{tabularx}
\end{center}

\begin{propriete}[Le suffixe -tion : du verbe au nom d'action]
Le suffixe \cle{-tion} (et ses variantes \eng{-sion}, \eng{-ation}) transforme un verbe en \cle{nom d'action}. Forme : verbe + \eng{-tion}, souvent avec une petite modification orthographique.
\begin{itemize}
\item \eng{pollute} $\rightarrow$ \eng{pollution} (le \eng{-e} final disparaît) ;
\item \eng{protect} $\rightarrow$ \eng{protection} ;
\item \eng{contaminate} $\rightarrow$ \eng{contamination} (le \eng{-ate} devient \eng{-ation}) ;
\item \eng{conserve} $\rightarrow$ \eng{conservation} ;
\item \eng{recycle} $\rightarrow$ \eng{recycling} (exception : ici on utilise \eng{-ing}, pas \eng{-tion}).
\end{itemize}
Comparaison avec le français : la correspondance est presque parfaite (\eng{pollution}, \eng{protection}, \eng{conservation}), ce qui est une chance pour les francophones — mais attention à la prononciation : en anglais, \eng{-tion} se prononce « cheune » (\eng{pollution} = « po-lou-cheune »), jamais « sion ».
\end{propriete}

\begin{propriete}[Les autres suffixes typiques de ce thème]
\begin{itemize}
\item \cle{-able} : possibilité, « qui peut être ». \eng{recycle} $\rightarrow$ \eng{recyclable} (recyclable) ; \eng{sustain} $\rightarrow$ \eng{sustainable} (durable, soutenable).
\item \cle{-ist} : la personne qui agit ou milite. \eng{environment} $\rightarrow$ \eng{environmentalist} (écologiste, défenseur de l'environnement) ; comparer \eng{activist}, \eng{scientist}.
\item \cle{-al} : formation d'adjectifs. \eng{environment} $\rightarrow$ \eng{environmental} (environnemental) ; \eng{globe} $\rightarrow$ \eng{global} (mondial) $\rightarrow$ \eng{globally} (adverbe en \eng{-ly} : mondialement).
\item \cle{-ity} : nom abstrait de qualité. \eng{sustainable} $\rightarrow$ \eng{sustainability} (durabilité).
\end{itemize}
\end{propriete}

\begin{exemplebox}[Une famille, trois phrases]
\eng{Air pollution is a major problem. Factories pollute the air. The air in the city is polluted.} = « La pollution de l'air est un problème majeur. Les usines polluent l'air. L'air de la ville est pollué. »

Autre famille : \eng{We must protect the forest. The protection of the forest is our duty. Protective laws already exist.} = « Nous devons protéger la forêt. La protection de la forêt est notre devoir. Des lois protectrices existent déjà. »
\end{exemplebox}

\courssub{Prononciation et orthographe : les mots pièges}

Ces mots sont fréquents à l'oral comme à l'écrit ; leur prononciation surprend souvent les francophones.

\begin{center}
\begin{tabularx}{\linewidth}{|l|l|X|}
\hline
\rowcolor{popblueL} \textbf{Mot} & \textbf{Prononciation (lettres françaises)} & \textbf{Piège à éviter} \\
\hline
\eng{environment} & « in-vaï-reun-ment » & ne pas oublier le \eng{n} avant le \eng{m} à l'écrit : \eng{enviroNment} \\
\hline
\eng{sustainable} & « seu-sté-né-beul » & ne pas dire « sou-sté-nable » comme en français \\
\hline
\eng{drought} & « draout » & le \eng{-gh} est muet ; sens : la sécheresse \\
\hline
\eng{government} & « gueu-veurn-ment » & le premier \eng{n} est presque muet à l'oral, mais obligatoire à l'écrit \\
\hline
\eng{pollution} & « po-lou-cheune » & \eng{-tion} = « cheune », pas « sion » \\
\hline
\eng{global warming} & « glo-beul ouor-ming » & \eng{warming} rime avec \eng{morning} \\
\hline
\end{tabularx}
\end{center}

\begin{attention}[economic, economical, ecologic : trois pièges]
\begin{itemize}
\item \eng{economic} = économique, \textbf{lié à l'économie} du pays : \eng{the economic crisis} (la crise économique), \eng{economic growth} (la croissance économique).
\item \eng{economical} = \textbf{qui fait des économies}, qui coûte peu : \eng{Recycling is economical.} (Le recyclage ne coûte pas cher.) \eng{an economical car} (une voiture économe).
\item \eng{ecologic} \textbf{n'existe pas} : il faut dire \eng{ecological} : \eng{an ecological disaster} (une catastrophe écologique). On emploie aussi très souvent \eng{environmental} : \eng{environmental problems}.
\end{itemize}
Ne traduisez donc jamais « une voiture économique » par \eng{an economic car} : dites \eng{an economical car}.
\end{attention}

\courssub{Emploi en contexte : s'entraîner à choisir la bonne forme}

\begin{methode}[Enrichir une copie en variant les formes d'une même famille]
Pour éviter les répétitions dans une composition sur l'environnement :
\begin{enumerate}
\item Repérez le mot que vous allez répéter (souvent \eng{pollution}).
\item Cherchez les autres membres de sa famille : \eng{pollute}, \eng{polluted}, \eng{pollutant}.
\item Changez la construction de la phrase : au lieu de \eng{Pollution is dangerous and pollution kills animals}, écrivez \eng{Polluted water kills animals, and factories that pollute rivers must be punished.}
\item Alternez verbe, nom et adjectif : cela montre à l'examinateur que vous maîtrisez la grammaire et le vocabulaire.
\end{enumerate}
\end{methode}

\begin{exoresolu}[Texte à trous : choisir la bonne forme]
Énoncé — Complétez avec la bonne forme du mot entre parenthèses : \eng{1. The government must (protection) \dots\ the forest. 2. Water (pollute) \dots\ kills fish in our rivers. 3. Plastic bags are not (recycle) \dots\ in our town. 4. (environment) \dots\ problems concern us all. 5. We need (sustain) \dots\ solutions for the future.}
\tcblower
Solution détaillée :
\begin{enumerate}
\item \eng{protect} : après le modal \eng{must}, il faut un \textbf{verbe} à l'infinitif.
\item \eng{pollution} : sujet du verbe \eng{kills}, il faut un \textbf{nom}.
\item \eng{recyclable} : après \eng{are not}, il faut un \textbf{adjectif} (suffixe \eng{-able}).
\item \eng{Environmental} : devant le nom \eng{problems}, il faut un \textbf{adjectif} (suffixe \eng{-al}).
\item \eng{sustainable} : devant le nom \eng{solutions}, adjectif en \eng{-able} ; le nom abstrait serait \eng{sustainability}.
\end{enumerate}
\end{exoresolu}

\begin{exoresolu}[Transformation : du nom au verbe]
Énoncé — Réécrivez chaque phrase en remplaçant le nom souligné par le verbe de la même famille : \eng{1. The \textbf{pollution} of the river by the factory is scandalous. 2. The \textbf{protection} of endangered species is necessary.}
\tcblower
Solution : 1. \eng{The factory pollutes the river; it is scandalous.} (nom \eng{pollution} $\rightarrow$ verbe \eng{pollute}, 3\textsuperscript{e} personne : \eng{pollutes}). 2. \eng{We must protect endangered species.} (nom \eng{protection} $\rightarrow$ verbe \eng{protect} après le modal \eng{must}). Remarquez que la version avec le verbe est plus directe et plus naturelle en anglais : l'anglais préfère souvent les verbes aux noms abstraits, contrairement au français.
\end{exoresolu}

\begin{textebox}[Model dialogue — a green project at school]
\eng{A: What is your school doing to go green?\\
B: We have started a recycling club and we plant trees every month.\\
A: That is wonderful! Who helps you?\\
B: An environmentalist from Douala visits us every term. She teaches us about sustainability and conservation.\\
A: And is the project expensive?\\
B: Not at all. It is very economical: we collect recyclable plastic and we sell it. With the money, we buy young trees.\\
A: Your school is a good example. I will propose the same project in my school in Bamenda.\\
B: Great idea! Remember: act locally, think globally.}

\textbf{Glossaire :} \eng{to go green} (exp., devenir écologique) ; \eng{every term} (exp., chaque trimestre) ; \eng{young trees} (n. pl., jeunes arbres) ; \eng{to propose} (v., proposer).
\end{textebox}

\begin{exercice}[À vous de jouer]
\begin{enumerate}
\item Complétez : \eng{The (globe) \dots\ temperature is rising because of (globe) \dots\ warming.}
\item Transformez : \eng{The contamination of the well is dangerous.} (utilisez le verbe \eng{contaminate}).
\item Prononcez à voix haute : \eng{environment}, \eng{drought}, \eng{government}, \eng{sustainable}.
\end{enumerate}
\end{exercice}

\begin{corrige}[Exercice « À vous de jouer »]
1. \eng{The global temperature is rising because of global warming.} (deux adjectifs en \eng{-al}). 2. \eng{Someone has contaminated the well; it is dangerous.} ou \eng{The well is contaminated.} 3. « in-vaï-reun-ment », « draout », « gueu-veurn-ment », « seu-sté-né-beul ».
\end{corrige}

\begin{experience}[Mini-production guidée : mes cinq phrases]
Rédigez \textbf{cinq phrases personnelles} sur votre ville ou votre école, en utilisant chacun de ces cinq mots cibles : \eng{pollution}, \eng{protect}, \eng{sustainable}, \eng{recyclable}, \eng{environmentalist}. Modèle : \eng{In Douala, air pollution is a serious problem.} Puis lisez vos phrases à votre voisin : il doit identifier la nature grammaticale de chaque mot cible (nom, verbe ou adjectif).
\end{experience}

\begin{aretenir}[L'essentiel de la section]
\begin{itemize}
\item Une racine = plusieurs mots : \eng{pollute} (v) $\rightarrow$ \eng{pollution}, \eng{pollutant} (n) $\rightarrow$ \eng{polluted} (adj) ; \eng{environment} (n) $\rightarrow$ \eng{environmental} (adj) $\rightarrow$ \eng{environmentalist} (n) ; \eng{sustain} (v) $\rightarrow$ \eng{sustainable} (adj) $\rightarrow$ \eng{sustainability} (n).
\item Suffixes-clés : \eng{-tion} (nom d'action), \eng{-able} (possibilité), \eng{-ist} (personne), \eng{-al} (adjectif), \eng{-ity} (qualité abstraite).
\item Prononciation : \eng{environment} « in-vaï-reun-ment », \eng{drought} « draout », \eng{-tion} « cheune ».
\item Pièges : \eng{economic} $\neq$ \eng{economical} ; \eng{ecologic} n'existe pas, dites \eng{ecological} ; n'oubliez pas le \eng{n} de \eng{environment}.
\item Dans une copie, variez verbe, nom et adjectif d'une même famille pour éviter les répétitions.
\end{itemize}
\end{aretenir}

\newpage

\coursec{Activité d'intégration}

\courssub{Objectifs de l'activité}

À la fin de cette activité d'intégration, l'élève sera capable de :

\begin{itemize}
\item réemployer à l'écrit et à l'oral le vocabulaire des trois champs lexicaux de la leçon : \cle{climate change and global warming}, \cle{pollution and waste}, \cle{solutions and protection} ;
\item utiliser correctement au moins trois \cle{collocations} étudiées (par exemple \eng{reduce your carbon footprint}, \eng{raise awareness}, \eng{tackle climate change}) ;
\item employer des mots de la même famille (\eng{pollute / pollution / polluted}, \eng{protect / protection / protective}) en choisissant la bonne nature grammaticale ;
\item éviter les \cle{faux amis} et les calques du français (\eng{environment} et non \eng{environnement}, \eng{to protect the environment} et non \eng{to protect the nature}) ;
\item produire un document authentique (une affiche de campagne) et le présenter oralement en deux minutes, avec une prononciation soignée.
\end{itemize}

\courssub{Situation}

\begin{textebox}[A letter from the Principal — Government High School Biyem-Assi, Yaoundé]
Dear students,

As you know, our school is taking part in the national \textbf{“Green Schools” competition} organised by the Ministry of Secondary Education. Every year, schools across Cameroon present an environmental project, and the best campaign wins a prize of 200,000 CFA francs to create a school garden.

This year, our school wants to win! Your mission, in groups of four, is to design a \textbf{“Go Green!” poster} for our school. Your poster must include:

-- a catchy slogan with a strong collocation, for example “Reduce your carbon footprint!” or “Say no to plastic waste!”;

-- five concrete actions that students can take every day, written in clear, correct English (for example: recycling plastic bottles, planting trees, saving energy, avoiding bush fires, using reusable bags).

Each group will then present its poster to the class in a \textbf{two-minute speech}. The class will vote for the best campaign, and the winning poster will be sent to the Ministry.

Remember: the jury will pay attention to your vocabulary. Use the words and expressions you have learnt about climate change, global warming, pollution and environmental protection. Avoid word-for-word translations from French!

Good luck, and may the greenest team win!

The Principal
\end{textebox}

\textbf{Glossaire :} \emph{catchy} (adj.) : accrocheur ; \emph{bush fire} (n.) : feu de brousse ; \emph{reusable} (adj.) : réutilisable ; \emph{jury} (n.) : jury ; \emph{word-for-word} (adv.) : mot à mot.

\courssub{Consignes}

\begin{enumerate}
\item \textbf{Comprendre la commande.} Lis attentivement la lettre du proviseur. Souligne en rouge ce que l'affiche doit contenir, en bleu les critères du jury. Reformule la tâche en une phrase à l'oral : \eng{We have to...}

\item \textbf{Mobiliser le lexique.} En groupe, faites un remue-méninges : listez au moins dix mots ou expressions de la leçon que vous pourriez utiliser (par exemple \eng{global warming}, \eng{greenhouse gases}, \eng{deforestation}, \eng{carbon footprint}, \eng{to recycle}, \eng{renewable energy}, \eng{to raise awareness}, \eng{to cut down on waste}). Classez-les dans les trois champs lexicaux : problème, pollution, solution.

\item \textbf{Choisir le slogan.} Rédigez un slogan court contenant une \cle{collocation} correcte. Vérifiez la collocation : on dit \eng{reduce your carbon footprint}, \eng{save energy}, \eng{protect the environment} — jamais \eng{reduce your foot carbon} ni \eng{economize the energy} (calque du français « économiser »).

\item \textbf{Rédiger les cinq actions.} Écrivez cinq phrases impératives ou avec \eng{should}, en utilisant des verbes précis : \eng{recycle}, \eng{plant}, \eng{sort}, \eng{switch off}, \eng{reuse}, \eng{avoid}. Variez les familles de mots : si vous utilisez le verbe \eng{to pollute} dans une phrase, utilisez le nom \eng{pollution} dans une autre.

\item \textbf{Relire et corriger.} Chassez les faux amis et les calques : \eng{environment} (et non \eng{environnement}), \eng{actually} signifie « en fait » (et non « actuellement »), \eng{a deceiver} n'existe pas pour « déchets » — on dit \eng{waste} ou \eng{rubbish}. Vérifiez aussi l'orthographe : \eng{environment} (avec un \emph{n} avant \emph{m}), \eng{government}, \eng{pollution}.

\item \textbf{Préparer la prise de parole.} Répartissez le discours de deux minutes entre les quatre membres : introduction (le problème), présentation du slogan, les cinq actions, appel à l'action final. Travaillez la prononciation des mots difficiles : \eng{climate} « klaï-mèt », \eng{pollution} « po-lou-cheune », \eng{sustainable} « se-sté-ne-beul », \eng{environment} « in-vaï-reune-mente ».

\item \textbf{Présenter et voter.} Chaque groupe présente son affiche. Pendant les présentations, les autres élèves complètent une grille d'écoute : slogan entendu, trois actions relevées, une collocation repérée, un éventuel calque du français. La classe vote ensuite pour la meilleure campagne.
\end{enumerate}

\courssub{Ressources à mobiliser}

\begin{aretenir}[Boîte à outils lexicale]
\begin{itemize}
\item \textbf{Problème :} \eng{climate change}, \eng{global warming}, \eng{greenhouse gases}, \eng{the greenhouse effect}, \eng{rising temperatures}, \eng{drought}, \eng{flooding}, \eng{deforestation}.
\item \textbf{Pollution :} \eng{air / water / soil pollution}, \eng{plastic waste}, \eng{rubbish}, \eng{to pollute}, \eng{to dump waste}, \eng{exhaust fumes}, \eng{bush fires}.
\item \textbf{Solutions :} \eng{to recycle}, \eng{to reuse}, \eng{to reduce}, \eng{to sort waste}, \eng{to plant trees}, \eng{to save energy}, \eng{to switch off the lights}, \eng{renewable energy}, \eng{solar panels}, \eng{reusable bags}, \eng{to protect the environment}.
\item \textbf{Collocations :} \eng{reduce your carbon footprint}, \eng{raise awareness}, \eng{tackle / fight climate change}, \eng{cut down on plastic}, \eng{take action}, \eng{make a difference}.
\item \textbf{Familles de mots :} \eng{to pollute} (verbe) / \eng{pollution} (nom) / \eng{polluted} (adjectif) ; \eng{to protect} / \eng{protection} / \eng{protective} ; \eng{to sustain} / \eng{sustainable} / \eng{sustainability}.
\item \textbf{Structures utiles :} impératif (\eng{Plant a tree!}), \eng{should + base verbale} (\eng{We should sort our waste.}), \eng{if + present, will + base verbale} (\eng{If we act now, we will save our planet.}).
\end{itemize}
\end{aretenir}

\begin{attention}[Pièges à éviter dans votre production]
\begin{itemize}
\item \eng{environment} s'écrit avec un \emph{n} : pas \eng{environement}.
\item On dit \eng{to protect the environment}, jamais \eng{to protect the nature} (calque de « protéger la nature »).
\item \eng{actually} = « en fait » ; « actuellement » se dit \eng{currently} ou \eng{at present}.
\item « Les déchets » = \eng{waste} (indénombrable) ou \eng{rubbish} ; jamais \eng{deceivers}.
\item « Sensibiliser » = \eng{to raise awareness}, pas \eng{to sensibilize}.
\end{itemize}
\end{attention}

\courssub{Proposition de production et corrigé commenté}

\begin{exoresolu}[Modèle d'affiche et de discours — Groupe « Green Eagles »]
\textbf{Énoncé.} Rédigez l'affiche (slogan + cinq actions) et le discours de présentation de deux minutes.

\tcblower

\textbf{Solution détaillée.}

\textbf{Affiche :}

\eng{GO GREEN, BIYEM-ASSI! Reduce your carbon footprint today!}

\eng{1. Sort your waste: plastic bottles are recycled, not burned!}

\eng{2. Plant a tree every year to fight deforestation.}

\eng{3. Switch off the lights and save energy.}

\eng{4. Use a reusable bag and say no to plastic waste.}

\eng{5. Walk or share a bendskin: cut down on exhaust fumes!}

\textbf{Discours :}

\eng{Good morning everyone. Our planet is in danger. Because of global warming, temperatures are rising and droughts are becoming more frequent in our country. Here in Yaoundé, plastic waste pollutes our streets and our rivers.}

\eng{That is why our group, the Green Eagles, has created this campaign: “Go Green, Biyem-Assi! Reduce your carbon footprint today!”}

\eng{We propose five simple actions. First, we should sort our waste: plastic bottles must be recycled, not burned. Second, every student should plant a tree every year to fight deforestation. Third, we must switch off the lights when we leave the classroom to save energy. Fourth, we should use reusable bags instead of plastic ones. Finally, if we walk or share transport, we will cut down on exhaust fumes and air pollution.}

\eng{If we act together, we will make a difference. Protect our environment, protect our future. Thank you!}

\textbf{Commentaire des choix de langue :}

\begin{itemize}
\item Le slogan contient la collocation \eng{reduce your carbon footprint}, exigée par la tâche ; l'impératif donne de la force à l'affiche.
\item Les cinq actions réemploient les trois champs lexicaux : problème (\eng{deforestation}, \eng{exhaust fumes}), pollution (\eng{plastic waste}), solutions (\eng{sort}, \eng{recycled}, \eng{plant}, \eng{switch off}, \eng{save energy}, \eng{reusable}).
\item Les familles de mots sont exploitées : verbe \eng{pollutes}, nom \eng{pollution}, participe \eng{recycled} ; verbe \eng{to protect} répété en finale pour l'effet oratoire.
\item Les structures grammaticales varient : impératif (\eng{Sort your waste}), \eng{should} et \eng{must} pour le conseil, et une phrase du premier conditionnel (\eng{If we act together, we will make a difference}) qui donne une conclusion convaincante.
\item Aucun calque du français : \eng{environment} est bien orthographié, \eng{reusable bags} et non \eng{reutilizable bags}, \eng{save energy} et non \eng{economize energy}.
\item Le contexte camerounais (\eng{bendskin}, Yaoundé, sécheresses) rend la production authentique sans artifice.
\end{itemize}
\end{exoresolu}

\courssub{Bilan de l'activité}

Cette activité a permis de vérifier l'acquisition des savoir-faire de la leçon dans une tâche authentique et motivante :

\begin{itemize}
\item \textbf{Lexique :} les élèves ont réemployé les trois champs lexicaux (changement climatique, pollution, solutions) dans un document à visée réelle — convaincre un jury et des camarades.
\item \textbf{Collocations et familles de mots :} le slogan et les actions ont imposé le choix de collocations exactes et la manipulation des natures de mots (verbe, nom, adjectif).
\item \textbf{Vigilance linguistique :} la phase de relecture collective a traqué les faux amis et les calques du français, point faible récurrent des apprenants francophones.
\item \textbf{Expression orale :} la présentation de deux minutes a entraîné la fluidité, la prononciation des mots du thème et la prise de parole en groupe, compétences exigées à l'épreuve orale du baccalauréat.
\item \textbf{Compréhension orale :} la grille d'écoute remplie pendant les présentations a maintenu l'attention de toute la classe et réinvesti le lexique à la réception.
\end{itemize}

Pour prolonger l'activité, l'affiche gagnante peut être reproduite et affichée dans la salle de classe, et chaque élève peut rédiger à la maison un court paragraphe intitulé \eng{My personal green commitment}, préparant ainsi la leçon d'expression écrite du chapitre.

\newpage

\coursec{Méthodes et exercices résolus}

\coursec{Lesson 33 — Vocabulary: Environmental Issues, Climate Change and Global Warming}

\begin{textebox}[Authentic context — Cameroon and the climate crisis]
“In the Far North region of Cameroon, the rainy season has become shorter and unpredictable. Prolonged droughts destroy millet and sorghum crops, pushing families into food insecurity. Further south, in Douala, heavy rains cause devastating floods because plastic waste blocks the drainage systems. Meanwhile, the mangroves of the Wouri estuary are disappearing due to overexploitation and rising sea levels. To tackle these challenges, the government has launched the Green Sahel initiative: a vast reforestation programme combined with the promotion of solar energy in rural areas. Young people are encouraged to raise awareness in their communities and to adopt sustainable habits such as recycling and reducing their carbon footprint.”
\end{textebox}

\courssub{Champ lexical 1 : Climate change and global warming}

\begin{definition}[Mots-clés : Causes et effets du changement climatique]
\begin{center}
\begin{tabularx}{\linewidth}{|X|X|X|}\hline
\rowcolor{popblueL} \textbf{English} & \textbf{Français} & \textbf{Collocation / Note} \\ \hline
\eng{climate change} & changement climatique & \eng{tackle climate change} \\ \hline
\eng{global warming} & réchauffement climatique & \eng{fight global warming} \\ \hline
\eng{greenhouse effect} & effet de serre & \eng{greenhouse gases} (gaz à effet de serre) \\ \hline
\eng{carbon emissions} / \eng{CO2 emissions} & émissions de carbone / CO2 & \eng{reduce / cut emissions} \\ \hline
\eng{fossil fuels} & combustibles fossiles & \eng{depend on fossil fuels} \\ \hline
\eng{deforestation} & déforestation & \eng{massive deforestation} \\ \hline
\eng{rising temperatures} & hausse des températures & \eng{temperatures are rising} \\ \hline
\eng{melting ice caps / glaciers} & fonte des calottes / glaciers & \eng{melting ice caps} \\ \hline
\eng{rising sea levels} & montée du niveau de la mer & \eng{threaten coastal areas} \\ \hline
\eng{drought} & sécheresse & \eng{prolonged / severe drought} \\ \hline
\eng{floods} & inondations & \eng{devastating floods} \\ \hline
\eng{unpredictable rainfall} & pluies imprévisibles & \eng{erratic rainfall} \\ \hline
\end{tabularx}
\end{center}
\textbf{Prononciation (repères français) :} \eng{climate} « klaï-mit », \eng{greenhouse} « gri-ne-haousse », \eng{drought} « draoute » (le \emph{gh} ne se prononce pas), \eng{flood} « flade ».
\end{definition}

\courssub{Champ lexical 2 : Pollution and waste}

\begin{definition}[Mots-clés : Pollution, déchets et impacts]
\begin{center}
\begin{tabularx}{\linewidth}{|X|X|X|}\hline
\rowcolor{poporangeL} \textbf{English} & \textbf{Français} & \textbf{Collocation / Note} \\ \hline
\eng{pollution} & pollution & \eng{air / water / soil pollution} \\ \hline
\eng{to pollute} & polluer & \eng{factories pollute the river} \\ \hline
\eng{waste} (U) & déchets & \eng{plastic waste, waste management, dump waste} \\ \hline
\eng{contamination} & contamination & \eng{water contamination} \\ \hline
\eng{carbon footprint} & empreinte carbone & \eng{reduce one's carbon footprint} \\ \hline
\eng{endangered species} & espèces menacées & \eng{protect endangered species} \\ \hline
\eng{natural disasters} & catastrophes naturelles & \eng{increase in natural disasters} \\ \hline
\eng{single-use plastic} & plastique à usage unique & \eng{ban single-use plastic} \\ \hline
\end{tabularx}
\end{center}
\emph{Attention :} \eng{waste} est indénombrable (pas de \emph{s} au pluriel). On dit \eng{much waste} / \eng{a lot of waste}.
\end{definition}

\courssub{Champ lexical 3 : Solutions and protection}

\begin{definition}[Mots-clés : Protection, énergies et développement durable]
\begin{center}
\begin{tabularx}{\linewidth}{|X|X|X|}\hline
\rowcolor{popgreenL} \textbf{English} & \textbf{Français} & \textbf{Collocation / Note} \\ \hline
\eng{to protect} & protéger & \eng{protect the environment \textbf{from} pollution} \\ \hline
\eng{protection} & protection & \eng{environmental protection} \\ \hline
\eng{to conserve} / \eng{conservation} & conserver / conservation & \eng{wildlife conservation} \\ \hline
\eng{renewable energy} & énergies renouvelables & \eng{solar / wind / hydro power} \\ \hline
\eng{reforestation} & reboisement & \eng{reforestation programme} (un seul \emph{f}) \\ \hline
\eng{to recycle} / \eng{recycling} & recycler / recyclage & \eng{recycle plastic bottles} \\ \hline
\eng{sustainable development} & développement durable & \eng{sustainable agriculture} \\ \hline
\eng{to raise awareness} & sensibiliser & \eng{raise awareness \textbf{about} climate change} \\ \hline
\eng{to take action / measures} & prendre des mesures & \eng{take urgent action} (jamais \eng{make measures}) \\ \hline
\end{tabularx}
\end{center}
\textbf{Prononciation :} \eng{renewable} « ri-niou-e-beul », \eng{sustainable} « so-sté-ne-beul », \eng{reforestation} « ri-fo-res-ta-cheune ».
\end{definition}

\courssub{Méthode 1 — Identifier et employer le vocabulaire du thème environnement}

\begin{methode}[Identifier et employer le vocabulaire du thème]
Pour maîtriser le lexique de l'environnement (\eng{climate change}, \eng{global warming}, \eng{pollution}, \eng{waste}, \eng{protection}), suis cette démarche en cinq étapes :
\begin{enumerate}
\item \textbf{Classe par champ lexical.} Range chaque mot dans l'une des trois familles : causes (\eng{greenhouse gases, carbon emissions, fossil fuels, deforestation}), effets (\eng{global warming, rising sea levels, drought, floods, melting ice caps}) et solutions (\eng{recycling, renewable energy, reforestation, conservation}).
\item \textbf{Apprends le mot avec sa collocation.} Ne mémorise jamais un mot seul : on dit \eng{to reduce emissions}, \eng{to tackle climate change}, \eng{to protect the environment \textbf{from}}, \eng{to clear forests} (ou \eng{cut down trees}), \eng{to dump waste}.
\item \textbf{Note la nature grammaticale.} \eng{pollute} (verbe), \eng{pollution} (nom), \eng{polluted} (adjectif). Cela évite les erreurs de forme dans la phrase.
\item \textbf{Vérifie le sens dans le contexte.} Le mot anglais ne recouvre pas toujours le mot français ressemblant : \eng{the environment} = « l'environnement », mais \eng{an issue} = « un problème / une question », pas « une issue » (\eng{a way out}).
\item \textbf{Réemploie immédiatement.} Construis une phrase personnelle avec chaque nouveau mot, en lien avec le Cameroun si possible : \eng{Floods in Douala are becoming more frequent because of climate change.}
\end{enumerate}
\textbf{Réflexe examen :} dans une question de vocabulaire, cite le mot exact du texte, puis explique-le en anglais simple (\eng{“drought” means a long period without rain}).
\end{methode}

\begin{exoresolu}[Vocabulary in context — reading comprehension]
\textbf{Énoncé.} Read the text and answer the questions in English.

\emph{“In the Far North region of Cameroon, temperatures are rising and rainfall is becoming unpredictable. Farmers say the droughts last longer every year. To tackle the problem, the government has launched a reforestation programme and encourages the use of solar energy, a clean and renewable source of power.”}

1. Find two words or expressions in the text related to the \emph{effects} of climate change.
2. What does “to tackle the problem” mean? Explain in your own words.
3. Give the noun corresponding to the verb “to reforest” and use it in a sentence.
\tcblower
\textbf{Solution détaillée.}

\textbf{Question 1.} On cherche les mots du champ lexical des \emph{effets} du changement climatique. Dans le texte : \eng{temperatures are rising} (ou \eng{rising temperatures}) et \eng{droughts} (on peut aussi citer \eng{unpredictable rainfall}). Réponse rédigée :

\eng{Two words related to the effects of climate change are “temperatures are rising” and “droughts”.}

\textbf{Question 2.} \eng{To tackle} est une collocation clé : \eng{to tackle a problem / an issue} = « s'attaquer à un problème ». Il faut reformuler en anglais simple :

\eng{“To tackle the problem” means to deal with it, to take action in order to solve it.}

\textbf{Question 3.} Famille de mots : verbe \eng{to reforest} $\rightarrow$ nom \eng{reforestation} (attention à l'orthographe : un seul \emph{f}, préfixe \emph{re-} + \emph{forest}). Le nom figure d'ailleurs dans le texte. Phrase personnelle :

\eng{The noun is “reforestation”. Example: Reforestation is essential to fight global warming in the Far North of Cameroon.}

\textbf{Conclusion.} Une bonne réponse de vocabulaire = le mot exact + sa nature + une phrase en anglais correct. Évite de traduire mot à mot : montre que tu comprends le sens en contexte.
\end{exoresolu}

\courssub{Méthode 2 — Réemployer le lexique à l'oral et à l'écrit}

\begin{methode}[Réemployer le lexique dans une production]
Pour réutiliser le vocabulaire de l'environnement dans un dialogue, un exposé ou une composition :
\begin{enumerate}
\item \textbf{Prépare une banque de 8 à 10 mots} avant de parler ou d'écrire : \eng{global warming, greenhouse effect, carbon footprint, waste, to recycle, renewable energy, to pollute, endangered species, sustainable development}.
\item \textbf{Structure ton propos} : problème (\eng{The main problem is...}) $\rightarrow$ causes (\eng{This is caused by... / This is due to...}) $\rightarrow$ conséquences (\eng{As a result, ...}) $\rightarrow$ solutions (\eng{To solve this, we should...}).
\item \textbf{Utilise des connecteurs} : \eng{first of all, moreover, however, as a result, therefore, in conclusion}.
\item \textbf{Varie les formulations} : ne répète pas \eng{pollution} dix fois ; alterne avec \eng{air pollution, water contamination, environmental damage}.
\item \textbf{À l'oral,} soigne la prononciation des mots clés : \eng{environment} « èn-vai-ron-ment », \eng{climate} « klaï-mit », \eng{earth} « eurth », \eng{sustainable} « so-sté-ne-beul ».
\end{enumerate}
\textbf{Structures à retenir :} \eng{We should + BV} (il faudrait), \eng{It is important to + BV}, \eng{If we do not act now, ... will + BV}.
\end{methode}

\begin{exoresolu}[Guided writing — a short paragraph]
\textbf{Énoncé.} Write a paragraph of about 80 words on the following topic: “What can young Cameroonians do to protect the environment?” Use at least five words from the lesson and two link words.
\tcblower
\textbf{Solution détaillée.}

\textbf{Étape 1 — Banque de mots.} \eng{to protect the environment, to recycle waste, to plant trees, renewable energy, pollution, to raise awareness, sustainable}.

\textbf{Étape 2 — Plan.} Introduction du problème $\rightarrow$ deux ou trois actions concrètes $\rightarrow$ conclusion avec \eng{therefore}.

\textbf{Étape 3 — Rédaction.} On utilise \eng{should} + base verbale pour les conseils, et les connecteurs \eng{first of all}, \eng{moreover}, \eng{therefore}.

\eng{Young Cameroonians can do a lot to protect the environment. First of all, they should recycle waste such as plastic bottles instead of throwing them into rivers, because water pollution kills fish and spreads diseases. Moreover, they can plant trees in their schools and neighbourhoods to fight global warming. They should also raise awareness among their friends and families about climate change. Therefore, if every young person acts today, Cameroon will become a cleaner and more sustainable country tomorrow.}

(Environ 80 mots ; cinq mots du lexique et trois connecteurs sont employés.)

\textbf{Conclusion.} Une production réussie = lexique du thème + structures de conseil (\eng{should}, \eng{it is important to}) + connecteurs + phrases simples et correctes. Mieux vaut une phrase simple juste qu'une phrase complexe fautive.
\end{exoresolu}

\courssub{Méthode 3 — Éviter les faux amis et les calques du français}

\begin{methode}[Repérer et corriger les faux amis]
\begin{enumerate}
\item \textbf{Méfie-toi des mots qui ressemblent au français.} Construis-toi une liste personnelle des faux amis du thème :
\begin{center}
\begin{tabularx}{\linewidth}{|X|X|X|}\hline
\rowcolor{popblueL} \textbf{Mot anglais} & \textbf{Sens réel} & \textbf{Piège français} \\ \hline
\eng{an issue} & un problème, une question & pas « une issue » (= \eng{a way out}) \\ \hline
\eng{to damage} & endommager & pas « damager » ; « un dégât » = \eng{damage} (invariable) \\ \hline
\eng{sensible} & raisonnable, sensé & « sensible » (qui ressent) = \eng{sensitive} \\ \hline
\eng{actually} & en fait, réellement & « actuellement » = \eng{currently} / \eng{nowadays} \\ \hline
\eng{to achieve} & accomplir, réaliser & « acheter » = \eng{to buy} \\ \hline
\end{tabularx}
\end{center}
\item \textbf{Ne traduis jamais mot à mot.} « Protéger l'environnement contre la pollution » se dit \eng{to protect the environment \textbf{from} pollution} (préposition \eng{from}, pas \eng{against} dans ce contexte courant).
\item \textbf{Attention aux calques de structure.} « Les gouvernements doivent prendre des mesures » = \eng{governments must \textbf{take} measures / take action} (et non \eng{make measures}).
\item \textbf{Vérifie les noms invariables.} \eng{damage}, \eng{information}, \eng{progress} n'ont pas de pluriel : jamais \eng{damages} au sens de « dégâts » (sauf sens juridique : \eng{damages} = dommages-intérêts).
\item \textbf{Relis en chassant les prépositions françaises.} \eng{to depend \textbf{on}} (pas \eng{of}), \eng{to be responsible \textbf{for}} (pas \eng{of}), \eng{to be interested \textbf{in}}.
\end{enumerate}
\end{methode}

\begin{exoresolu}[Error correction — faux amis et calques]
\textbf{Énoncé.} Each sentence contains one mistake (false friend or word-for-word translation from French). Correct it and justify briefly.

1. \eng{Actually, many species are endangered because of deforestation.}
2. \eng{Factories make a lot of damages to the environment.}
3. \eng{We must protect the environment of pollution.}
4. \eng{The government should make measures to reduce carbon emissions.}
\tcblower
\textbf{Solution détaillée.}

\textbf{1.} Faux ami : \eng{actually} signifie « en fait », pas « actuellement ». Correction : \eng{Currently, many species are endangered because of deforestation.} (« Actuellement, de nombreuses espèces sont menacées à cause de la déforestation. »)

\textbf{2.} Double erreur : \eng{damage} est invariable au sens de « dégâts » et la collocation correcte est \eng{to cause damage} (ou \eng{to do damage}), pas \eng{make damages}. Correction : \eng{Factories cause a lot of damage to the environment.}

\textbf{3.} Calque du français « protéger de » : la préposition anglaise est \eng{from}. Correction : \eng{We must protect the environment from pollution.}

\textbf{4.} Calque de « prendre des mesures » : en anglais on dit \eng{to take measures} ou \eng{to take action}. Correction : \eng{The government should take measures to reduce carbon emissions.}

\textbf{Conclusion.} À chaque relecture, interroge-toi sur trois points : le mot ressemble-t-il au français ? la préposition est-elle la bonne ? la collocation verbe + nom est-elle naturelle en anglais ?
\end{exoresolu}

\courssub{Méthode 4 — Familles de mots, prononciation et orthographe}

\begin{methode}[Exploiter les familles de mots]
\begin{enumerate}
\item \textbf{Construis le tableau de la famille} dès que tu rencontres un mot nouveau : verbe, nom, adjectif.
\begin{center}
\begin{tabular}{|l|l|l|}\hline
\rowcolor{popblueL} \textbf{Verbe} & \textbf{Nom} & \textbf{Adjectif} \\ \hline
\eng{to pollute} & \eng{pollution} & \eng{polluted / polluting} \\ \hline
\eng{to protect} & \eng{protection} & \eng{protective} \\ \hline
\eng{to conserve} & \eng{conservation} & \eng{conservationist} (\eng{conservative} = politique) \\ \hline
\eng{to warm} & \eng{warmth / warming} & \eng{warm} \\ \hline
\eng{to sustain} & \eng{sustainability} & \eng{sustainable} \\ \hline
\eng{to emit} & \eng{emission} & \eng{emissive} \\ \hline
\eng{to destroy} & \eng{destruction} & \eng{destructive} \\ \hline
\end{tabular}
\end{center}
\item \textbf{Choisis la forme selon la fonction dans la phrase.} Après un article ou un adjectif possessif, il faut un nom : \eng{Air pollution is dangerous.} Après un auxiliaire ou un sujet, un verbe : \eng{Factories pollute the air.}
\item \textbf{Soigne l'orthographe des mots pièges} : \eng{environment} (avec un \emph{n} avant \emph{-ment}), \eng{government} (avec un \emph{n}), \eng{which} (pas \eng{wich}), \eng{because} (pas \eng{because of} devant un verbe conjugué).
\item \textbf{Prononciation en lettres françaises} : \eng{pollution} « po-lou-cheune », \eng{renewable} « ri-niou-e-beul », \eng{greenhouse} « gri-ne-haousse », \eng{drought} « draoute ».
\item \textbf{Mémorise les suffixes} : \emph{-tion} (nom d'action : \eng{pollution, protection, emission}), \emph{-able} (possibilité : \eng{renewable, sustainable}), \emph{-al} (adjectif : \eng{environmental, global}), \emph{-ity} (qualité : \eng{sustainability}).
\end{enumerate}
\end{methode}

\begin{exoresolu}[Word formation — complétion]
\textbf{Énoncé.} Complete each sentence with the correct form of the word in brackets.

1. \eng{Air \ldots\ (pollute) is a serious problem in big cities like Douala.}
2. \eng{We must use \ldots\ (renew) energy such as solar power.}
3. \eng{The \ldots\ (protect) of forests is essential for our future.}
4. \eng{Climate change is a \ldots\ (globe) issue.}
5. \eng{Many factories \ldots\ (pollution) the rivers every day.}
\tcblower
\textbf{Solution détaillée.}

\textbf{1.} Après l'adjectif \eng{air} et devant le verbe \eng{is}, il faut un \textbf{nom} : \eng{Air pollution is a serious problem in big cities like Douala.} (suffixe \emph{-tion} : pollute $\rightarrow$ pollution).

\textbf{2.} Devant le nom \eng{energy}, il faut un \textbf{adjectif} : \eng{We must use renewable energy such as solar power.} (suffixe \emph{-able} = « renouvelable »).

\textbf{3.} Après l'article \eng{the}, il faut un \textbf{nom} : \eng{The protection of forests is essential for our future.} (protect $\rightarrow$ protection).

\textbf{4.} Devant le nom \eng{issue}, il faut un \textbf{adjectif} : \eng{Climate change is a global issue.} (globe $\rightarrow$ global, suffixe \emph{-al}).

\textbf{5.} Il manque le \textbf{verbe} de la phrase ; sujet pluriel \eng{factories}, présent simple sans \emph{-s} : \eng{Many factories pollute the rivers every day.} (attention : \eng{pollution} est le nom, pas le verbe).

\textbf{Conclusion.} La méthode infaillible : repérer la fonction grammaticale (sujet ? verbe ? devant un nom ?), puis choisir la forme de la famille avec le bon suffixe.
\end{exoresolu}

\courssub{Méthode 5 — Collocations et expressions idiomatiques en contexte}

\begin{methode}[Employer les collocations justes]
\begin{enumerate}
\item \textbf{Apprends les couples verbe + nom par cœur} : \eng{to tackle climate change, to reduce emissions, to cut down trees, to dump waste, to save energy, to raise awareness, to fight global warming, to ban single-use plastic}.
\item \textbf{Apprends les couples adjectif + nom} : \eng{global warming, renewable energy, natural disasters, endangered species, carbon footprint, fossil fuels, acid rain, severe drought}.
\item \textbf{Connais quelques expressions idiomatiques} : \eng{to go green} (adopter un mode de vie écologique), \eng{the tip of the iceberg} (la partie visible d'un problème bien plus grand), \eng{to be in the same boat} (être tous concernés par le même problème).
\item \textbf{Teste la collocation} : si tu hésites entre deux verbes, demande-toi lequel les anglophones emploient naturellement : on \eng{raises} awareness (pas \eng{rises}), on \eng{meets} a target (pas \eng{touches}), on \eng{reaches} an agreement.
\item \textbf{En contexte,} une collocation bien placée vaut mieux qu'un long discours : \eng{Planting trees is only the tip of the iceberg; we must also reduce carbon emissions.}
\end{enumerate}
\end{methode}

\begin{exoresolu}[Collocations — choix et réemploi]
\textbf{Énoncé.}

A. Choose the correct verb: \eng{We must (do / make / take) action to save the planet.}

B. Match and then write a sentence with each pair: \eng{to raise / to cut down / to dump} with \eng{awareness / trees / waste}.

C. Explain in English the expression “to go green”.
\tcblower
\textbf{Solution détaillée.}

\textbf{A.} La collocation correcte est \eng{to take action} : \eng{We must take action to save the planet.} (\eng{do} et \eng{make} ne se combinent pas avec \eng{action} au sens de « agir ».)

\textbf{B.} Les collocations sont : \eng{to raise awareness} (sensibiliser), \eng{to cut down trees} (abattre des arbres), \eng{to dump waste} (déverser des déchets). Phrases :

\eng{Schools should raise awareness about climate change among students.}

\eng{People who cut down trees without planting new ones destroy the forest.}

\eng{Some companies dump waste into rivers, which pollutes the water.}

\textbf{C.} \eng{“To go green” means to adopt an environmentally friendly lifestyle, for example by recycling, saving energy and using renewable energy.} (« Adopter un mode de vie écologique. »)

\textbf{Conclusion.} Les collocations s'acquièrent par blocs : note toujours le verbe avec son nom habituel, et réemploie le bloc entier dans une phrase personnelle liée à ton environnement quotidien.
\end{exoresolu}

\begin{savaistu}[Le Cameroun et l'action climatique]
Le Cameroun est partie prenante de l'Accord de Paris (2015) et a soumis sa Contribution Déterminée au Niveau National (CDN / \eng{NDC}). Le pays vise une réduction de 35\% de ses émissions de gaz à effet de serre d'ici 2030. Les initiatives phares incluent : le projet \eng{Green Sahel} (reboisement du Grand Nord), le développement de l'hydroélectricité (barrages de Lom Pangar, Nachtigal, Memve'ele) et la promotion du solaire hors-réseau (\eng{off-grid solar}) dans les zones rurales. Le mot \eng{sustainable} (\eng{sustainable development}) est au cœur de la Stratégie Nationale de Développement (SND30).
\end{savaistu}

\begin{attention}[Les 5 erreurs qui coûtent des points au Bac]
\begin{enumerate}
\item \textbf{Orthographe de \eng{environment} et \eng{government}.} On écrit \eng{enviro\textbf{n}ment} et \eng{gover\textbf{n}ment} avec un \emph{n}. Faute très fréquente et très sanctionnée dans une copie sur l'écologie.
\item \textbf{Faux ami \eng{actually}.} \eng{Actually} = « en fait » ; pour « actuellement », écris \eng{currently} ou \eng{nowadays}. \eng{Actually, the climate is changing} ne veut pas dire « actuellement ».
\item \textbf{Calque « prendre des mesures ».} On dit \eng{to \textbf{take} measures / to take action}, jamais \eng{to make measures}. De même : \eng{to \textbf{raise} awareness} (pas \eng{to rise}).
\item \textbf{Prépositions.} \eng{to protect the environment \textbf{from} pollution}, \eng{to depend \textbf{on} fossil fuels}, \eng{to be responsible \textbf{for} pollution}, \eng{to be interested \textbf{in} ecology}. Ne transpose jamais la préposition française.
\item \textbf{Confusion nom / verbe et pluriels interdits.} \eng{pollution} est un nom, le verbe est \eng{to pollute} ; \eng{damage} et \eng{progress} sont invariables (jamais \eng{damages}, \eng{progresses} au sens courant). Vérifie aussi l'accord au présent : \eng{Climate change \textbf{affects} us all.}
\end{enumerate}
\end{attention}

\courssub{Exercices d'entraînement — Bac Blanc}

\begin{exercice}[Vocabulary \& Grammar Consolidation]
\textbf{A. Complete the sentences with the correct word from the box. Two words are extra.}
\eng{carbon footprint / drought / renewable / to dump / endangered / sustainable / emissions / to raise}

1. The long \ldots\ has destroyed the harvest in the Far North.
2. Factories must reduce their CO2 \ldots\ immediately.
3. Solar power is a \ldots\ source of energy.
4. We must not \ldots\ chemical waste into the river.
5. The gorilla is an \ldots\ species in Central Africa.
6. Young people should \ldots\ awareness about recycling.

\textbf{B. Choose the correct option (a, b or c).}
1. The government has decided to \ldots\ measures against illegal logging.
   a) make \quad b) take \quad c) do
2. We must protect the forest \ldots\ destruction.
   a) of \quad b) from \quad c) against
3. \ldots\, the situation is improving thanks to new laws.
   a) Actually \quad b) Currently \quad c) Really
4. \eng{Deforestation} causes a lot of \ldots\ to the ecosystem.
   a) damages \quad b) damage \quad c) damagings

\textbf{C. Word Formation. Complete with the right form of the word in brackets.}
1. \eng{The \ldots\ (destroy) of the mangroves increases the risk of floods.}
2. \eng{We need more \ldots\ (ecology) farming methods.}
3. \eng{This NGO works for the \ldots\ (protect) of wildlife.}
4. \eng{The \ldots\ (warm) of the planet is accelerating.}
\end{exercice}

\begin{corrige}[Exercice Bac Blanc]
\textbf{A.}
1. \eng{drought} (nom, effet climatique).
2. \eng{emissions} (collocation \eng{reduce emissions}).
3. \eng{renewable} (adjectif devant \eng{source}).
4. \eng{dump} (verbe + \eng{waste}, collocation).
5. \eng{endangered} (adjectif composé \eng{endangered species}).
6. \eng{raise} (collocation \eng{raise awareness}).

\textbf{B.}
1. \textbf{b) take} (\eng{take measures/action}).
2. \textbf{b) from} (\eng{protect ... from}).
3. \textbf{b) Currently} (\eng{actually} = en fait).
4. \textbf{b) damage} (invariable).

\textbf{C.}
1. \eng{destruction} (nom, suffixe \emph{-tion}, après \eng{the}).
2. \eng{ecological} (adjectif, suffixe \emph{-al}, qualifie \eng{farming methods}).
3. \eng{protection} (nom, suffixe \emph{-tion}, après \eng{the}).
4. \eng{warming} (nom, \eng{global warming}, après \eng{the}).
\end{corrige}

\newpage

\coursec{Exercices}

\courssub{Je vérifie mes connaissances}

\begin{exercice}[QCM : la bonne collocation]
Pour chaque phrase, choisis la bonne réponse (a, b, c ou d).
\begin{enumerate}
\item We must \_\_\_\_\_ climate change before it is too late. \quad a) make \quad b) tackle \quad c) do \quad d) build
\item Factories \_\_\_\_\_ harmful gases into the atmosphere. \quad a) recycle \quad b) emit \quad c) plant \quad d) melt
\item We should \_\_\_\_\_ plastic bottles instead of throwing them away. \quad a) pollute \quad b) waste \quad c) recycle \quad d) burn
\item Solar and wind power are \_\_\_\_\_ energies. \quad a) renewable \quad b) fossil \quad c) poisonous \quad d) extinct
\item The \_\_\_\_\_ of the earth is rising every decade. \quad a) weather \quad b) temperature \quad c) season \quad d) forecast
\end{enumerate}
\end{exercice}

\begin{exercice}[Vrai ou faux — justifie ta réponse]
Dis si les affirmations suivantes sont vraies (true) ou fausses (false) et justifie ta réponse en une phrase en anglais.
\begin{enumerate}
\item Global warming only affects cold countries.
\item Solar power is a renewable energy.
\item Deforestation helps to reduce carbon dioxide in the atmosphere.
\item A drought is a long period without rain.
\item Burning waste in the open air is good for the environment.
\end{enumerate}
\end{exercice}

\begin{exercice}[Vocabulaire : trouve le mot]
Complète chaque définition avec le mot anglais qui convient, choisi parmi : \eng{flood — drought — greenhouse effect — landfill — endangered species — carbon footprint}.
\begin{enumerate}
\item A long period with no rain at all: \_\_\_\_\_\_\_\_
\item Too much water covering the land, for example in Douala during the rainy season: \_\_\_\_\_\_\_\_
\item The process by which gases trap heat around the earth: the \_\_\_\_\_\_\_\_
\item A place where rubbish is buried: a \_\_\_\_\_\_\_\_
\item Animals that may disappear forever, like the Cross River gorilla: \_\_\_\_\_\_\_\_
\item The amount of carbon dioxide produced by a person or an activity: \_\_\_\_\_\_\_\_
\end{enumerate}
\end{exercice}

\begin{exercice}[Familles de mots : transforme le mot entre parenthèses]
Complète chaque phrase avec la forme correcte du mot entre parenthèses.
\begin{enumerate}
\item The government must take action against air \_\_\_\_\_\_\_\_ (POLLUTE).
\item \_\_\_\_\_\_\_\_ (WARM) temperatures are melting the ice caps.
\item We need to find \_\_\_\_\_\_\_\_ (SUSTAIN) solutions for our cities.
\item The \_\_\_\_\_\_\_\_ (DESTROY) of the rainforest is a serious problem.
\item Many \_\_\_\_\_\_\_\_ (ENVIRONMENT) organisations work in Cameroon.
\end{enumerate}
\end{exercice}

\courssub{Je m'entraîne}

\begin{exercice}[Traduction : du français vers l'anglais]
Traduis les phrases suivantes en anglais en utilisant le vocabulaire de la leçon.
\begin{enumerate}
\item Le réchauffement climatique menace les récoltes de cacao dans la région du Centre.
\item Nous devons réduire nos émissions de gaz à effet de serre.
\item Les énergies renouvelables comme le soleil et le vent ne polluent pas.
\item Chaque année, les inondations détruisent des maisons à Douala.
\item Il faut protéger les espèces menacées de nos forêts.
\end{enumerate}
\end{exercice}

\begin{exercice}[Corrige les faux amis et les calques]
Chaque phrase contient une erreur typique d'un élève francophone. Souligne l'erreur et réécris la phrase correctement.
\begin{enumerate}
\item We must protect the nature.
\item The government took important mesures against pollution.
\item Climate change is actually the biggest problem in the world.
\item The floods caused many damages in the city.
\item We should economise water during the dry season.
\end{enumerate}
\end{exercice}

\begin{exercice}[Texte à trous : les inondations de Douala]
Complète le texte suivant avec dix des mots de la liste : \eng{drought — flood — emissions — waste — recycle — sustainable — climate — rainfall — drainage — residents — pollution — renewable}.

\begin{textebox}[Floods in Douala]
Every year, when heavy \_\_\_\_\_\_\_\_ falls on Douala, the streets fill with water. Last October, a terrible \_\_\_\_\_\_\_\_ destroyed shops and homes in several neighbourhoods. Experts say that \_\_\_\_\_\_\_\_ change is making the rains more violent. Plastic \_\_\_\_\_\_\_\_ blocks the gutters, so the city needs a better \_\_\_\_\_\_\_\_ system. Many \_\_\_\_\_\_\_\_ have lost their belongings. To solve the problem, the city must \_\_\_\_\_\_\_\_ plastic, reduce \_\_\_\_\_\_\_\_ from old vehicles, and build \_\_\_\_\_\_\_\_ housing. Meanwhile, other regions of Cameroon suffer from the opposite problem: a long \_\_\_\_\_\_\_\_ that kills crops and cattle.
\end{textebox}
\end{exercice}

\begin{exercice}[Appariements : expressions idiomatiques]
Associe chaque expression anglaise (1--6) à son équivalent français (a--f).
\begin{enumerate}
\item to go green
\item to leave a carbon footprint
\item to be in the same boat
\item to turn a blind eye to pollution
\item the tip of the iceberg
\item to weather the storm
\end{enumerate}
\begin{itemize}
\item[a)] fermer les yeux sur la pollution
\item[b)] la partie visible du problème
\item[c)] adopter un mode de vie écologique
\item[d)] surmonter la crise
\item[e)] être tous logés à la même enseigne
\item[f)] laisser une empreinte carbone
\end{itemize}
\end{exercice}

\courssub{Je me prépare au Bac}

\begin{exercice}[Compréhension écrite et langue — type Bac]
Lis le texte suivant, puis réponds aux questions en anglais, avec tes propres mots.

\begin{textebox}[A greener future for Cameroon?]
In the village of Nkometou, thirty kilometres from Yaoundé, a group of young farmers has decided to fight back against climate change. Five years ago, their maize and cassava fields were drying up because the rainy season had become unpredictable. “We planted in March like our parents, but the rain simply did not come,” explains Amina, the leader of the group. “We lost almost everything.”

With the help of a local association, the farmers planted two thousand trees around their fields and started using drought-resistant seeds. They also built small tanks to collect rainwater and learnt to produce compost instead of burning their waste. Today, their harvests have improved, and the village has become a model for the whole region.

However, the challenges remain enormous. Across the country, deforestation continues, plastic waste invades the streets of Douala and Yaoundé, and greenhouse gas emissions from old lorries and generators keep rising. “Planting trees is not enough,” Amina admits. “The government, companies and ordinary citizens must all work together. The earth is the only home we have.”
\end{textebox}

\textbf{Comprehension questions}
\begin{enumerate}
\item Why did the farmers of Nkometou lose their crops five years ago? (2 marks)
\item Mention three solutions the farmers adopted to adapt to climate change. (3 marks)
\item According to the text, what environmental problems still exist in Cameroon? (3 marks)
\item What does Amina mean when she says “The earth is the only home we have”? (2 marks)
\end{enumerate}

\textbf{Language work}
\begin{enumerate}
\item Find in the text a synonym of “unpredictable” and an antonym of “improved”. (2 marks)
\item Put the verbs in brackets in the correct tense: “If the government (invest) \_\_\_\_\_\_ in renewable energy, emissions (decrease) \_\_\_\_\_\_.” (2 marks)
\item Turn into the passive voice: “The farmers planted two thousand trees.” (1 mark)
\end{enumerate}
\end{exercice}

\begin{exercice}[Traduction — type Bac]
Traduis le passage suivant en anglais.

\begin{textebox}[Texte à traduire]
Le changement climatique est un défi majeur pour le Cameroun. Dans le Grand Nord, la sécheresse détruit les récoltes, tandis que les inondations frappent régulièrement les habitants de Douala. Pour lutter contre ce phénomène, il faut réduire les émissions de gaz, recycler les déchets plastiques et développer les énergies renouvelables. Chaque citoyen peut agir : planter des arbres, économiser l'eau et éviter de brûler les ordures. Protéger l'environnement, c'est protéger notre avenir.
\end{textebox}
\end{exercice}

\begin{exercice}[Expression écrite — type Bac]
Write a composition of \textbf{150--180 words} on the following topic:

\begin{center}
\emph{“What can young Cameroonians do to fight climate change?”}
\end{center}

In your composition, you should:
\begin{itemize}
\item explain why climate change is a threat to Cameroon (droughts, floods, poor harvests);
\item present at least three concrete actions that young people can take (planting trees, recycling waste, saving water, raising awareness at school);
\item use \textbf{at least eight words or expressions} from this lesson (e.g. \eng{global warming, emissions, renewable energy, sustainable, endangered, carbon footprint, to tackle, to go green});
\item organise your text into a short introduction, a development and a conclusion.
\end{itemize}
\end{exercice}

\newpage

\coursec{Corrigés des exercices}

\begin{corrige}[Exercice 1 — QCM : la bonne collocation]
\begin{enumerate}
\item \textbf{b) tackle} — \eng{We must tackle climate change before it is too late.} « Nous devons nous attaquer au changement climatique avant qu'il ne soit trop tard. » La collocation correcte est \eng{to tackle a problem / climate change}. On ne dit jamais \eng{make climate change} ni \eng{do climate change}.
\item \textbf{b) emit} — \eng{Factories emit harmful gases into the atmosphere.} « Les usines émettent des gaz nocifs dans l'atmosphère. » Collocation : \eng{to emit gases} ; le nom correspondant est \eng{emissions}.
\item \textbf{c) recycle} — \eng{We should recycle plastic bottles instead of throwing them away.} « Nous devrions recycler les bouteilles en plastique au lieu de les jeter. » Noter le phrasal verb \eng{to throw away} (« jeter »).
\item \textbf{a) renewable} — \eng{Solar and wind power are renewable energies.} « L'énergie solaire et l'énergie éolienne sont des énergies renouvelables. » \eng{Fossil} désigne les énergies fossiles (pétrole, charbon) ; \eng{exhaustible} signifie « épuisable », c'est l'inverse du sens attendu.
\item \textbf{b) temperature} — \eng{The temperature of the earth is rising every decade.} « La température de la terre augmente chaque décennie. » Attention : \eng{weather} = le temps qu'il fait ; \eng{climate} = le climat d'une région ; seul \eng{temperature} convient avec \eng{rising}.
\end{enumerate}
\end{corrige}

\begin{corrige}[Exercice 2 — Vrai ou faux]
\begin{enumerate}
\item \textbf{False.} \eng{Global warming affects the whole planet: some regions suffer from droughts, others from floods or rising sea levels.} (Le réchauffement touche toute la planète, pas seulement les pays froids.)
\item \textbf{True.} \eng{Solar power is renewable because the sun's energy will never run out.} (Le soleil est une source inépuisable ; \eng{to run out} = « s'épuiser ».)
\item \textbf{False.} \eng{Deforestation increases carbon dioxide because trees absorb CO2 when they grow.} (Les arbres absorbent le CO2 ; les couper aggrave le problème.)
\item \textbf{True.} \eng{A drought is a long period of time without rain, which kills crops and cattle.} (Définition exacte de \eng{drought} ; \eng{cattle} est un pluriel indénombrable : « le bétail ».)
\item \textbf{False.} \eng{Burning waste in the open air is harmful because it releases toxic gases into the atmosphere.} (Brûler les ordures pollue l'air ; \eng{to release} = « rejeter, libérer ».)
\end{enumerate}
\textbf{Point de méthode :} la justification doit être une phrase complète avec sujet + verbe ; une seule phrase suffit. Commencez par \eng{True} ou \eng{False}, puis justifiez avec vos propres mots.
\end{corrige}

\begin{corrige}[Exercice 3 — Vocabulaire : trouve le mot]
\begin{enumerate}
\item \textbf{drought} (« sécheresse » — prononciation « draoute »)
\item \textbf{flood} (« inondation » — prononciation « fleud »)
\item \textbf{greenhouse effect} (« effet de serre »)
\item \textbf{landfill} (« décharge »)
\item \textbf{endangered species} (« espèces menacées » — \eng{species} a la même forme au singulier et au pluriel)
\item \textbf{carbon footprint} (« empreinte carbone »)
\end{enumerate}
\textbf{Attention à l'orthographe :} \eng{drought} s'écrit avec \emph{-ough} mais se prononce « draoute » ; \eng{flood} se prononce « fleud » (et non « floude ») ; \eng{greenhouse} s'écrit en un seul mot.
\end{corrige}

\begin{corrige}[Exercice 4 — Familles de mots]
\begin{enumerate}
\item \textbf{pollution} — \eng{The government must take action against air pollution.} (nom : verbe \emph{pollute} $\rightarrow$ nom \emph{pollution})
\item \textbf{Warming} — \eng{Warming temperatures are melting the ice caps.} (adjectif participial en \emph{-ing} : « des températures qui se réchauffent »)
\item \textbf{sustainable} — \eng{We need to find sustainable solutions for our cities.} (adjectif : \emph{sustain} $\rightarrow$ \emph{sustainable} ; collocation \eng{sustainable solutions})
\item \textbf{destruction} — \eng{The destruction of the rainforest is a serious problem.} (nom : \emph{destroy} $\rightarrow$ \emph{destruction}, forme irrégulière à mémoriser ; \eng{rainforest} s'écrit en un seul mot)
\item \textbf{environmental} — \eng{Many environmental organisations work in Cameroon.} (adjectif : \emph{environment} $\rightarrow$ \emph{environmental} ; ne pas oublier le \emph{-al} ni le \emph{n} final de \emph{environment})
\end{enumerate}
\end{corrige}

\begin{corrige}[Exercice 5 — Traduction : du français vers l'anglais]
\begin{enumerate}
\item \eng{Global warming is threatening cocoa harvests in the Centre region.} — Présent continu pour une menace en cours ; \eng{harvests} au pluriel ; pas d'article devant \eng{global warming}.
\item \eng{We must reduce our greenhouse gas emissions.} — Ordre des mots : \eng{greenhouse gas emissions} (les compléments précèdent le nom principal \emph{emissions}).
\item \eng{Renewable energies such as the sun and the wind do not pollute.} — Autre formulation correcte : \eng{Renewable energies, like solar and wind power, do not pollute.}
\item \eng{Every year, floods destroy houses in Douala.} — \eng{Floods} au pluriel sans article (vérité générale) ; \eng{every year} en tête de phrase, suivi d'une virgule.
\item \eng{We must protect the endangered species of our forests.} — Autre formulation correcte : \eng{Endangered species in our forests must be protected.} (passif avec \emph{must be} + participe passé).
\end{enumerate}
\textbf{Points de vigilance :} ne pas traduire « il faut » par \eng{it must} ; utiliser \eng{we must} ou le passif. Pas de \emph{s} à \eng{must}, même à la troisième personne.
\end{corrige}

\begin{corrige}[Exercice 6 — Corrige les faux amis et les calques]
\begin{enumerate}
\item Erreur : \eng{the nature}. Correction : \eng{We must protect nature.} — En anglais, \eng{nature} ne prend jamais d'article quand il désigne l'environnement en général.
\item Erreur : \eng{mesures}. Correction : \eng{The government took important measures against pollution.} — Orthographe anglaise : \emph{measures} (avec un \emph{a}). Collocation : \eng{to take measures}.
\item Erreur : \eng{actually}. Correction : \eng{Climate change is currently the biggest problem in the world.} — Faux ami : \eng{actually} = « en fait » ; « actuellement » = \eng{currently} ou \eng{at present}.
\item Erreur : \eng{many damages}. Correction : \eng{The floods caused a lot of damage in the city.} — \eng{Damage} (dégâts) est indénombrable : jamais de \emph{s} ni de \emph{many}. Attention : \eng{damages} avec un \emph{s} existe, mais signifie « dommages et intérêts » au tribunal.
\item Erreur : \eng{economise}. Correction : \eng{We should save water during the dry season.} — Faux ami : « économiser » = \eng{to save} ; \eng{to economise} existe mais est rare et soutenu.
\end{enumerate}
\end{corrige}

\begin{corrige}[Exercice 7 — Texte à trous : les inondations de Douala]
Texte complété :

\eng{Every year, when heavy \textbf{rainfall} falls on Douala, the streets fill with water. Last October, a terrible \textbf{flood} destroyed shops and homes in several neighbourhoods. Experts say that \textbf{climate} change is making the rains more violent. Plastic \textbf{waste} blocks the gutters, so the city needs a better \textbf{drainage} system. Many \textbf{residents} have lost their belongings. To solve the problem, the city must \textbf{recycle} plastic, reduce \textbf{emissions} from old vehicles, and build \textbf{sustainable} housing. Meanwhile, other regions of Cameroon suffer from the opposite problem: a long \textbf{drought} that kills crops and cattle.}

\textbf{Justifications :} \eng{heavy rainfall} (collocation : une pluie « forte » est \emph{heavy}, jamais \emph{strong}) ; \eng{climate change} (collocation figée) ; \eng{plastic waste} (\emph{waste} indénombrable) ; \eng{reduce emissions} (collocation) ; \eng{sustainable housing} (adjectif + nom). Les mots non utilisés sont \eng{pollution} et \eng{renewable}.
\end{corrige}

\begin{corrige}[Exercice 8 — Appariements : expressions idiomatiques]
\begin{center}
\begin{tabular}{|l|l|}\hline
\rowcolor{popblueL} Expression & Équivalent \\ \hline
1. to go green & c) adopter un mode de vie écologique \\ \hline
2. to leave a carbon footprint & f) laisser une empreinte carbone \\ \hline
3. to be in the same boat & e) être tous logés à la même enseigne \\ \hline
4. to turn a blind eye to pollution & a) fermer les yeux sur la pollution \\ \hline
5. the tip of the iceberg & b) la partie visible du problème \\ \hline
6. to weather the storm & d) surmonter la crise \\ \hline
\end{tabular}
\end{center}
\textbf{Remarque :} dans \eng{to weather the storm}, \eng{weather} est un verbe (« résister à ») et non le nom « temps ». Ces expressions ne se traduisent jamais mot à mot : apprenez-les comme des blocs.
\end{corrige}

\begin{corrige}[Exercice 9 — Compréhension écrite et langue — type Bac]
\textbf{Barème indicatif : 15 points (10 compréhension + 5 langue).}

\textbf{Comprehension questions — réponses attendues}
\begin{enumerate}
\item (2 pts) \eng{They lost their crops because the rainy season had become unpredictable: the rain did not come and their fields dried up.}
\item (3 pts, 1 pt par solution) \eng{They planted two thousand trees, they used drought-resistant seeds, and they built tanks to collect rainwater.} (On accepte aussi : \eng{they learnt to produce compost instead of burning waste}.)
\item (3 pts) \eng{Deforestation continues, plastic waste invades the streets of big cities, and greenhouse gas emissions from old lorries and generators keep rising.}
\item (2 pts) \eng{She means that we have no other planet to live on, so everyone — the government, companies and citizens — must protect the earth together.}
\end{enumerate}

\textbf{Language work}
\begin{enumerate}
\item (2 pts) Synonyme de \eng{unpredictable} : \eng{uncertain} ou \eng{changeable} (« imprévisible »). Antonyme de \eng{improved} : \eng{worsened} (« empiré, dégradé »). Toute réponse cohérente et correctement orthographiée est acceptée.
\item (2 pts) \eng{If the government \textbf{invested} in renewable energy, emissions \textbf{would decrease}.} — Conditionnel de type 2 : \emph{if + prétérit, would + base verbale} (hypothèse irréelle au présent).
\item (1 pt) \eng{Two thousand trees were planted by the farmers.} — Passif au prétérit : \emph{were} + participe passé \emph{planted}.
\end{enumerate}

\textbf{Points de vigilance :} répondre par des phrases complètes, avec ses propres mots (ne pas recopier le texte) ; au passif, ne pas oublier l'accord \emph{were} avec le pluriel \emph{trees} ; au conditionnel de type 2, ne jamais mettre \emph{would} dans la proposition en \emph{if}.
\end{corrige}

\begin{corrige}[Exercice 10 — Traduction — type Bac]
\textbf{Barème indicatif : 10 points (2 points par phrase).}

\textbf{Proposition de traduction :}

\eng{Climate change is a major challenge for Cameroon. In the Far North, drought destroys crops, while floods regularly strike the residents of Douala. To fight this phenomenon, we must reduce gas emissions, recycle plastic waste and develop renewable energies. Every citizen can act: plant trees, save water and avoid burning rubbish. Protecting the environment means protecting our future.}

\textbf{Points de vigilance :}
\begin{itemize}
\item « un défi majeur » = \eng{a major challenge} (et non \emph{a big defi}).
\item « tandis que » = \eng{while} ou \eng{whereas}.
\item « il faut » = \eng{we must} ou \eng{it is necessary to} — jamais \eng{it must}.
\item « économiser l'eau » = \eng{save water} (faux ami \emph{economise}).
\item « ordures » = \eng{rubbish} (GB) ou \eng{garbage} (US), indénombrables.
\item Dernière phrase : gérondif sujet \eng{Protecting... means protecting...} — structure élégante à valoriser.
\end{itemize}
\end{corrige}

\begin{corrige}[Exercice 11 — Expression écrite — type Bac]
\textbf{Barème indicatif : 20 points — respect du sujet et de la tâche (4), organisation en introduction/développement/conclusion (4), richesse du vocabulaire de la leçon, au moins 8 mots (6), correction grammaticale (4), longueur 150--180 mots (2).}

\textbf{Proposition de rédaction modèle :}

\eng{Climate change is now a serious threat to Cameroon. In the Far North, long droughts destroy crops, while floods regularly hit Douala. Global warming is making harvests poorer and poorer, so young people cannot turn a blind eye to this problem.}

\eng{Fortunately, there are many concrete actions we can take. First, we should plant trees at school and in our neighbourhoods, because trees absorb greenhouse gases. Secondly, we must recycle plastic waste instead of burning it, in order to reduce pollution and our carbon footprint. Thirdly, we can save water and electricity at home. Finally, young Cameroonians can raise awareness through clubs, debates and social media, and encourage their families to go green.}

\eng{To conclude, if we all adopt a sustainable lifestyle, we will help to tackle climate change. The earth is the only home we have: protecting it means protecting our future.}

(Environ 150 mots. Mots de la leçon employés : \eng{climate change, droughts, floods, global warming, greenhouse gases, recycle, waste, pollution, carbon footprint, save water, go green, sustainable, tackle}.)

\textbf{Points de vigilance :}
\begin{itemize}
\item Compter les mots : en dessous de 140 ou au-dessus de 200, on perd des points.
\item Pas d'article devant \eng{nature} ; \eng{damage} et \eng{waste} indénombrables.
\item « sensibiliser » = \eng{to raise awareness} (et non \emph{sensibilize}).
\item Varier les articulateurs : \eng{first, secondly, finally, to conclude}.
\end{itemize}
\end{corrige}

\newpage

\coursec{Fiche bilan}

\begin{popfigure}
\begin{tikzpicture}[mindmap, grow cyclic, every node/.style={concept, font=\small},
root concept/.append style={concept color=popblue, font=\bfseries},
level 1/.append style={level distance=3.2cm, sibling angle=51},
level 2/.append style={level distance=1.9cm, font=\scriptsize}]
\node[root concept]{Environment}
  child[concept color=popgreen]{ node{Climate change}
    child{ node{global warming} }
    child{ node{greenhouse gases} }
  }
  child[concept color=poporange]{ node{Pollution}
    child{ node{air / water / soil} }
    child{ node{waste, litter} }
  }
  child[concept color=poppurple]{ node{Causes}
    child{ node{deforestation} }
    child{ node{fossil fuels} }
  }
  child[concept color=poppink]{ node{Effects}
    child{ node{droughts, floods} }
    child{ node{rising sea levels} }
  }
  child[concept color=popgold]{ node{Solutions}
    child{ node{recycle, reuse} }
    child{ node{renewable energy} }
  }
  child[concept color=popblue!60]{ node{Collocations}
    child{ node{carbon footprint} }
    child{ node{raise awareness} }
  }
  child[concept color=popmuted]{ node{Faux amis}
    child{ node{sensible $\neq$ sensible} }
  };
\end{tikzpicture}
\legende{Carte mentale : le vocabulaire de l'environnement}
\end{popfigure}

\begin{bilan}
\textbf{1. Lexique clé (English -- Français)}
\begin{itemize}
  \item \eng{climate change} : le changement climatique ; \eng{global warming} : le réchauffement climatique
  \item \eng{greenhouse gases} : les gaz à effet de serre ; \eng{carbon emissions} : les émissions de carbone
  \item \eng{fossil fuels} (coal, oil, gas) : les énergies fossiles ; \eng{renewable energy} (solar, wind) : les énergies renouvelables
  \item \eng{pollution} : la pollution ; \eng{waste / rubbish} : les déchets ; \eng{litter} : les ordures jetées au sol
  \item \eng{deforestation} : la déforestation ; \eng{drought} : la sécheresse ; \eng{flood} : l'inondation
  \item \eng{to recycle / to reuse / to reduce} : recycler / réutiliser / réduire ; \eng{to protect the environment} : protéger l'environnement
  \item \eng{endangered species} : les espèces menacées ; \eng{sustainable development} : le développement durable
\end{itemize}

\textbf{2. Collocations à connaître par c\oe ur}
\begin{center}
\begin{tabularx}{\linewidth}{|X|X|}
\hline
\rowcolor{popblueL} Collocation anglaise & Traduction \\
\hline
\eng{to cause / to combat climate change} & provoquer / combattre le changement climatique \\
\hline
\eng{to reduce one's carbon footprint} & réduire son empreinte carbone \\
\hline
\eng{to raise awareness (of / about)} & sensibiliser (à) \\
\hline
\eng{to run out of natural resources} & épuiser les ressources naturelles \\
\hline
\eng{to cut down on waste} & réduire les déchets \\
\hline
\eng{to ban plastic bags} & interdire les sacs plastiques \\
\hline
\eng{global temperatures are rising} & les températures mondiales augmentent \\
\hline
\end{tabularx}
\end{center}

\textbf{3. Points de langue à retenir}
\begin{itemize}
  \item \eng{environment} est \textbf{invariable} au sens de « milieu naturel » : on dit \eng{the environment} (avec \emph{the}), jamais \eng{environments} dans ce sens.
  \item \eng{pollution}, \eng{waste}, \eng{evidence}, \eng{information} sont \textbf{indénombrables} : pas de \emph{a}, pas de pluriel (\eng{some pollution}, \eng{a piece of information}).
  \item Prépositions : \eng{to protect \textbf{from}}, \eng{to suffer \textbf{from}}, \eng{to depend \textbf{on}}, \eng{to be responsible \textbf{for}}, \eng{to be aware \textbf{of}}.
  \item Familles de mots : \eng{to pollute} (verbe) $\rightarrow$ \eng{pollution} (nom) $\rightarrow$ \eng{polluted / polluting} (adjectifs) ; \eng{to sustain} $\rightarrow$ \eng{sustainable} $\rightarrow$ \eng{sustainability}.
  \item Prononciation : \eng{climate} « klaï-mèt », \eng{drought} « draoute », \eng{threat} « srète », \eng{environment} « in-vaï-ron-mente » (le \emph{n} avant \emph{m} se prononce).
\end{itemize}

\textbf{4. Faux amis et pièges}
\begin{itemize}
  \item \eng{sensible} = raisonnable ; « sensible » = \eng{sensitive}.
  \item \eng{actually} = en fait ; « actuellement » = \eng{currently}.
  \item \eng{a waste} : attention, \eng{to waste} = gaspiller ; « un déchet » = \eng{a piece of waste / waste} (indénombrable).
  \item \eng{the nature} : on dit simplement \eng{nature} (sans \emph{the}) pour « la nature ».
\end{itemize}

\textbf{5. Méthode express}
\begin{itemize}
  \item À l'oral : annoncer le problème (\eng{One major issue is...}), donner une cause (\eng{This is caused by...}), proposer une solution (\eng{We should...}).
  \item À l'écrit : introduire (\eng{Nowadays...}), développer avec \eng{firstly, moreover, however}, conclure avec \eng{In conclusion, it is high time we acted}.
\end{itemize}
\end{bilan}

\begin{aretenir}[Checklist avant le Bac]
\begin{itemize}
  \item[$\square$] Je connais au moins 20 mots du champ lexical de l'environnement avec leur traduction.
  \item[$\square$] Je sais distinguer \eng{climate change} et \eng{global warming}.
  \item[$\square$] Je maîtrise 5 collocations : \eng{carbon footprint}, \eng{raise awareness}, \eng{run out of}, \eng{cut down on}, \eng{ban plastic bags}.
  \item[$\square$] Je sais que \eng{pollution}, \eng{waste} et \eng{information} sont indénombrables.
  \item[$\square$] J'emploie les bonnes prépositions : \eng{protect from}, \eng{depend on}, \eng{responsible for}, \eng{aware of}.
  \item[$\square$] Je ne confonds plus \eng{sensible} / \eng{sensitive} et \eng{actually} / \eng{currently}.
  \item[$\square$] Je sais construire les familles de mots : \eng{pollute -- pollution -- polluted}.
  \item[$\square$] Je prononce correctement \eng{climate}, \eng{drought}, \eng{threat}, \eng{environment}.
  \item[$\square$] Je sais rédiger un paragraphe d'opinion sur l'environnement avec des connecteurs (\eng{firstly, moreover, however, in conclusion}).
  \item[$\square$] Je peux proposer des solutions avec \eng{should}, \eng{must}, \eng{it is high time + prétérit}.
\end{itemize}
\end{aretenir}

\findecours

\end{document}
